% Printed 428; principal PDF 435; Google PDF 437.
\par\clearpage
\noindent\makebox[0pt][l]{\fontsize{13}{16}\selectfont 428}\hfill{\small\textbf{NACHLASS. ZUR HANNOVERSCHEN TRIANGULATION.}}\hfill\null\par\medskip
\begin{center}
\begingroup\setlength{\tabcolsep}{4pt}\renewcommand{\arraystretch}{1.25}
\begin{tabular}{|l|rrr|rrrr|}
\hline
\multicolumn{1}{|c|}{Dreieckspunkt}&\multicolumn{3}{c|}{Breite}&\multicolumn{4}{c|}{Länge v. Gött. Merid.}\\\hline
Timpenberg&$53^\circ$&$7'$&$22''{,}103$&$-$&$0^\circ$&$19'$&$35''{,}387$\\
Nindorf&53&8&56,063&$-$&0&20&32,170\\
Bullerberg&53&9&47,220&$+$&0&31&45,928\\
Wilsede&53&10&9,647&$+$&0&0&11,329\\
Lüneburg, Mich.-Turm&53&15&4,641&$-$&0&27&29,462\\
Brüttendorf&53&15&57,245&$+$&0&40&42,934\\
Garlste&53&15&58,486&$+$&1&13&37,091\\
Zeven&53&17&55,135&$+$&0&39&43,446\\
Litberg&53&23&20,206&$+$&0&19&45,061\\
Varel&53&23&56,979&$+$&1&48&24,812\\
Brillit&53&24&47,035&$+$&0&57&0,384\\
Hamburg, Mich.-Turm&53&33&0,900&$-$&0&2&8,755\\
Bremerlehe&53&34&7,239&$+$&1&21&2,134\\
Jever&53&34&26,427&$+$&2&2&26,615\\
Langwarden&53&36&20,450&$+$&1&38&7,093\\
\end{tabular}\endgroup
\end{center}
\par\noindent\hspace*{\parindent}Die Angabe für die Göttinger Sternwarte bezieht sich auf den Theodolitstandpunkt. Er liegt im Meridian durch die Mitte der Achse des \textsc{Reichenbach}schen Meridiankreises, aber 5,507 m weiter nördlich.

\par\noindent\hspace*{\parindent}Die Breite der Achse des Meridiankreises ist aufgrund der in meiner »Bestimmung des Breitenunterschiedes u. s. w.« angeführten Beobachtungen $=51^\circ31'47''{,}85$ angenommen worden. Im Frühjahr 1828 habe ich neue, noch zahlreichere Beobachtungen angestellt. Ihre endgültige Berechnung kann erst erfolgen, wenn die Teilungsfehler der bestimmten Teilstriche, auf die sich die Beobachtungen bezogen, mit größter Genauigkeit bestimmt worden sind. Es lässt sich aber schon absehen, dass das Ergebnis höchstens um ein paar Zehntelbogensekunden vom obigen abweichen wird.

\par\noindent\hspace*{\parindent}Bei der Berechnung der obigen Breiten und Längen sind noch \textsc{Walbecks} Dimensionen der Erde zugrunde gelegt worden.
\begin{center}\rule{2.6cm}{0.4pt}\end{center}

% Printed 429; principal PDF 436; Google PDF 438.
\par\clearpage
\noindent\null\hfill{\small\textbf{BEMERKUNGEN. ZUR HANNOVERSCHEN TRIANGULATION.}}\hfill\makebox[0pt][r]{\fontsize{13}{16}\selectfont 429}\par\medskip
\begin{center}\textbf{BEMERKUNGEN.}\end{center}
\begingroup\small
\par\noindent\hspace*{\parindent}Der Plan zu dem von \textsc{Gauss} beabsichtigten Werk über die trigonometrischen Messungen in Hannover steht auf einem kleinen Blatt. Die Einleitung und der unvollendete erste Abschnitt dieses Werkes sind 2 Blättern entnommen, die in Buchform zusammengelegt waren. Es ist nicht ersichtlich, wann \textsc{Gauss} mit ihrer Ausarbeitung begonnen hat. Nach der Bemerkung über die kurhessischen Messungen oben auf S. 403 muss dies jedoch vor 1835 gewesen sein (denn in diesem Jahr wurden die Arbeiten für die kurhessische Triangulation, die seit dem Frühjahr 1824 eingestellt waren, wieder aufgenommen. \textsc{Gerling}, Beiträge zur Geographie Kurhessens etc. S. V und VIII). Von den Berichten an das hannoversche Kabinettsministerium sind die beiden ersten aus den Jahren 1820 und 1821 über die notwendigen Instrumente und die Vorarbeiten für die Gradmessung bereits (im Auszug) in Band IV, S. 485/486 und S. 487/489 abgedruckt. Insgesamt sind im Gauss-Archiv gegen 40 meist sehr umfangreiche und ausführliche Berichte über den Fortgang der Triangulationsarbeiten vorhanden, zum Teil mit Übersichtskarten. Auszüge aus einigen dieser Berichte, die sich auf den allgemeinen Teil der Triangulation beziehen, stehen unter [2]. Auch die Originalakten des amtlichen Schriftwechsels zur Gradmessung und zur Landesvermessung befinden sich im Gauss-Archiv. Einschließlich der Arbeitsberichte stammen etwa 90 Schriftstücke von \textsc{Gauss}' Hand.

\par\noindent\hspace*{\parindent}Im Jahr 1816 hatte \textsc{Schumacher} von König \textsc{Friedrich} VI. von Dänemark den Auftrag erhalten, eine Gradmessung auf dem Meridian von Skagen bis Lauenburg und auf dem Breitenkreis von Copenhagen bis zur Westküste Jütlands durchzuführen. Der zu messende Meridianbogen erstreckte sich über $4\frac12^\circ$. Die Fortsetzung dieses Bogens durch Hannover, wodurch er eine Länge von $6\frac12^\circ$ erreichen würde, regte \textsc{Schumacher} bei \textsc{Gauss} sofort an. \textsc{Gauss} interessierte sich zwar sehr für die »herrliche große Unternehmung«, glaubte aber bei der hannoverschen Regierung noch keine entsprechenden Wünsche äußern zu dürfen (vgl. S. 345), wohl weil der Bau seiner Sternwarte erst kurz zuvor abgeschlossen worden war und ihre Ausstattung noch Ausgaben erforderte. Da wandte sich \textsc{Schumacher} im Juli 1817 persönlich an den Minister \textsc{von Arnswaldt} in Hannover und erreichte, dass \textsc{Gauss} zunächst zu einem »Memoire« über die Fortsetzung der dänischen Breitengradmessung durch Hannover aufgefordert wurde. Als \textsc{Schumacher} 1818 seine südlichsten Dreiecke maß, bat er \textsc{Gauss}, in Lüneburg Anschlussmessungen vorzunehmen. Der dortige Michaelisthurm war von den dänischen Dreieckspunkten Hamburg (Michaelisthurm), Hohenhorn und Lauenburg aus sichtbar. \textsc{Gauss} hatte jedoch Bedenken, die Erlaubnis seiner Regierung einzuholen, weil über die Gradmessung in Hannover noch nichts beschlossen war. In einem Brief an \textsc{Schumacher} vom 12. August 1818 sagt er:

\par\noindent\hspace*{\parindent}»Schon im vergangenen Herbst, gleich nachdem ich Ihre Notizen erhalten hatte, habe ich ein Memoire über Ihre Gradmessung verfasst und die vielfältigen Vorteile, die ihre künftige Fortsetzung durch Hannover haben würde, so weit wie möglich hervorgehoben, sodass ich nun nichts mehr hinzuzufügen wüsste. Ich habe dieses Memoire eingesandt, aber bis heute ist nichts weiter darauf erfolgt. Unter allen schwierigen Künsten ist die Kunst des Drängens diejenige, für die ich --- freilich zu meinem großen Nachteil --- am wenigsten Talent habe und am wenigsten geeignet bin. Unter den gegebenen Umständen kann ich deshalb wohl nicht schriftlich auf die Sache quaestionis zurückkommen.«

\par\noindent\hspace*{\parindent}Erneut wandte sich \textsc{Schumacher} an den Minister \textsc{von Arnswaldt}, und \textsc{Gauss} erhielt den Auftrag, die zur Verbindung der hannoverschen und dänischen Triangulation notwendigen Messungen in Lüneburg vorzunehmen (vgl. S. 347 oben). Diese Messungen führte \textsc{Gauss} gemeinsam mit \textsc{Schumacher} in der ersten Hälfte des Oktobers 1818 aus. Dabei gab ihm ein von der Sonne beleuchtetes Fenster des Michaelisthurms in Hamburg, das ihn bei seinen Beobachtungen störte, auch die erste Anregung zur Erfindung des Heliotrops im Herbst 1820.

% Printed 430; principal PDF 437; Google PDF 439.
\par\clearpage
\noindent\makebox[0pt][l]{\fontsize{13}{16}\selectfont 430}\hfill\textbf{BEMERKUNGEN.}\hfill\null\par\medskip
\par\noindent\hspace*{\parindent}Über die Durchführung einer hannoverschen Triangulation wurde aber noch immer nicht entschieden. In einem Brief an \textsc{Schumacher} vom 25. November 1818 sagt \textsc{Gauss}:

\par\noindent\hspace*{\parindent}»Den Bericht über meine Reise [nach Lüneburg] habe ich schon vor längerer Zeit nach Hannover geschickt und darin auch dargelegt, dass ein größerer Theodolit rechtzeitig bestellt werden muss. Eine Antwort habe ich bisher aber nicht erhalten. Weiter drängen kann ich nicht und \textls[120]{mag} ich nicht, denn ich kann überhaupt nur dann eine Arbeit erwarten, die mir Freude macht, wenn man sich \textls[120]{gern} darauf einlässt. Andernfalls, und wenn allerlei einengende Rücksichten zu nehmen wären, würde ich keine Freude daran haben. Ich werde also den Ausgang ruhig abwarten.«

\par\noindent\hspace*{\parindent}\textsc{Schumacher} nutzte nun seinen Aufenthalt in England im April 1819, wo er den bei der englischen Triangulation benutzten \textsc{Ramsden}schen Zenitsektor abholen wollte, der ihm für seine astronomischen Messungen geliehen worden war, um in London maßgebende Persönlichkeiten wie Sir \textsc{Joseph Banks} für die hannoversche Gradmessung zu gewinnen. Er veranlasste außerdem den dänischen Gesandten, in dieser Sache bei Graf \textsc{Münster}, dem Minister für die hannoverschen Angelegenheiten in London, Schritte zu unternehmen. Auf \textsc{Schumachers} Anregung richtete \textsc{Gauss} daraufhin zwei Berichte an Graf \textsc{Münster} und an \textsc{v. Arnswaldt}. Beide sind in Band IV abgedruckt, der erste auf S. 482/483, der zweite auf S. 484/485. Der Vermittlung \textsc{Schumachers}, der im Juni 1819 erneut bei \textsc{v. Arnswaldt} in Hannover vorstellig wurde, war es außerdem zu verdanken, dass \textsc{Gauss} Ende Juni und Anfang Juli 1819 an den Beobachtungen mit dem Zenitsektor in Lauenburg teilnehmen konnte. Im folgenden Jahr wollte \textsc{Gauss} der Basismessung in Holstein beiwohnen. Am 18. Januar 1820 schrieb er an \textsc{Schumacher}: »Es würde mir große Freude machen, wenn es möglich wäre, in diesem Jahr nochmals einige Wochen in Ihrer Gesellschaft und bei Ihren Arbeiten zu verbringen. Zum einen käme es aber darauf an, wann Sie sich wieder in jener Gegend aufhalten werden, zum anderen fürchte ich, mich lästig zu machen, wenn ich zum dritten Mal in Hannover eine Reise beantrage, die nur in einem gewissen Zusammenhang mit einer \textls[120]{möglichen}, aber vielleicht noch weit entfernten Unternehmung in unserem Königreich zu stehen scheinen muss.«\enspace .\ .\ .\ .\enspace Da \textsc{Gauss} von \textsc{Schumacher} um eine Besprechung der dänischen Gradmessung gebeten worden war, heißt es im selben Brief weiter:

\par\noindent\hspace*{\parindent}»Teilen Sie mir bitte die (wenn auch vorerst nur vorläufigen) \textls[120]{Ergebnisse} der Beobachtungen mit dem Sektor mit, die Sie im Januar und Februar dieses Jahres in Copenhagen anstellen, und zwar hauptsächlich solche, zu denen ich hier noch entsprechende Beobachtungen machen kann, entweder mit dem \textsc{Reichenbach}schen Kreis oder, falls sich das Eintreffen der neuen Klemmarme wider Erwarten bis in den Februar hinein verzögern sollte, vorerst mit dem \textsc{Repsold}schen. Ich werde diese Ergebnisse dann in unseren Gelehrten Anzeigen bekannt machen und die Gelegenheit nutzen, Ihrer Gradmessung insgesamt die gebührende Aufmerksamkeit zu schenken. Dazu muss ich Sie aber aus dem oben genannten Grund bitten, mir eine knappe Darstellung der Hauptpunkte zu schicken, umso mehr, als Sie vielleicht den einen oder anderen Umstand vorerst noch nicht erwähnt wissen möchten. Auf ein paar Wochen früher oder später wird es wohl nicht ankommen. Dass das Ganze auf eine möglichst \textls[120]{ungezwungene} Art zur Sprache kommt, ist auch \textls[120]{mir} wichtig, weil ich auf keinen Fall den Anschein erwecken möchte, als wollte ich unsere Regierung auf diesem Umweg wieder an die Sache erinnern. Denn so bereitwillig ich die Dreiecke bis Göttingen usw. fortführe, wenn die nötigen Mittel auf angemessene Weise bereitgestellt werden, liegt mir dabei allein am wissenschaftlichen Interesse. Persönlich sehe ich es vielmehr als ein Opfer an, das ich unter der oben genannten Voraussetzung aber sehr gern bringe.«

\par\noindent\hspace*{\parindent}Auf \textsc{Schumachers} Betreiben erhielt das dänische Außenministerium von seinem König den Auftrag, die hannoversche Regierung um \textsc{Gauss}' Anwesenheit bei der Braaker Basis-\par

% Printed 431; principal PDF 438; Google PDF 440.
\clearpage
\noindent\null\hfill\textbf{ZUR HANNOVERSCHEN TRIANGULATION.}\hfill\makebox[0pt][r]{\fontsize{13}{16}\selectfont 431}\par\medskip
\noindent messung zu bitten. Über seine Teilnahme (siehe auch Band IV, S. 486/487), die einschließlich der Reise vom 12. September bis zum 25. Oktober 1820 dauerte, liegt dem Kabinettsministerium ein Bericht vom 1. November 1820 vor, in dem es heißt:

\par\noindent\hspace*{\parindent}»Während meines Aufenthalts in Holstein habe ich gemeinsam mit Professor \textsc{Schumacher} die vereinbarten Beobachtungen, Messungen und Versuche angestellt. Ich habe ferner einem Teil der Basismessung beigewohnt und schließlich alle Vereinbarungen und Vorkehrungen getroffen, die unsere Fortsetzung der Gradmessung betreffen.\enspace .\ .\ .\ .\enspace Der von \textsc{Repsold} in Hamburg für diese Basismessung angefertigte Apparat übertrifft an Genauigkeit, Solidität und Zweckmäßigkeit alle anderen, die bei ähnlichen Gelegenheiten verwendet wurden.«

\par\noindent\hspace*{\parindent}Am 9. Mai 1820 erging schließlich die Kabinettsorder \textsc{Georgs} IV., des Königs von Großbritannien etc. und Hannover, an das Ministerium, die \textsc{Gauss} mit »der Fortsetzung der dänischen Gradmessung« durch Hannover betraute (vgl. S. 347). Die Mitteilung des Ministeriums an \textsc{Gauss} erfolgte unter dem Datum des 30. Juni 1820 (vgl. S. 418). Die Messungen für die Gradmessung führte \textsc{Gauss} in den Jahren 1821, 1822 und 1823 aus.

\par\noindent\hspace*{\parindent}Der dänisch-hannoversche Meridianbogen ließ sich nach Süden erheblich ausdehnen. Er konnte bis zur Insel Elba fortgeführt werden und sich über etwa $16^\circ$ erstrecken (vgl. S. 402). Über die Punkte Brocken und Inselsberg war die hannoversche Gradmessung mit den Dreiecken des preußischen Generalstabs verbunden. An die Seite Hohehagen-Inselsberg grenzten die kurhessischen Dreiecke, deren Messung 1822 \textsc{Gerling} übertragen worden war. Südlich davon waren in Hessen-Darmstadt, Württemberg, Bayern und Österreich die geodätischen Arbeiten teils im Gange, teils geplant. Getrennt von diesem Dreieckssystem bestand der große englisch-französische Bogen auf dem Meridian von Paris, der entlang der Ostküste Spaniens bis zur Insel Formentera reichte (vgl. S. 419). Im Norden schlossen sich an den französischen Bogen die 1801--1811 von Generalleutnant \textsc{von Krayenhoff} in Belgien, den Niederlanden, Ostfriesland und Oldenburg gemessenen Dreiecke an (Précis historique des opérations géodésiques et astronomiques, faites en Hollande etc. La Haye 1815). Das \textsc{Krayenhoff}sche Dreiecksnetz war zwar in seinem südöstlichen Teil zwischen Nederweert und Nimwegen mit dem \textsc{Tranchot}schen verbunden und dieses über die Seite Nürburg-Fleckert mit der \textsc{Müffling}schen Dreieckskette »von Berlin nach dem Rhein«. Über die beiden Letzteren war jedoch nichts veröffentlicht (vgl. S. 365), und auch später wurde nichts darüber veröffentlicht. Zudem sollten die \textsc{Tranchot}schen Dreiecke nur als Grundlage für eine Karte dienen. Die so hergestellte Querverbindung der beiden großen Dreieckssysteme konnte daher nicht als ausreichend gelten. In einem Promemoria an den Bremer Senat, das dieser an die hannoversche Regierung weiterleitete, schlug \textsc{Olbers} gegen Ende des Jahres 1823, wohl im Einvernehmen mit \textsc{Gauss}, vor, den dänisch-hannoverschen Bogen von den Seiten Hamburg-Wilsede und Wilsede-Falkenberg aus über Bremen und Oldenburg mit der \textsc{Krayenhoff}schen Seite Varel-Jever und dadurch mit der englisch-französischen Gradmessung zu verbinden. In einer ausführlicheren Darlegung vom 7. Januar 1824, die ebenso wie das Promemoria von \textsc{Olbers} im Gauss-Archiv liegt, erklärte \textsc{Gauss} seine Zustimmung zur Fortsetzung der Gradmessung nach Jever. Ihre Ausführung wurde dann durch ein Reskript des Grafen \textsc{Münster} vom 15. Februar 1824 angeordnet und ihm am 8. März desselben Jahres vom Kabinettsministerium übertragen (vgl. S. 419). Man hätte den Anschluss auch weiter südlich an die \textsc{Krayenhoff}sche Seite Bentheim-Kirchhesepe herstellen können. Obwohl der nördlichere Weg kürzer war, bot er größere Schwierigkeiten im Gelände als der südlichere. \textsc{Gauss} zog ihn dennoch vor, weil sich auf diesem Weg die Dreiecke mit der Nordsee verbinden ließen und sich dadurch die relativen Höhen seiner Dreieckspunkte in absolute Höhen über der Meeresoberfläche umrechnen ließen (vgl. S. 419). Als weitere Begründung für eine Fortsetzung der Gradmessung nach Varel-Jever gibt \textsc{Gauss} in dem
% Printed 432; principal PDF 439; Google PDF 441.
\par\clearpage
\noindent\makebox[0pt][l]{\fontsize{13}{16}\selectfont 432}\hfill\textbf{BEMERKUNGEN.}\hfill\null\par\medskip
\noindent erwähnten Bericht vom 7. Januar 1824 noch Folgendes an: »Insofern eine solche Verbindung, die quer von Ost nach West verläuft, größtenteils über hannoversches Gebiet führt, liegt der Vorteil für die Geografie des Königreichs ebenso auf der Hand. Es ist jetzt allgemein anerkannt, dass eine genaue Landesvermessung ohne eine gehörige Triangulation unmöglich ist. Bloße Detailmessungen lassen sich niemals mit Sicherheit zu einem unverzerrten Ganzen verbinden. Aber auch abgesehen von der unvergleichlich größeren Genauigkeit gewinnt eine Detailaufnahme, die sich auf eine gute vorangehende Triangulation stützt, in ihrem gesamten Plan und Ablauf so sehr an Leichtigkeit, Einfachheit, Sicherheit und Überprüfbarkeit jedes einzelnen Teils, dass die Hälfte der Zeit und Kosten gespart wird. Die Gradmessungsdreiecke umspannen bereits einen sehr beträchtlichen Teil des Königreichs. Quer verlaufende Verbindungsdreiecke würden den umspannten Raum fast verdoppeln.«

\par\noindent\hspace*{\parindent}Die Beobachtungen für die Fortsetzung der Gradmessung nach Jever fanden 1824 und 1825 statt.

\par\noindent\hspace*{\parindent}Über die Seite Varel-Jever war der Anschluss an die \textsc{Krayenhoff}schen Dreiecke hergestellt. Ihr nördlicher Teil erwies sich jedoch als sehr ungenau, wie eine direkte Prüfung durch Nachmessung von Winkeln in Jever zeigte und auch eine Ausgleichung durch \textsc{Gauss}, Suppl. theor. comb. Art. 23, bestätigte (vgl. die Briefe an \textsc{Bessel} vom 12. März und 20. November 1826, S. 360/362, an \textsc{Olbers} vom 4. Juli 1824, S. 370, und vom 14. Mai 1826, S. 321/322, und an \textsc{Gerling} vom 12. September und 14. November 1838, S. 391 und S. 393). Damit war auch diese Verbindung der beiden großen Meridianbögen nur von geringem Wert. Auf der Rückreise von den Messungen Ende Juli 1825 beriet sich deshalb \textsc{Gauss} (nach seinem Bericht vom 11. März 1827 über die trigonometrischen Arbeiten im Jahr 1825) in Bremen mit \textsc{Olbers}, wie dem abzuhelfen sei. Da die südlichen \textsc{Krayenhoff}schen Dreiecke genauer waren als die nördlichen (vgl. S. 361), schlug \textsc{Olbers} vor, »eine neue Reihe von Dreiecken anzufangen, die von Bremen aus südwestwärts durch das Osnabrückische verlaufen und bei Bentheim eine neue Verbindung mit den \textsc{Krayenhoff}schen Dreiecken herstellen würden. Er fügte die Versicherung hinzu, dass der Bremer Senat aus Interesse an wissenschaftlichen Unternehmungen die Mithilfe des Gehilfen \textsc{Klüver} [den die Stadt Bremen bereits im Vorjahr \textsc{Gauss} zur Unterstützung bei den Messungen zugeteilt hatte] für den Rest des Jahres dazu gern bewilligen würde«. \textsc{Gauss} hatte jedoch Bedenken, auf diesen Vorschlag einzugehen, da er »seinen Auftrag überschritten hätte« und ihm außerdem »das Gelände von Bremen bis Bentheim ganz unbekannt war«. Diese Verbindung von der Seite Bremen-Steinberg bis zur Seite Kirchhesepe-Bentheim wurde dann später, im Jahr 1829, von seinem Sohn, dem Leutnant \textsc{Joseph Gauss}, hergestellt (vgl. den Brief an \textsc{Bessel} vom 9. April 1830, S. 363, sowie auch S. 422), allerdings mit geringerer Genauigkeit (vgl. S. 342). Das noch bestehende Vorhaben (nach dem Arbeitsbericht für 1824), gemeinsam mit \textsc{Schumacher} über Wangeroog, Neuwerk und einen dänischen Dreieckspunkt an der schleswig-holsteinischen Küste die Insel Helgoland, deren Längenunterschied gegenüber Greenwich, Altona und Bremen 1824 chronometrisch bestimmt worden war (vgl. Band VI, S. 455/459), an die Triangulation anzuschließen, wurde nicht mehr ausgeführt, ebenso wenig wie eine geplante neue Verbindung der dänischen und hannoverschen Gradmessung über die Seite Litberg-Hamburg.

\par\noindent\hspace*{\parindent}Im Frühjahr 1827 beobachtete \textsc{Gauss} den Breitenunterschied zwischen Göttingen und Altona mit dem \textsc{Ramsden}schen Zenitsektor, den die englische Regierung auch ihm geliehen hatte. Dabei unterstützte ihn der dänische Ingenieurleutnant \textsc{v. Nehus}. Als Hauptmann \textsc{Müller} den Zenitsektor nach England zurückbrachte, richtete er am 3. Oktober 1827 ein Promemoria an Graf \textsc{Münster}, in dem er die Erweiterung der Triangulation auf das ganze Königreich Hannover vorschlug. Zugleich teilte er mit, dass \textsc{Gauss} bereit sei, die Leitung der Arbeiten zu übernehmen, wenn sein Gesundheitszustand dies zuließe. Hierzu hatte sich \textsc{Gauss} schon früher in einem Schreiben an \textsc{Münster} erboten (nach dem Brief an \textsc{Olbers} vom 2. April 1826, S. 376). Am 8. November wurde \textsc{Gauss} daraufhin vom

% Printed 433; principal PDF 440; Google PDF 442.
\par\clearpage
\noindent\null\hfill\textbf{ZUR HANNOVERSCHEN TRIANGULATION.}\hfill\makebox[0pt][r]{\fontsize{13}{16}\selectfont 433}\par\medskip
\noindent Ministerium aufgefordert, entsprechende Vorschläge zu machen. Aufgrund seines auf S. 413/416 im Auszug wiedergegebenen Berichts wurde durch eine Kabinettsorder König \textsc{Georgs} IV. vom 25. März 1828 an das Ministerium befohlen, »die Triangulation fortzusetzen und zu vervollständigen, um sie zur Verbesserung der Landesgeografie und zur Anfertigung genauer Karten zu verwenden«. Das Ministerium übermittelte \textsc{Gauss} den Auftrag, die Landesvermessung zu leiten, durch Reskript vom 14. April 1828.

\par\noindent\hspace*{\parindent}Die Entstehung des auf S. 418/425 mitgeteilten historischen Berichts ist wohl auf die Trennung Englands und Hannovers im Jahr 1837 zurückzuführen.
\par\bigskip
\par\noindent\hspace*{\parindent}Nach den Beobachtungsbüchern verteilen sich die von \textsc{Gauss} selbst ausgeführten Messungen zeitlich wie folgt auf die einzelnen Stationen.

\par\noindent\hspace*{\parindent}1821 wurde beobachtet: auf der Göttinger Sternwarte vom 24. Juni bis 13. Juli, auf dem nördlichen Meridianzeichen vom 13. bis 17. Juli, auf dem Hohehagen vom 19. bis 29. Juli, auf dem Hils vom 7. bis 27. August und auf dem Brocken vom 2. bis 23. September.

\par\noindent\hspace*{\parindent}1822 wurde gemessen: in Lichtenberg vom 18. Juni bis 4. Juli, auf dem Deister vom 6. bis 16. Juli, in Garssen vom 18. Juli bis 3. August, in Falkenberg vom 5. August bis 6. September, in Hauselberg vom 7. bis 12. September, in Breithorn vom 13. bis 16. September, in Wulfsode vom 18. bis 21. und am 23. und 24. September, in Timpenberg am 22. September, in Wilsede vom 26. September bis 7. Oktober und in Scharnhorst (Eschede) vom 10. bis 13. Oktober.

\par\noindent\hspace*{\parindent}1823 fanden die Beobachtungen statt: in Timpenberg vom 30. Mai bis 2. Juni, in Nindorf vom 3. bis 10. Juni, in Lüneburg vom 11. bis 24. Juni und in Hamburg vom 27. Juni bis 11. Juli. Einige Messungen wurden am 13. Juli in Blankenese angestellt (wohl weil \textsc{Gauss} nach einem Brief an \textsc{Schumacher} vom 20. Juli 1823 daran dachte, sein »Dreieckssystem ganz von der unbequemen und unsicheren Station Michaelis in Hamburg zu trennen«), ferner am 14. Juli in Altona und am 19. Juli auf dem Ägydiusturm in Hannover. Weitere Beobachtungen erfolgten vom 22. August bis 3. September auf der Göttinger Sternwarte, vom 13. bis 27. September auf dem Brocken und vom 5. bis 16. Oktober auf dem Hohehagen.

\par\noindent\hspace*{\parindent}Die Messungen des Jahres 1824 fanden statt: in Falkenberg vom 21. bis 23. Mai, in Elmhorst vom 24. Mai bis 5. Juni, in Bullerberg vom 7. bis 18. Juni, in Bottel (Everser Feld) vom 19. bis 24. Juni, in Brüttendorf vom 28. Juni bis 10. Juli, in Bremen vom 13. Juli bis 20. August, in Garlste (Langeberg) vom 23. bis 25. August, in Brillit vom 27. bis 29. August, in Zeven vom 30. August bis 15. September, in Steinberg vom 17. bis 24. September, in Bottel am 25. September, in Litberg vom 26. September bis 3. Oktober und in Wilsede vom 5. bis 24. Oktober.

\par\noindent\hspace*{\parindent}1825 beobachtete \textsc{Gauss}: in Brüttendorf am 25. April, in Zeven vom 28. April bis 10. Mai, in Bremen vom 12. bis 22. Mai, in Garlste vom 28. Mai bis 5. Juni, in Bremerlehe vom 7. bis 13. Juni, in Varel vom 15. bis 26. Juni, in Langwarden vom 27. Juni bis 12. Juli, in Jever vom 14. bis 19. Juli, in Brillit vom 25. Juli bis 2. August und in Zeven am 4. und 5. August.

\par\noindent\hspace*{\parindent}Die vorstehende Übersicht gibt nur die Dauer der Beobachtungen an. Die Dauer der Aufenthalte auf den Stationen geht daraus nicht hervor.

\par\noindent\hspace*{\parindent}Auf dem Inselsberg erfolgten die Messungen 1821 durch \textsc{Encke}, 1823 durch \textsc{Gerling}.

\par\noindent\hspace*{\parindent}An den Feldarbeiten für die hannoversche Landesvermessung, die von 1828 bis 1844 dauerten, nahm \textsc{Gauss} persönlich nicht teil. Nur einmal, am 7. September 1828, beobachtete er auf der Station Hohehagen. Die Messungen führten die Gehilfen bei der Gradmessung aus: Hauptmann (später Major) \textsc{Müller}, Leutnant (später Hauptmann) \textsc{Hartmann} und Leutnant (später Baurat) \textsc{Joseph Gauss}. Sie waren jedoch nicht dauerhaft für diese Aufgabe abkommandiert, sondern jeweils nur während einiger Sommermonate, soweit sie im militärischen Dienst entbehrlich waren. \textsc{Hartmann} starb bereits 1834, \textsc{Müller} 1843. Die Bearbeitung der Beobachtungen
\par\medskip\noindent\makebox[0pt][l]{IX.}\makebox[.86\linewidth][r]{55}\par

% Printed 434; principal PDF 441; Google PDF 443.
\clearpage
\noindent\makebox[0pt][l]{\fontsize{13}{16}\selectfont 434}\hfill\textbf{BEMERKUNGEN. ZUR HANNOVERSCHEN TRIANGULATION.}\hfill\null\par\medskip
\noindent hatte, wie bereits auf S. 241 mitgeteilt, \textsc{Gauss} allein übernommen. Nur 2 Monate des Winters 1830/1831 unterstützte ihn dabei sein Sohn \textsc{Joseph}. Die rechnerische Bearbeitung wurde erst im Jahr 1848 abgeschlossen. Am 15. März 1848 sandte \textsc{Gauss} dem Innenministerium 35 Hefte mit den Stationsbeobachtungen, zum Teil in Abschrift, 6 Hefte mit Stationsskizzen und 1 Heft mit dem allgemeinen Koordinatenverzeichnis. Diese Hefte befinden sich jetzt im Gauss-Archiv. Das Begleitschreiben dazu ist im Auszug in Band IV, S. 481, abgedruckt.

\par\noindent\hspace*{\parindent}Dieser Sendung lag auf einem einzelnen Blatt das unter [3] wiedergegebene Verzeichnis der geografischen Positionen der Hauptdreieckspunkte bei. Wie aus einem Beobachtungs- und Rechnungsheft zur Gradmessung hervorgeht, wurden Breite und Länge nach den Formeln auf S. 80/81 von Punkt zu Punkt übertragen.

\par\noindent\hspace*{\parindent}Die Zahl der durch die Gradmessung und die Landesvermessung und zum Teil unter Heranziehung von Messungen in den Nachbarstaaten festgelegten Punkte beträgt 2578. Ihre ebenen rechtwinkligen Koordinaten sind im allgemeinen Koordinatenverzeichnis angegeben. Aus den Teilverzeichnissen kommen die Doppelbestimmungen von 411 dieser Punkte hinzu. Siehe Band IV, S. 415/449. \textsc{Gauss} selbst bestimmte über 500 Punkte (vgl. S. 414).
\par\bigskip
\par\noindent\hspace*{\parindent}Es sei noch erwähnt, dass \textsc{Gauss} 1803 aus eigenem Antrieb an eine Vermessung des Herzogtums Braunschweig gedacht und dafür Winkelmessungen ausgeführt hatte. An \textsc{Olbers} schrieb er am 8. April 1803:

\par\noindent\hspace*{\parindent}.\ .\ .\ .\enspace »Ich habe den Plan, einst das ganze Land mit einem Dreiecksnetz zu überziehen, wofür meine jetzigen Messungen nur eine Vorübung sind.«\enspace .\ .\ .\ .

\par\noindent\hspace*{\parindent}Im selben Jahr beteiligte er sich auch an Beobachtungen von Pulversignalen zur Bestimmung von Längen sowie an Breitenbestimmungen für die von \textsc{v. Zach} geleitete trigonometrische und astronomische Vermessung Thüringens (Monatliche Correspondenz, X. Band, S. 302). Für die Vermessung Westfalens, die kurz zuvor unter dem preußischen Generalmajor \textsc{von Lecoq} stattgefunden hatte, lieferte er Berechnungen astronomischer Bestimmungen. Im Aufsatz \textsc{Lecoqs}, Über die trigonometrische Aufnahme in Westphalen (Monatliche Correspondenz, VIII. Band, S. 139), heißt es: »Im astronomischen Teil ist mir Dr. \textsc{Gauss} von großem Nutzen gewesen. Seine Berechnungen und Briefe haben zu meinem Unterricht viel beigetragen.«
\par\bigskip
\par\noindent\hspace*{\parindent}Für die nebenstehende Übersichtskarte der hannoverschen Hauptdreiecke diente eine Karte aus \textsc{Papen}s Topographischem Atlas des Königreichs Hannover und Herzogthums Braunschweig als Vorlage. Auf der Rückseite dieser Karte. die \textsc{Gauss} der letzten Sendung vom 15. März 1848 beigelegt hatte, vermerkte er: »Bei der Farbgebung wurde die Dreiecksseite vom Hohehagen zum Köterberg versehentlich als Teil des Hauptsystems VI gekennzeichnet, obwohl diese Verbindung nur in der preußischen Katastervermessung hergestellt wurde. Stattdessen hätte die Seite Köterberg-Hils als südliche Begrenzung des Systems VI farbig markiert werden sollen. Abgesehen von diesem geringfügigen Umstand finde ich in dieser zweckmäßig angeordneten Übersichtskarte eine getreue Darstellung der Hauptdreiecke und ihrer Verbindungen mit den Messungen in den Nachbarstaaten.« Das Original der Dreieckskarte wurde von Leutnant J. \textsc{Gauss} entworfen. Der genannte Fehler wurde in der nebenstehenden Übersichtskarte berichtigt. Ferner wurde die fehlende Seite Brüttendorf-Litberg eingetragen und die Seite Steinberg-Bottel, die in der Vorlage als einseitig beobachtet eingezeichnet war, durchgezogen. Außerdem war im Original, wie aus einem Bericht und einer Karte von \textsc{Gauss} aus dem Jahr 1843 hervorgeht, das \textsc{Krayenhoff}sche Viereck Emden-Pilsum-Hage-Aurich fälschlich als Teil des ostfriesischen Netzes eingezeichnet worden. In diesem Netz fehlten auch die Seiten Hage-Baltrum und Langeoog-Esens.\hfill\textsc{Krüger}.
\begin{center}\rule{2.6cm}{0.4pt}\end{center}
\par\endgroup

% Unnumbered plate following printed 434.
% Original map devised by Lieutenant J. Gauss. The accompanying signed
% account by Krüger identifies his editorial corrections to this plate.
% Works IX (Leipzig, 1903), SUB Göttingen 4 MATH I, 3253:9.
\clearpage
\begingroup
\edef\HannoverPageWidth{\the\pdfpagewidth}
\edef\HannoverPageHeight{\the\pdfpageheight}
\newgeometry{layoutwidth=420mm,layoutheight=594mm,margin=20mm}
\pdfpagewidth=420mm\relax
\pdfpageheight=594mm\relax
\begin{landscape}
\thispagestyle{empty}
\begingroup
\definecolor{HannoverHaupt}{RGB}{190,94,93}
\definecolor{HannoverWest}{RGB}{77,123,147}
\definecolor{HannoverOst}{RGB}{82,74,53}
\definecolor{HannoverBremen}{RGB}{205,182,76}
\definecolor{HannoverLueneburg}{RGB}{86,119,76}
\definecolor{HannoverMittel}{RGB}{200,104,64}
\definecolor{HannoverEichs}{RGB}{75,66,52}
\definecolor{HannoverHarz}{RGB}{77,111,134}
\definecolor{HannoverOber}{RGB}{167,147,111}
\definecolor{HannoverAller}{RGB}{163,92,91}
\begin{center}
\begin{tikzpicture}[
  x=.14mm,y=-.14mm,
  mapMain/.style={line width=.40pt},
  mapFine/.style={line width=.20pt},
  mapDash/.style={line width=.19pt,dash pattern=on 1.30pt off .67pt},
  mapOutline/.style={draw=black!86,line width=.17pt,line cap=round,line join=round},
  mapBand/.style={line width=.80mm,line cap=round,line join=round},
  mapPoint/.style={draw=black,line width=.35pt,fill=white},
  mapTown/.style={font=\fontsize{8.1}{8.7}\selectfont\bfseries,inner sep=0pt,align=center},
  mapMinor/.style={font=\fontsize{7.3}{8}\selectfont\bfseries\itshape,inner sep=0pt,align=center},
  mapLegend/.style={font=\fontsize{6.5}{7.2}\selectfont\bfseries\itshape,inner sep=0pt,align=left},
  mapTick/.style={font=\fontsize{5.3}{5.9}\selectfont\bfseries,inner sep=0pt},
  mapCaption/.style={font=\fontsize{10.1}{11.2}\selectfont\bfseries,inner sep=0pt}
]
\path[use as bounding box] (32,56) rectangle (3056,2571);
% Latin Modern TS1 slot 115 is the long s (U+017F).
\newcommand{\HannoverLongS}{{\fontencoding{TS1}\fontfamily{lmr}\selectfont\char115\relax}}
\newcommand{\HannoverCircle}[2]{\draw[mapPoint] #1 circle[radius=#2];}
\newcommand{\HannoverCross}[1]{%
  \draw[mapFine] #1 ++(0,-8) -- ++(0,-14);
  \draw[mapFine] #1 ++(-5,-17) -- ++(10,0);
  \fill #1 ++(0,-23) circle[radius=1.4];
  \fill #1 ++(-5,-17) circle[radius=1.15];
  \fill #1 ++(5,-17) circle[radius=1.15];}
\newcommand{\HannoverLight}[1]{%
  \draw[mapFine] #1 ++(0,-8) -- ++(0,-22);
  \draw[mapFine] #1 ++(-8,-21) -- ++(16,0);
  \draw[mapFine] #1 ++(-6,-27) -- ++(12,12);
  \draw[mapFine] #1 ++(6,-27) -- ++(-12,12);}
\newcommand{\HannoverCastle}[1]{%
  \draw[mapFine] #1 ++(0,-8) -- ++(0,-17);
  \fill #1 ++(0,-26) -- ++(-5,11) -- ++(5,-3) -- ++(5,3) -- cycle;}
\newcommand{\HannoverTower}[1]{%
  \draw[mapFine] #1 ++(-3,-8) -- ++(0,-10) -- ++(6,0) -- ++(0,10);
  \draw[mapFine] #1 ++(-7,-18) -- ++(14,0);}

\coordinate (neuwerk) at (1212,160);
\coordinate (wangeroog) at (883,259);
\coordinate (altenwalde) at (1296,244);
\coordinate (borkum) at (267,404);
\coordinate (juist) at (438,334);
\coordinate (norderney) at (544,311);
\coordinate (baltrum) at (624,304);
\coordinate (langeoog) at (687,287);
\coordinate (spiekeroog) at (807,269);
\coordinate (dornum) at (656,371);
\coordinate (esens) at (753,376);
\coordinate (hage) at (582,407);
\coordinate (pilsum) at (461,510);
\coordinate (jever) at (900,449);
\coordinate (langwarden) at (1107,425);
\coordinate (haye) at (1302,429);
\coordinate (bremerlehe) at (1254,466);
\coordinate (holssel) at (1274,374);
\coordinate (krempel) at (1322,336);
\coordinate (bederkesa-i) at (1375,416);
\coordinate (bederkesa) at (1387,422);
\coordinate (wuestenwohlde) at (1423,458);
\coordinate (dolosenberg) at (1512,441);
\coordinate (silberberg) at (1495,332);
\coordinate (basdahl) at (1459,575);
\coordinate (brillit) at (1461,601);
\coordinate (varel) at (1018,604);
\coordinate (garlste) at (1314,725);
\coordinate (westerstede) at (906,718);
\coordinate (aurich) at (681,530);
\coordinate (emden) at (532,604);
\coordinate (leer) at (654,733);
\coordinate (bremen) at (1349,891);
\coordinate (zeven) at (1607,706);
\coordinate (bruettendorf) at (1599,734);
\coordinate (bottel) at (1606,925);
\coordinate (steinberg-n) at (1598,969);
\coordinate (asendorf) at (1455,1161);
\coordinate (twistringen) at (1262,1129);
\coordinate (wildeshausen) at (1164,1043);
\coordinate (oldenburg) at (1054,826);
\coordinate (krapendorf) at (950,1084);
\coordinate (windberg) at (687,1043);
\coordinate (kirchhesepe) at (519,1257);
\coordinate (quekenberg) at (790,1340);
\coordinate (mordkuhlenberg) at (1035,1337);
\coordinate (knickberg) at (1374,1326);
\coordinate (osterberg-w) at (1580,1292);
\coordinate (eckberg) at (1640,1345);
\coordinate (hornhuizen) at (92,567);
\coordinate (uithuiser) at (276,559);
\coordinate (holwierda) at (359,610);
\coordinate (groningen) at (193,718);
\coordinate (midwolde) at (427,753);
\coordinate (onstwedde) at (434,895);
\coordinate (osterwolde) at (40,906);
\coordinate (rolde) at (225,919);
\coordinate (gieten) at (281,907);
\coordinate (beilen) at (154,1034);
\coordinate (ter-apel) at (446,1032);
\coordinate (sleen) at (299,1115);
\coordinate (koeverden) at (260,1216);
\coordinate (lemelerberg) at (76,1367);
\coordinate (uelsen) at (322,1345);
\coordinate (crempe) at (1723,245);
\coordinate (stade) at (1711,442);
\coordinate (hamburg) at (1967,491);
\coordinate (segeberg) at (2144,160);
\coordinate (luebeck) at (2332,223);
\coordinate (syk) at (2131,425);
\coordinate (basis-n) at (2086,417);
\coordinate (basis-s) at (2113,453);
\coordinate (hohenhorn) at (2165,559);
\coordinate (lauenburg) at (2268,643);
\coordinate (litberg) at (1777,632);
\coordinate (wilsede) at (1945,823);
\coordinate (lueneburg) at (2184,750);
\coordinate (bretze) at (2337,738);
\coordinate (glienitz) at (2455,782);
\coordinate (hoe-beck) at (2727,903);
\coordinate (taetendorf) at (2259,925);
\coordinate (hohenmechtin) at (2449,926);
\coordinate (nindorf) at (2122,838);
\coordinate (timpenberg) at (2115,863);
\coordinate (wulfsode) at (2101,909);
\coordinate (elmhorst) at (1785,974);
\coordinate (falkenberg) at (1900,1100);
\coordinate (hauselberg) at (2070,1093);
\coordinate (breithorn) at (2109,1119);
\coordinate (scharnhorst) at (2109,1211);
\coordinate (garssen) at (2049,1266);
\coordinate (celle) at (2015,1298);
\coordinate (holxerberg) at (2226,1074);
\coordinate (wieren) at (2306,1087);
\coordinate (pugelatz) at (2398,1095);
\coordinate (hankensbuettel) at (2281,1209);
\coordinate (bullerberg) at (1674,821);
\coordinate (brelingerberg) at (1802,1337);
\coordinate (osterberg-e) at (1934,1395);
\coordinate (oldenzaal) at (345,1519);
\coordinate (bentheim) at (464,1534);
\coordinate (harikerberg) at (136,1571);
\coordinate (groll) at (165,1737);
\coordinate (ahaus) at (375,1719);
\coordinate (bocholt) at (156,1911);
\coordinate (doerenberg) at (937,1661);
\coordinate (nonnenstein) at (1166,1595);
\coordinate (wittekindstein) at (1384,1614);
\coordinate (huenenburg) at (1159,1811);
\coordinate (koeterberg) at (1607,1959);
\coordinate (hausheide) at (1423,2068);
\coordinate (spitzwarte) at (1130,2237);
\coordinate (desenberg) at (1535,2263);
\coordinate (deisselberg) at (1630,2199);
\coordinate (dommel) at (1236,2404);
\coordinate (hannover) at (1828,1511);
\coordinate (deister-ii) at (1708,1624);
\coordinate (deister-i) at (1756,1636);
\coordinate (osterwald) at (1760,1725);
\coordinate (hils) at (1874,1930);
\coordinate (wohlenberg) at (2183,1447);
\coordinate (lichtenberg) at (2116,1735);
\coordinate (festberg) at (2303,1705);
\coordinate (huyseburg) at (2504,1871);
\coordinate (brocken) at (2295,2013);
\coordinate (bocksberg) at (2153,1963);
\coordinate (fuhrenberg) at (2045,1995);
\coordinate (haspelkopf) at (2158,2079);
\coordinate (eversberg) at (2296,2117);
\coordinate (lange-ecke) at (2265,2138);
\coordinate (schalliethe) at (2449,2156);
\coordinate (tockenburg) at (1965,2132);
\coordinate (wulften) at (2060,2112);
\coordinate (hauseberg) at (2006,2199);
\coordinate (hellberg) at (2094,2209);
\coordinate (bazenberg) at (2079,2267);
\coordinate (sonnenstein) at (2168,2272);
\coordinate (fahlenstollen) at (1783,2104);
\coordinate (weper) at (1873,2128);
\coordinate (meridianzeichen) at (1931,2207);
\coordinate (kleper) at (1949,2246);
\coordinate (goettingen) at (1932,2246);
\coordinate (hohehagen) at (1835,2293);
\coordinate (gieseberg) at (1891,2342);
\coordinate (steinberg-s) at (1797,2382);
\coordinate (meisner) at (1878,2508);
\coordinate (staufenberg) at (1723,2263);

% Outlines, coastlines and waterways traced from the visible plate strokes.
\draw[mapOutline] (1435,136) -- (1437,138) -- (1438,142) -- (1446,150) -- (1449,160) -- (1453,164) -- (1458,166) -- (1461,171) -- (1466,171) -- (1467,170) -- (1474,177) -- (1486,177) -- (1488,176) -- (1493,182) -- (1510,187) -- (1516,191) -- (1531,190) -- (1538,186) -- (1540,191) -- (1560,190) -- (1561,189) -- (1570,191) -- (1579,191) -- (1581,193) -- (1596,195) -- (1605,202) -- (1615,202) -- (1619,204) -- (1633,218) -- (1636,219) -- (1641,224) -- (1644,231) -- (1656,243) -- (1659,250) -- (1659,253) -- (1665,257) -- (1668,257);
\draw[mapOutline] (1501,207) -- (1505,207);
\draw[mapOutline] (1511,207) -- (1512,206) -- (1515,207);
\draw[mapOutline] (1520,207) -- (1533,207);
\draw[mapOutline] (1539,207) -- (1543,207);
\draw[mapOutline] (1545,207) -- (1552,208);
\draw[mapOutline] (1576,207) -- (1581,207) -- (1582,208);
\draw[mapOutline] (1490,208) -- (1481,210);
\draw[mapOutline] (1492,208) -- (1496,208);
\draw[mapOutline] (1558,208) -- (1562,208);
\draw[mapOutline] (1585,208) -- (1589,209);
\draw[mapOutline] (1479,210) -- (1475,211);
\draw[mapOutline] (1594,210) -- (1598,211);
\draw[mapOutline] (1443,211) -- (1438,212);
\draw[mapOutline] (1448,211) -- (1451,210) -- (1456,211);
\draw[mapOutline] (1435,212) -- (1431,214);
\draw[mapOutline] (1466,212) -- (1470,212);
\draw[mapOutline] (1603,213) -- (1608,214);
\draw[mapOutline] (1370,214) -- (1367,214) -- (1365,216);
\draw[mapOutline] (1375,214) -- (1382,214);
\draw[mapOutline] (1384,214) -- (1389,214);
\draw[mapOutline] (1391,214) -- (1392,215) -- (1397,215);
\draw[mapOutline] (1400,215) -- (1410,217);
\draw[mapOutline] (1425,215) -- (1422,215) -- (1418,217) -- (1412,217);
\draw[mapOutline] (1611,216) -- (1627,226);
\draw[mapOutline] (1243,224) -- (1245,222) -- (1245,218) -- (1251,211) -- (1254,210) -- (1259,205) -- (1261,198) -- (1267,191) -- (1273,190) -- (1277,187) -- (1288,186) -- (1290,184) -- (1305,182) -- (1306,185) -- (1304,187) -- (1309,192) -- (1310,196) -- (1313,199) -- (1315,199) -- (1318,202) -- (1323,204) -- (1323,210) -- (1329,215) -- (1334,214) -- (1339,217) -- (1352,229);
\draw[mapOutline] (1292,225) -- (1291,224) -- (1292,221) -- (1296,221) -- (1298,225) -- (1303,225);
\draw[mapOutline] (1331,229) -- (1335,229);
\draw[mapOutline] (1425,229) -- (1434,229) -- (1438,227) -- (1463,226) -- (1468,227) -- (1472,230);
\draw[mapOutline] (1478,229) -- (1485,220) -- (1495,220) -- (1500,217) -- (1505,217) -- (1509,219) -- (1520,219) -- (1530,223) -- (1539,223) -- (1540,224) -- (1541,223) -- (1551,224) -- (1553,222) -- (1572,222) -- (1573,221) -- (1578,221) -- (1581,219) -- (1596,221) -- (1599,223) -- (1606,224) -- (1610,227) -- (1613,227) -- (1629,240) -- (1631,247) -- (1635,252) -- (1634,259) -- (1642,270);
\draw[mapOutline] (1633,233) -- (1639,238);
\draw[mapOutline] (1642,243) -- (1644,246);
\draw[mapOutline] (1646,251) -- (1648,254);
\draw[mapOutline] (912,255) -- (916,257);
\draw[mapOutline] (910,256) -- (902,258);
\draw[mapOutline] (1246,260) -- (1249,270) -- (1260,267);
\draw[mapOutline] (1284,260) -- (1279,264);
\draw[mapOutline] (1651,260) -- (1654,265);
\draw[mapOutline] (786,263) -- (783,266);
\draw[mapOutline] (907,266) -- (916,267);
\draw[mapOutline] (1265,268) -- (1272,269) -- (1275,273);
\draw[mapOutline] (1671,277) -- (1673,288) -- (1681,294) -- (1681,298) -- (1683,300) -- (1683,307) -- (1687,316) -- (1700,331) -- (1708,337) -- (1712,338);
\draw[mapOutline] (778,278) -- (777,279) -- (778,283);
\draw[mapOutline] (1223,278) -- (1220,285) -- (1221,286) -- (1221,293) -- (1219,297) -- (1218,307) -- (1214,316) -- (1207,323) -- (1201,337);
\draw[mapOutline] (749,284) -- (756,286) -- (755,290) -- (758,292) -- (757,293) -- (749,293) -- (747,291) -- (740,289) -- (738,290) -- (723,289);
\draw[mapOutline] (1650,285) -- (1651,289) -- (1657,296) -- (1663,309) -- (1667,313) -- (1667,317) -- (1670,323) -- (1680,331) -- (1682,337) -- (1685,340) -- (1687,340) -- (1691,343) -- (1692,349) -- (1697,354) -- (1698,363) -- (1695,366) -- (1693,366);
\draw[mapOutline] (719,288) -- (700,290);
\draw[mapOutline] (672,299) -- (658,300);
\draw[mapOutline] (551,300) -- (555,302) -- (561,302) -- (562,303) -- (573,303);
\draw[mapOutline] (658,305) -- (657,307) -- (649,311) -- (646,311) -- (643,308);
\draw[mapOutline] (518,307) -- (511,308) -- (509,310) -- (509,312) -- (511,313);
\draw[mapOutline] (603,311) -- (604,312) -- (603,315) -- (598,313) -- (592,315);
\draw[mapOutline] (1675,312) -- (1681,320);
\draw[mapOutline] (1681,322) -- (1692,334);
\draw[mapOutline] (445,323) -- (448,323) -- (450,325);
\draw[mapOutline] (847,329) -- (852,329);
\draw[mapOutline] (852,329) -- (854,327) -- (864,327) -- (872,330) -- (884,329);
\draw[mapOutline] (894,330) -- (895,331) -- (900,330) -- (911,331) -- (921,328) -- (925,330) -- (929,330) -- (932,327) -- (940,327) -- (944,330) -- (951,329) -- (958,332) -- (965,333) -- (968,336) -- (969,341) -- (968,342) -- (965,341) -- (962,344) -- (963,354) -- (961,358) -- (961,364) -- (963,367) -- (965,376) -- (970,383) -- (971,387) -- (970,391);
\draw[mapOutline] (831,332) -- (830,333) -- (816,333) -- (815,332) -- (811,333) -- (806,330) -- (793,330) -- (789,332) -- (780,333);
\draw[mapOutline] (767,335) -- (764,335) -- (760,339) -- (745,339) -- (736,343);
\draw[mapOutline] (415,338) -- (413,339) -- (408,338) -- (404,339) -- (402,338) -- (399,341) -- (394,341) -- (393,342) -- (381,342) -- (380,341) -- (375,341) -- (373,339);
\draw[mapOutline] (1694,338) -- (1702,347);
\draw[mapOutline] (646,339) -- (661,339);
\draw[mapOutline] (636,340) -- (629,339) -- (624,341) -- (620,341) -- (617,338) -- (614,341);
\draw[mapOutline] (673,341) -- (678,340) -- (681,341) -- (684,345);
\draw[mapOutline] (1722,342) -- (1724,344) -- (1728,345) -- (1735,353) -- (1736,357) -- (1740,359) -- (1741,370) -- (1745,374) -- (1744,379) -- (1746,381) -- (1746,385) -- (1743,389) -- (1742,394) -- (1743,397) -- (1748,401) -- (1747,405) -- (1754,411) -- (1755,414) -- (1763,420) -- (1769,421) -- (1777,429) -- (1777,436) -- (1774,438) -- (1774,441) -- (1776,443) -- (1776,445) -- (1779,448) -- (1782,449);
\draw[mapOutline] (593,344) -- (578,342) -- (573,343) -- (566,347) -- (564,346);
\draw[mapOutline] (712,347) -- (714,348) -- (722,348) -- (723,347);
\draw[mapOutline] (695,350) -- (701,353);
\draw[mapOutline] (1706,352) -- (1709,355) -- (1709,357);
\draw[mapOutline] (1687,364) -- (1690,366);
\draw[mapOutline] (1692,368) -- (1697,371) -- (1704,370) -- (1708,373) -- (1710,373);
\draw[mapOutline] (1720,369) -- (1724,376);
\draw[mapOutline] (1169,373) -- (1165,377);
\draw[mapOutline] (1198,373) -- (1201,380) -- (1202,393) -- (1205,400) -- (1208,403) -- (1208,406);
\draw[mapOutline] (1725,379) -- (1729,387) -- (1729,390);
\draw[mapOutline] (1161,381) -- (1155,387);
\draw[mapOutline] (306,383) -- (303,381) -- (300,381) -- (296,377) -- (293,377) -- (289,381) -- (285,383) -- (281,383) -- (277,386);
\draw[mapOutline] (1578,383) -- (1578,386) -- (1580,389);
\draw[mapOutline] (1719,387) -- (1723,393) -- (1723,400) -- (1722,401) -- (1723,404) -- (1727,410) -- (1729,410) -- (1733,413) -- (1734,418) -- (1738,422) -- (1738,427) -- (1744,431) -- (1744,435) -- (1747,438) -- (1747,442) -- (1753,445);
\draw[mapOutline] (654,388) -- (635,391) -- (632,393) -- (627,393) -- (626,394);
\draw[mapOutline] (316,390) -- (311,393);
\draw[mapOutline] (1156,390) -- (1157,394);
\draw[mapOutline] (855,392) -- (858,395) -- (861,396) -- (862,395);
\draw[mapOutline] (311,393) -- (309,396) -- (303,397) -- (301,399) -- (297,398) -- (294,399);
\draw[mapOutline] (1730,393) -- (1734,401);
\draw[mapOutline] (257,394) -- (254,398) -- (253,406);
\draw[mapOutline] (852,394) -- (851,395) -- (852,400);
\draw[mapOutline] (294,399) -- (294,403) -- (291,404);
\draw[mapOutline] (1160,401) -- (1162,405);
\draw[mapOutline] (495,404) -- (490,406) -- (486,413) -- (477,420) -- (475,423) -- (475,432) -- (486,444) -- (488,444) -- (496,452) -- (503,456) -- (507,455) -- (510,451) -- (515,450);
\draw[mapOutline] (253,409) -- (253,415) -- (258,417);
\draw[mapOutline] (870,419) -- (871,424) -- (869,427);
\draw[mapOutline] (1167,423) -- (1171,429) -- (1177,435) -- (1179,435);
\draw[mapOutline] (989,428) -- (993,433);
\draw[mapOutline] (1180,437) -- (1184,440);
\draw[mapOutline] (1076,442) -- (1067,451) -- (1067,462) -- (1064,471) -- (1065,472) -- (1064,477) -- (1065,488) -- (1063,491) -- (1063,496) -- (1056,499) -- (1057,501) -- (1063,501) -- (1064,500);
\draw[mapOutline] (1135,446) -- (1151,461) -- (1157,462) -- (1161,465) -- (1174,466) -- (1180,470) -- (1184,470) -- (1186,472) -- (1198,477) -- (1207,484) -- (1212,485);
\draw[mapOutline] (1001,449) -- (1008,459) -- (1008,464) -- (1010,469) -- (1010,477) -- (1013,480) -- (1013,483) -- (1021,492) -- (1025,500) -- (1026,506) -- (1025,507) -- (1021,506);
\draw[mapOutline] (1786,449) -- (1790,452);
\draw[mapOutline] (867,450) -- (864,454) -- (862,454);
\draw[mapOutline] (862,454) -- (860,456) -- (860,458);
\draw[mapOutline] (255,457) -- (265,467);
\draw[mapOutline] (1793,463) -- (1805,465) -- (1810,467);
\draw[mapOutline] (1789,465) -- (1805,478);
\draw[mapOutline] (1215,466) -- (1218,468);
\draw[mapOutline] (1810,467) -- (1813,471) -- (1818,469);
\draw[mapOutline] (1810,467) -- (1811,466) -- (1812,467) -- (1817,467);
\draw[mapOutline] (1822,468) -- (1818,469);
\draw[mapOutline] (1818,469) -- (1822,471);
\draw[mapOutline] (856,471) -- (859,473);
\draw[mapOutline] (1788,471) -- (1793,475) -- (1794,478) -- (1802,487) -- (1804,491) -- (1807,494) -- (1810,495) -- (1814,500) -- (1822,500) -- (1826,502) -- (1831,502) -- (1836,500) -- (1840,496) -- (1853,495) -- (1858,497) -- (1861,500) -- (1863,500) -- (1866,503) -- (1872,502) -- (1876,507) -- (1881,508) -- (1884,511) -- (1885,516) -- (1892,521) -- (1898,521) -- (1904,518);
\draw[mapOutline] (1222,472) -- (1226,476);
\draw[mapOutline] (485,475) -- (481,479);
\draw[mapOutline] (2167,475) -- (2163,478) -- (2159,478) -- (2156,480) -- (2152,486);
\draw[mapOutline] (859,476) -- (859,484) -- (858,485);
\draw[mapOutline] (1831,476) -- (1843,477) -- (1844,478) -- (1849,478) -- (1850,477) -- (1858,478) -- (1864,482) -- (1873,483) -- (1883,489) -- (1888,489) -- (1889,488) -- (1894,491) -- (1901,491) -- (1904,493) -- (1907,493) -- (1911,496) -- (1920,496) -- (1924,498) -- (1932,497) -- (1934,499) -- (1943,499);
\draw[mapOutline] (1859,476) -- (1864,476) -- (1874,479) -- (1879,479) -- (1895,483) -- (1900,483) -- (1921,488) -- (1926,488) -- (1951,493);
\draw[mapOutline] (515,479) -- (518,480) -- (517,494) -- (519,496) -- (519,501);
\draw[mapOutline] (279,480) -- (288,489) -- (291,490);
\draw[mapOutline] (479,480) -- (475,484);
\draw[mapOutline] (1238,481) -- (1245,494) -- (1245,498) -- (1247,502) -- (1249,525) -- (1242,533) -- (1240,533) -- (1238,536) -- (1236,536);
\draw[mapOutline] (1810,481) -- (1813,484) -- (1822,486);
\draw[mapOutline] (1232,484) -- (1235,489) -- (1237,497);
\draw[mapOutline] (2033,484) -- (2027,498) -- (2023,501) -- (2022,504);
\draw[mapOutline] (472,486) -- (468,490) -- (464,489);
\draw[mapOutline] (1827,486) -- (1831,486);
\draw[mapOutline] (1836,487) -- (1843,487);
\draw[mapOutline] (1851,488) -- (1858,490);
\draw[mapOutline] (855,489) -- (851,493);
\draw[mapOutline] (2403,489) -- (2405,490) -- (2406,495);
\draw[mapOutline] (1223,491) -- (1229,497) -- (1229,503) -- (1226,506) -- (1224,512) -- (1217,514) -- (1210,522) -- (1208,522) -- (1201,527) -- (1197,534) -- (1197,540) -- (1193,544) -- (1187,555) -- (1185,567) -- (1186,568) -- (1187,590) -- (1189,595);
\draw[mapOutline] (2145,491) -- (2140,495);
\draw[mapOutline] (293,492) -- (297,495) -- (303,497);
\draw[mapOutline] (1863,492) -- (1867,493);
\draw[mapOutline] (1869,494) -- (1877,498);
\draw[mapOutline] (850,496) -- (850,498) -- (853,500);
\draw[mapOutline] (2406,497) -- (2398,507) -- (2398,509);
\draw[mapOutline] (879,498) -- (887,496) -- (893,501) -- (896,500);
\draw[mapOutline] (926,498) -- (921,499) -- (918,501) -- (917,505) -- (915,507) -- (909,505);
\draw[mapOutline] (305,499) -- (309,502);
\draw[mapOutline] (492,499) -- (500,501) -- (506,501) -- (507,500) -- (509,501) -- (515,501) -- (516,500) -- (519,501);
\draw[mapOutline] (877,499) -- (872,500);
\draw[mapOutline] (1879,499) -- (1884,503);
\draw[mapOutline] (1237,500) -- (1240,507);
\draw[mapOutline] (870,501) -- (864,505);
\draw[mapOutline] (1074,502) -- (1087,503) -- (1091,498) -- (1095,498) -- (1100,502) -- (1099,506) -- (1100,507) -- (1107,508) -- (1109,511);
\draw[mapOutline] (853,503) -- (852,505) -- (854,510) -- (860,509);
\draw[mapOutline] (1890,504) -- (1898,503) -- (1899,504) -- (1904,504) -- (1907,506) -- (1916,505) -- (1918,507) -- (1922,507) -- (1925,512) -- (1928,512) -- (1931,508) -- (1936,507);
\draw[mapOutline] (1019,505) -- (1008,507) -- (1005,512) -- (1001,512) -- (997,509) -- (994,509) -- (991,512) -- (987,512) -- (985,510) -- (985,508) -- (983,506) -- (981,506) -- (976,509) -- (975,513) -- (970,521);
\draw[mapOutline] (1019,505) -- (1020,504) -- (1021,506);
\draw[mapOutline] (1886,505) -- (1888,513) -- (1891,516) -- (1903,516);
\draw[mapOutline] (314,506) -- (328,515);
\draw[mapOutline] (1889,506) -- (1891,509) -- (1899,511);
\draw[mapOutline] (2029,507) -- (2033,509);
\draw[mapOutline] (1239,510) -- (1236,517);
\draw[mapOutline] (1914,511) -- (1908,512);
\draw[mapOutline] (1901,512) -- (1908,512);
\draw[mapOutline] (2135,513) -- (2128,520) -- (2125,520) -- (2120,515) -- (2118,515) -- (2116,517) -- (2109,534) -- (2107,534) -- (2102,529) -- (2099,530);
\draw[mapOutline] (2390,514) -- (2384,516);
\draw[mapOutline] (1110,515) -- (1110,524) -- (1108,528) -- (1108,542) -- (1109,544);
\draw[mapOutline] (1908,515) -- (1903,516);
\draw[mapOutline] (1908,515) -- (1911,515) -- (1915,519) -- (1915,521);
\draw[mapOutline] (333,517) -- (336,520) -- (349,525) -- (355,531);
\draw[mapOutline] (1237,518) -- (1240,517) -- (1242,519) -- (1242,521) -- (1239,527) -- (1234,532);
\draw[mapOutline] (1904,518) -- (1909,519) -- (1911,521) -- (1911,525);
\draw[mapOutline] (2010,519) -- (2008,520) -- (2008,524) -- (2005,528);
\draw[mapOutline] (1952,520) -- (1958,520);
\draw[mapOutline] (438,521) -- (437,526);
\draw[mapOutline] (1229,522) -- (1220,523) -- (1207,533) -- (1204,542) -- (1198,550) -- (1197,555) -- (1199,558) -- (1201,558) -- (1209,550) -- (1220,542) -- (1223,538) -- (1233,533);
\draw[mapOutline] (1229,522) -- (1230,527);
\draw[mapOutline] (2386,525) -- (2382,533);
\draw[mapOutline] (284,529) -- (291,529) -- (293,527) -- (295,527) -- (297,529) -- (303,530) -- (305,528) -- (314,525) -- (320,525) -- (321,526) -- (328,525) -- (332,526) -- (338,533) -- (341,533) -- (343,535);
\draw[mapOutline] (2047,531) -- (2051,532);
\draw[mapOutline] (1234,533) -- (1236,536);
\draw[mapOutline] (2380,533) -- (2375,536);
\draw[mapOutline] (1939,536) -- (1943,539) -- (1946,539) -- (1949,537);
\draw[mapOutline] (1236,537) -- (1227,541) -- (1222,545) -- (1216,552) -- (1208,559) -- (1205,560) -- (1199,568) -- (1200,592) -- (1197,595) -- (1197,598);
\draw[mapOutline] (1954,537) -- (1957,540) -- (1958,543) -- (1958,549) -- (1957,550) -- (1958,551);
\draw[mapOutline] (940,538) -- (936,543);
\draw[mapOutline] (2374,538) -- (2369,547);
\draw[mapOutline] (438,539) -- (435,542) -- (434,548) -- (432,551) -- (432,557) -- (430,560) -- (431,572) -- (428,575) -- (429,587) -- (427,591) -- (429,594) -- (430,600) -- (433,602);
\draw[mapOutline] (1995,539) -- (1997,540) -- (1997,544);
\draw[mapOutline] (358,542) -- (356,544) -- (359,548) -- (358,558) -- (365,565) -- (364,574) -- (360,579);
\draw[mapOutline] (1236,544) -- (1237,550);
\draw[mapOutline] (2083,544) -- (2082,545) -- (2076,545);
\draw[mapOutline] (364,545) -- (365,550);
\draw[mapOutline] (2076,545) -- (2076,548) -- (2079,551) -- (2082,552) -- (2082,556) -- (2078,563) -- (2069,561) -- (2064,568) -- (2059,568) -- (2058,569) -- (2054,569) -- (2051,566) -- (2046,564) -- (2043,565) -- (2040,569);
\draw[mapOutline] (936,546) -- (929,548);
\draw[mapOutline] (2361,547) -- (2364,546) -- (2366,548);
\draw[mapOutline] (1917,548) -- (1923,554);
\draw[mapOutline] (2092,548) -- (2100,556);
\draw[mapOutline] (926,550) -- (927,551) -- (925,553) -- (924,557);
\draw[mapOutline] (1238,553) -- (1238,558) -- (1237,559) -- (1238,567) -- (1236,571) -- (1237,579) -- (1235,581) -- (1233,581);
\draw[mapOutline] (2358,553) -- (2356,555) -- (2353,553) -- (2350,555);
\draw[mapOutline] (366,554) -- (366,557) -- (369,562);
\draw[mapOutline] (1010,554) -- (1006,557) -- (1002,557) -- (1000,559) -- (1001,563) -- (1005,567);
\draw[mapOutline] (1010,554) -- (1010,552) -- (1012,551) -- (1013,554) -- (1011,555);
\draw[mapOutline] (1956,554) -- (1954,559) -- (1951,562);
\draw[mapOutline] (2347,554) -- (2344,555) -- (2345,562) -- (2343,567);
\draw[mapOutline] (1011,555) -- (1011,559) -- (1019,564) -- (1018,571) -- (1022,574);
\draw[mapOutline] (1927,556) -- (1932,558);
\draw[mapOutline] (1988,556) -- (1982,559) -- (1981,561) -- (1991,563);
\draw[mapOutline] (1940,557) -- (1947,562) -- (1948,561);
\draw[mapOutline] (1938,558) -- (1934,559);
\draw[mapOutline] (2104,558) -- (2109,557) -- (2113,561);
\draw[mapOutline] (1104,559) -- (1106,562);
\draw[mapOutline] (924,562) -- (914,564) -- (908,571);
\draw[mapOutline] (1106,562) -- (1102,565) -- (1098,565) -- (1096,567) -- (1096,569) -- (1092,571) -- (1092,575) -- (1096,576) -- (1097,580) -- (1095,585) -- (1089,591) -- (1085,598) -- (1082,600) -- (1077,600) -- (1070,607) -- (1065,607) -- (1064,606) -- (1061,607);
\draw[mapOutline] (370,566) -- (371,572) -- (369,575) -- (369,580) -- (368,581) -- (368,592);
\draw[mapOutline] (2006,567) -- (2014,570) -- (2019,576) -- (2022,578) -- (2026,578) -- (2028,581);
\draw[mapOutline] (2323,570) -- (2317,575) -- (2309,574) -- (2306,580);
\draw[mapOutline] (2335,573) -- (2328,575);
\draw[mapOutline] (1229,576) -- (1232,578);
\draw[mapOutline] (2325,576) -- (2321,579) -- (2324,582) -- (2327,583);
\draw[mapOutline] (898,578) -- (888,583);
\draw[mapOutline] (1223,582) -- (1228,582);
\draw[mapOutline] (1230,582) -- (1233,586) -- (1233,589);
\draw[mapOutline] (2305,582) -- (2304,584) -- (2305,587);
\draw[mapOutline] (1034,583) -- (1038,587);
\draw[mapOutline] (2340,583) -- (2336,586) -- (2330,585);
\draw[mapOutline] (886,584) -- (881,587);
\draw[mapOutline] (2029,584) -- (2032,591) -- (2037,595) -- (2039,595);
\draw[mapOutline] (875,589) -- (873,591) -- (871,591);
\draw[mapOutline] (2304,590) -- (2301,596);
\draw[mapOutline] (2040,591) -- (2039,595);
\draw[mapOutline] (1206,593) -- (1211,593) -- (1213,591) -- (1217,590) -- (1219,591) -- (1219,593);
\draw[mapOutline] (2082,593) -- (2084,596) -- (2080,600) -- (2078,608) -- (2074,611);
\draw[mapOutline] (2148,593) -- (2151,597) -- (2154,598);
\draw[mapOutline] (2177,593) -- (2179,594) -- (2179,597);
\draw[mapOutline] (2039,595) -- (2041,597) -- (2045,597) -- (2049,599) -- (2056,610) -- (2062,616) -- (2069,620);
\draw[mapOutline] (2174,596) -- (2179,597);
\draw[mapOutline] (1220,597) -- (1225,598);
\draw[mapOutline] (2179,597) -- (2208,624) -- (2214,625);
\draw[mapOutline] (874,598) -- (878,604);
\draw[mapOutline] (2300,598) -- (2300,603) -- (2301,604);
\draw[mapOutline] (2122,600) -- (2121,601) -- (2114,601);
\draw[mapOutline] (1003,601) -- (1003,605);
\draw[mapOutline] (1196,602) -- (1196,606);
\draw[mapOutline] (2112,602) -- (2106,607) -- (2104,614);
\draw[mapOutline] (880,607) -- (884,614) -- (886,615);
\draw[mapOutline] (1197,610) -- (1196,616);
\draw[mapOutline] (503,611) -- (498,615) -- (492,623) -- (490,623) -- (486,620) -- (473,619) -- (469,617) -- (462,622) -- (458,621) -- (450,628) -- (446,629) -- (439,629) -- (438,628) -- (439,626) -- (438,612);
\draw[mapOutline] (372,615) -- (384,625);
\draw[mapOutline] (887,618) -- (892,623);
\draw[mapOutline] (369,619) -- (382,631) -- (384,635) -- (384,640) -- (386,642) -- (389,643) -- (398,650) -- (407,649) -- (421,660) -- (426,660) -- (434,664) -- (440,662) -- (445,665) -- (450,665) -- (451,666) -- (456,666) -- (461,661) -- (466,661) -- (467,660) -- (470,661) -- (468,665) -- (459,672) -- (457,678) -- (453,679) -- (446,685) -- (446,690) -- (440,696) -- (443,708) -- (445,710) -- (447,710) -- (448,713);
\draw[mapOutline] (1197,619) -- (1198,629);
\draw[mapOutline] (2103,620) -- (2099,624) -- (2090,622);
\draw[mapOutline] (514,621) -- (512,623) -- (512,625);
\draw[mapOutline] (2216,625) -- (2222,627);
\draw[mapOutline] (893,627) -- (895,630);
\draw[mapOutline] (2224,628) -- (2227,633);
\draw[mapOutline] (388,629) -- (390,632);
\draw[mapOutline] (1199,632) -- (1199,636);
\draw[mapOutline] (392,636) -- (397,639);
\draw[mapOutline] (892,636) -- (890,639);
\draw[mapOutline] (534,637) -- (540,645) -- (544,645) -- (547,643) -- (548,644) -- (557,644) -- (559,646) -- (562,646) -- (565,648);
\draw[mapOutline] (1200,639) -- (1201,646);
\draw[mapOutline] (399,640) -- (403,643) -- (405,643);
\draw[mapOutline] (2241,641) -- (2245,642);
\draw[mapOutline] (2337,643) -- (2339,642) -- (2341,646);
\draw[mapOutline] (411,644) -- (413,644) -- (425,655) -- (432,658) -- (437,658);
\draw[mapOutline] (2298,644) -- (2304,644);
\draw[mapOutline] (2307,644) -- (2312,644);
\draw[mapOutline] (2317,644) -- (2325,646) -- (2333,644);
\draw[mapOutline] (2294,645) -- (2290,647);
\draw[mapOutline] (887,646) -- (879,658) -- (876,660);
\draw[mapOutline] (2287,648) -- (2282,648) -- (2281,649);
\draw[mapOutline] (2343,649) -- (2349,652);
\draw[mapOutline] (556,652) -- (554,654) -- (551,654) -- (548,657) -- (547,671) -- (546,672) -- (541,671) -- (538,676) -- (538,679) -- (541,680) -- (545,685) -- (545,694) -- (544,695) -- (543,706) -- (538,710) -- (538,713) -- (530,724) -- (528,724) -- (525,727) -- (520,727);
\draw[mapOutline] (466,653) -- (470,654);
\draw[mapOutline] (2352,653) -- (2353,652) -- (2355,653) -- (2356,658) -- (2361,658) -- (2363,660) -- (2363,663);
\draw[mapOutline] (475,654) -- (479,654);
\draw[mapOutline] (481,654) -- (482,655) -- (488,655);
\draw[mapOutline] (442,656) -- (446,657) -- (449,660) -- (454,660) -- (455,659);
\draw[mapOutline] (504,657) -- (508,657);
\draw[mapOutline] (2452,659) -- (2453,658) -- (2457,660);
\draw[mapOutline] (2451,661) -- (2448,663) -- (2447,666);
\draw[mapOutline] (2460,661) -- (2463,663);
\draw[mapOutline] (874,664) -- (872,667);
\draw[mapOutline] (2463,666) -- (2462,667) -- (2463,668) -- (2465,668) -- (2472,672) -- (2475,671);
\draw[mapOutline] (2445,669) -- (2442,672) -- (2439,672);
\draw[mapOutline] (869,671) -- (865,676);
\draw[mapOutline] (2475,671) -- (2481,671);
\draw[mapOutline] (2436,672) -- (2431,672) -- (2429,674) -- (2430,680);
\draw[mapOutline] (2486,673) -- (2489,675) -- (2488,679) -- (2486,681) -- (2489,683);
\draw[mapOutline] (2379,682) -- (2381,685) -- (2384,685);
\draw[mapOutline] (858,684) -- (856,688) -- (854,689);
\draw[mapOutline] (2387,687) -- (2393,692) -- (2400,693);
\draw[mapOutline] (1198,690) -- (1199,691) -- (1197,696);
\draw[mapOutline] (851,692) -- (846,697);
\draw[mapOutline] (2406,694) -- (2412,695);
\draw[mapOutline] (2414,694) -- (2418,695);
\draw[mapOutline] (2420,695) -- (2425,698);
\draw[mapOutline] (518,697) -- (519,701);
\draw[mapOutline] (2437,698) -- (2439,704);
\draw[mapOutline] (845,700) -- (834,711);
\draw[mapOutline] (1197,700) -- (1197,706);
\draw[mapOutline] (2425,701) -- (2425,707) -- (2427,708);
\draw[mapOutline] (520,706) -- (520,710);
\draw[mapOutline] (2442,706) -- (2448,707);
\draw[mapOutline] (2429,709) -- (2437,716) -- (2440,716);
\draw[mapOutline] (2446,710) -- (2441,714);
\draw[mapOutline] (832,713) -- (828,717);
\draw[mapOutline] (454,724) -- (459,724) -- (460,728) -- (464,732) -- (468,733) -- (471,736) -- (475,736) -- (480,740);
\draw[mapOutline] (2511,725) -- (2512,728) -- (2515,731);
\draw[mapOutline] (520,729) -- (516,729) -- (511,733) -- (501,734);
\draw[mapOutline] (520,729) -- (521,735);
\draw[mapOutline] (815,731) -- (809,735);
\draw[mapOutline] (805,737) -- (798,740);
\draw[mapOutline] (796,740) -- (788,743) -- (778,752);
\draw[mapOutline] (2517,741) -- (2519,746);
\draw[mapOutline] (1186,742) -- (1185,746);
\draw[mapOutline] (2842,751) -- (2846,751);
\draw[mapOutline] (2830,752) -- (2831,753) -- (2832,752) -- (2837,752);
\draw[mapOutline] (1184,753) -- (1186,758);
\draw[mapOutline] (2828,753) -- (2823,755) -- (2818,760);
\draw[mapOutline] (779,755) -- (780,759);
\draw[mapOutline] (523,756) -- (524,760);
\draw[mapOutline] (1188,760) -- (1194,767);
\draw[mapOutline] (524,763) -- (521,767) -- (522,769);
\draw[mapOutline] (2815,763) -- (2811,767);
\draw[mapOutline] (2807,769) -- (2804,772);
\draw[mapOutline] (527,773) -- (526,776) -- (522,780) -- (518,778);
\draw[mapOutline] (1197,773) -- (1199,778);
\draw[mapOutline] (784,776) -- (785,780);
\draw[mapOutline] (2801,777) -- (2799,780) -- (2796,780);
\draw[mapOutline] (747,778) -- (749,781);
\draw[mapOutline] (2792,778) -- (2783,787) -- (2779,788);
\draw[mapOutline] (746,780) -- (745,781) -- (746,786);
\draw[mapOutline] (754,781) -- (759,787) -- (762,788);
\draw[mapOutline] (1201,781) -- (1203,785);
\draw[mapOutline] (2541,782) -- (2546,783);
\draw[mapOutline] (786,784) -- (781,788) -- (780,787);
\draw[mapOutline] (2548,784) -- (2554,788);
\draw[mapOutline] (770,785) -- (767,788);
\draw[mapOutline] (773,785) -- (778,788);
\draw[mapOutline] (1579,788) -- (1579,794);
\draw[mapOutline] (2558,789) -- (2559,790) -- (2560,789) -- (2566,789);
\draw[mapOutline] (2569,789) -- (2570,795) -- (2573,798) -- (2576,799);
\draw[mapOutline] (512,790) -- (511,795);
\draw[mapOutline] (1209,790) -- (1213,794) -- (1216,795);
\draw[mapOutline] (746,791) -- (745,795);
\draw[mapOutline] (1219,797) -- (1224,800);
\draw[mapOutline] (2582,800) -- (2586,800);
\draw[mapOutline] (511,802) -- (510,808);
\draw[mapOutline] (2589,802) -- (2591,805);
\draw[mapOutline] (1230,803) -- (1234,804);
\draw[mapOutline] (1241,803) -- (1246,803);
\draw[mapOutline] (1252,803) -- (1259,806);
\draw[mapOutline] (2769,803) -- (2764,808) -- (2761,806) -- (2757,807);
\draw[mapOutline] (1773,805) -- (1775,808) -- (1775,810) -- (1773,811);
\draw[mapOutline] (2742,807) -- (2738,808);
\draw[mapOutline] (2745,807) -- (2746,808) -- (2754,808);
\draw[mapOutline] (1318,809) -- (1311,817) -- (1308,816) -- (1305,821);
\draw[mapOutline] (1322,809) -- (1325,812) -- (1327,812);
\draw[mapOutline] (1261,810) -- (1263,813);
\draw[mapOutline] (2736,811) -- (2732,816) -- (2730,814);
\draw[mapOutline] (745,812) -- (745,816);
\draw[mapOutline] (2594,812) -- (2595,816);
\draw[mapOutline] (2730,814) -- (2729,814) -- (2728,819);
\draw[mapOutline] (1358,816) -- (1362,815) -- (1364,817);
\draw[mapOutline] (1267,817) -- (1272,818);
\draw[mapOutline] (1290,817) -- (1296,820);
\draw[mapOutline] (1369,818) -- (1364,819);
\draw[mapOutline] (1284,819) -- (1280,820);
\draw[mapOutline] (2594,820) -- (2591,826);
\draw[mapOutline] (506,821) -- (511,826);
\draw[mapOutline] (2726,822) -- (2725,825) -- (2722,828) -- (2717,826);
\draw[mapOutline] (744,824) -- (744,828);
\draw[mapOutline] (2589,827) -- (2587,829) -- (2591,834) -- (2591,836);
\draw[mapOutline] (510,829) -- (512,832);
\draw[mapOutline] (1260,829) -- (1261,836);
\draw[mapOutline] (1377,831) -- (1381,839) -- (1384,842) -- (1388,843);
\draw[mapOutline] (2611,832) -- (2618,835);
\draw[mapOutline] (2607,833) -- (2603,834) -- (2598,839);
\draw[mapOutline] (1417,834) -- (1421,834);
\draw[mapOutline] (2704,835) -- (2702,835) -- (2700,837) -- (2701,842);
\draw[mapOutline] (1405,841) -- (1400,843);
\draw[mapOutline] (1390,844) -- (1398,844);
\draw[mapOutline] (2619,844) -- (2626,847);
\draw[mapOutline] (520,846) -- (519,850);
\draw[mapOutline] (1267,846) -- (1271,849);
\draw[mapOutline] (1445,846) -- (1447,850);
\draw[mapOutline] (738,854) -- (737,858);
\draw[mapOutline] (2637,855) -- (2643,859) -- (2644,858);
\draw[mapOutline] (2690,855) -- (2686,857) -- (2684,856);
\draw[mapOutline] (1271,856) -- (1272,860);
\draw[mapOutline] (518,857) -- (517,862);
\draw[mapOutline] (1446,857) -- (1451,859) -- (1454,863);
\draw[mapOutline] (2662,858) -- (2665,858) -- (2670,861);
\draw[mapOutline] (2681,858) -- (2681,860) -- (2678,861);
\draw[mapOutline] (2646,860) -- (2650,861) -- (2660,860);
\draw[mapOutline] (2672,861) -- (2676,861);
\draw[mapOutline] (1274,862) -- (1281,865);
\draw[mapOutline] (735,864) -- (734,868);
\draw[mapOutline] (1282,868) -- (1287,872);
\draw[mapOutline] (517,870) -- (517,874);
\draw[mapOutline] (732,874) -- (731,880);
\draw[mapOutline] (517,881) -- (518,885);
\draw[mapOutline] (2645,885) -- (2651,888);
\draw[mapOutline] (1299,889) -- (1301,892);
\draw[mapOutline] (2657,890) -- (2663,893);
\draw[mapOutline] (518,892) -- (519,896);
\draw[mapOutline] (730,893) -- (735,897);
\draw[mapOutline] (738,899) -- (744,903);
\draw[mapOutline] (1301,899) -- (1301,906) -- (1300,907) -- (1300,912) -- (1301,913);
\draw[mapOutline] (2746,899) -- (2749,900) -- (2763,913) -- (2767,914);
\draw[mapOutline] (519,902) -- (518,907);
\draw[mapOutline] (746,905) -- (751,909);
\draw[mapOutline] (754,911) -- (758,914) -- (760,914);
\draw[mapOutline] (2781,911) -- (2785,911);
\draw[mapOutline] (517,912) -- (517,919);
\draw[mapOutline] (2678,912) -- (2682,913);
\draw[mapOutline] (2776,913) -- (2772,914);
\draw[mapOutline] (2790,915) -- (2792,918);
\draw[mapOutline] (762,916) -- (764,918) -- (766,918);
\draw[mapOutline] (1389,919) -- (1384,921);
\draw[mapOutline] (1413,919) -- (1411,920) -- (1413,922) -- (1417,922) -- (1420,924) -- (1421,923) -- (1427,924);
\draw[mapOutline] (1332,921) -- (1329,923) -- (1330,924);
\draw[mapOutline] (1382,921) -- (1381,925) -- (1383,927) -- (1383,929) -- (1386,930) -- (1391,936);
\draw[mapOutline] (2797,921) -- (2802,921);
\draw[mapOutline] (516,924) -- (519,932);
\draw[mapOutline] (772,925) -- (777,928);
\draw[mapOutline] (2804,925) -- (2804,927) -- (2799,930);
\draw[mapOutline] (780,930) -- (778,933) -- (778,937);
\draw[mapOutline] (1354,930) -- (1363,939);
\draw[mapOutline] (1385,935) -- (1387,937) -- (1391,936);
\draw[mapOutline] (1374,938) -- (1379,942) -- (1381,942) -- (1382,940);
\draw[mapOutline] (777,939) -- (784,942);
\draw[mapOutline] (2796,939) -- (2795,942) -- (2796,945);
\draw[mapOutline] (795,940) -- (800,940);
\draw[mapOutline] (1262,940) -- (1259,945);
\draw[mapOutline] (1265,940) -- (1272,946);
\draw[mapOutline] (792,942) -- (793,943) -- (792,944) -- (787,943);
\draw[mapOutline] (811,942) -- (807,943);
\draw[mapOutline] (812,944) -- (815,946);
\draw[mapOutline] (1256,947) -- (1255,951);
\draw[mapOutline] (2795,948) -- (2794,953);
\draw[mapOutline] (1309,949) -- (1313,954) -- (1316,954) -- (1318,952) -- (1320,953);
\draw[mapOutline] (821,950) -- (825,953);
\draw[mapOutline] (1260,953) -- (1261,957) -- (1258,959);
\draw[mapOutline] (512,954) -- (510,959);
\draw[mapOutline] (1447,954) -- (1444,953) -- (1440,956) -- (1439,955);
\draw[mapOutline] (1446,956) -- (1445,960);
\draw[mapOutline] (2792,957) -- (2792,959) -- (2788,962);
\draw[mapOutline] (1437,959) -- (1436,963) -- (1432,968) -- (1432,970);
\draw[mapOutline] (1256,960) -- (1252,962);
\draw[mapOutline] (2765,960) -- (2763,963);
\draw[mapOutline] (827,963) -- (830,969);
\draw[mapOutline] (1247,965) -- (1233,976);
\draw[mapOutline] (508,966) -- (507,971);
\draw[mapOutline] (1496,966) -- (1494,965) -- (1484,969) -- (1480,968);
\draw[mapOutline] (1496,966) -- (1496,970) -- (1494,974);
\draw[mapOutline] (1473,967) -- (1478,967);
\draw[mapOutline] (1471,968) -- (1468,975) -- (1466,976);
\draw[mapOutline] (1450,969) -- (1453,968) -- (1457,969);
\draw[mapOutline] (1457,971) -- (1464,976);
\draw[mapOutline] (832,972) -- (831,977) -- (834,981);
\draw[mapOutline] (1449,972) -- (1446,975) -- (1443,974);
\draw[mapOutline] (1437,975) -- (1439,978) -- (1439,981);
\draw[mapOutline] (1230,976) -- (1224,970) -- (1220,969) -- (1218,978) -- (1216,979) -- (1213,977) -- (1211,979) -- (1211,981);
\draw[mapOutline] (505,978) -- (503,984);
\draw[mapOutline] (1491,978) -- (1491,981) -- (1493,984);
\draw[mapOutline] (1439,981) -- (1440,982) -- (1439,983) -- (1438,982);
\draw[mapOutline] (833,983) -- (834,987);
\draw[mapOutline] (1211,983) -- (1208,986) -- (1204,987) -- (1201,986);
\draw[mapOutline] (1434,985) -- (1439,990) -- (1441,997) -- (1437,1001) -- (1437,1006) -- (1436,1007) -- (1437,1008);
\draw[mapOutline] (502,988) -- (499,991);
\draw[mapOutline] (1200,989) -- (1200,993) -- (1196,997) -- (1190,997) -- (1189,998) -- (1189,1003) -- (1178,1015);
\draw[mapOutline] (834,990) -- (837,995);
\draw[mapOutline] (1492,992) -- (1489,994) -- (1489,997);
\draw[mapOutline] (494,995) -- (492,998);
\draw[mapOutline] (2756,996) -- (2756,1001);
\draw[mapOutline] (1485,998) -- (1481,1000);
\draw[mapOutline] (490,1001) -- (480,1016);
\draw[mapOutline] (837,1001) -- (836,1006);
\draw[mapOutline] (2759,1003) -- (2767,1011) -- (2769,1011);
\draw[mapOutline] (1480,1007) -- (1475,1009);
\draw[mapOutline] (1434,1011) -- (1430,1016);
\draw[mapOutline] (1470,1011) -- (1463,1015) -- (1462,1014);
\draw[mapOutline] (2758,1014) -- (2755,1016);
\draw[mapOutline] (2763,1014) -- (2766,1015) -- (2768,1014);
\draw[mapOutline] (1460,1016) -- (1454,1021) -- (1451,1027);
\draw[mapOutline] (1427,1017) -- (1422,1022) -- (1419,1022) -- (1415,1025);
\draw[mapOutline] (477,1018) -- (472,1026) -- (465,1033);
\draw[mapOutline] (1176,1018) -- (1176,1025) -- (1175,1026);
\draw[mapOutline] (2753,1018) -- (2752,1019) -- (2747,1019);
\draw[mapOutline] (2745,1019) -- (2740,1018) -- (2737,1021) -- (2735,1025) -- (2735,1028) -- (2731,1032) -- (2726,1032) -- (2719,1037) -- (2718,1039) -- (2720,1042);
\draw[mapOutline] (1415,1028) -- (1416,1030) -- (1420,1031);
\draw[mapOutline] (1451,1033) -- (1451,1037);
\draw[mapOutline] (461,1038) -- (458,1042);
\draw[mapOutline] (2513,1041) -- (2514,1040) -- (2518,1041);
\draw[mapOutline] (2545,1041) -- (2546,1044) -- (2548,1045);
\draw[mapOutline] (2521,1042) -- (2524,1044) -- (2527,1044);
\draw[mapOutline] (2720,1042) -- (2719,1043) -- (2720,1044) -- (2721,1043);
\draw[mapOutline] (2503,1043) -- (2507,1044);
\draw[mapOutline] (2534,1043) -- (2530,1044);
\draw[mapOutline] (2717,1044) -- (2713,1046);
\draw[mapOutline] (456,1045) -- (451,1050);
\draw[mapOutline] (2584,1045) -- (2588,1045) -- (2593,1047);
\draw[mapOutline] (2584,1045) -- (2581,1050);
\draw[mapOutline] (2494,1047) -- (2489,1056) -- (2489,1059);
\draw[mapOutline] (2703,1047) -- (2698,1048) -- (2696,1052) -- (2693,1052) -- (2690,1054);
\draw[mapOutline] (2549,1048) -- (2556,1054);
\draw[mapOutline] (2596,1048) -- (2603,1051);
\draw[mapOutline] (2708,1048) -- (2704,1049);
\draw[mapOutline] (821,1049) -- (816,1053);
\draw[mapOutline] (2578,1051) -- (2574,1052);
\draw[mapOutline] (2605,1051) -- (2609,1053);
\draw[mapOutline] (1424,1052) -- (1423,1054) -- (1425,1060);
\draw[mapOutline] (1449,1052) -- (1446,1059);
\draw[mapOutline] (2572,1053) -- (2567,1055);
\draw[mapOutline] (1431,1054) -- (1428,1059);
\draw[mapOutline] (2558,1055) -- (2564,1056);
\draw[mapOutline] (447,1056) -- (446,1057) -- (447,1059);
\draw[mapOutline] (2689,1056) -- (2685,1057);
\draw[mapOutline] (2612,1057) -- (2611,1060) -- (2614,1062);
\draw[mapOutline] (2683,1057) -- (2678,1059);
\draw[mapOutline] (1444,1059) -- (1439,1060);
\draw[mapOutline] (2648,1060) -- (2651,1060) -- (2653,1062);
\draw[mapOutline] (2674,1060) -- (2670,1063);
\draw[mapOutline] (2616,1063) -- (2621,1064) -- (2623,1063);
\draw[mapOutline] (450,1064) -- (451,1067) -- (450,1073);
\draw[mapOutline] (1180,1064) -- (1179,1069);
\draw[mapOutline] (2654,1064) -- (2651,1066) -- (2652,1069) -- (2660,1066);
\draw[mapOutline] (797,1065) -- (792,1067) -- (786,1072);
\draw[mapOutline] (2668,1065) -- (2662,1066);
\draw[mapOutline] (2490,1066) -- (2485,1072) -- (2481,1073) -- (2477,1078);
\draw[mapOutline] (2639,1066) -- (2640,1067) -- (2642,1066);
\draw[mapOutline] (2637,1070) -- (2633,1071) -- (2631,1073) -- (2625,1073) -- (2624,1072);
\draw[mapOutline] (784,1074) -- (778,1078);
\draw[mapOutline] (449,1075) -- (447,1079);
\draw[mapOutline] (1181,1079) -- (1180,1083);
\draw[mapOutline] (778,1081) -- (778,1085);
\draw[mapOutline] (2475,1082) -- (2473,1086);
\draw[mapOutline] (443,1084) -- (439,1088) -- (435,1096);
\draw[mapOutline] (778,1088) -- (773,1094);
\draw[mapOutline] (2472,1089) -- (2472,1091) -- (2470,1093) -- (2466,1094);
\draw[mapOutline] (2428,1091) -- (2425,1091) -- (2423,1093);
\draw[mapOutline] (2414,1094) -- (2420,1094);
\draw[mapOutline] (2446,1094) -- (2444,1094) -- (2442,1097);
\draw[mapOutline] (2458,1094) -- (2464,1094);
\draw[mapOutline] (1181,1096) -- (1179,1101);
\draw[mapOutline] (2449,1096) -- (2450,1095) -- (2456,1096);
\draw[mapOutline] (771,1097) -- (767,1102) -- (767,1104) -- (765,1105);
\draw[mapOutline] (2426,1099) -- (2430,1100);
\draw[mapOutline] (2441,1100) -- (2440,1101) -- (2435,1101);
\draw[mapOutline] (435,1103) -- (434,1105) -- (435,1109);
\draw[mapOutline] (1178,1105) -- (1177,1110);
\draw[mapOutline] (2372,1106) -- (2369,1108);
\draw[mapOutline] (772,1108) -- (780,1111);
\draw[mapOutline] (783,1114) -- (790,1123);
\draw[mapOutline] (434,1115) -- (434,1119);
\draw[mapOutline] (1176,1117) -- (1176,1120) -- (1166,1129) -- (1167,1135);
\draw[mapOutline] (2380,1119) -- (2380,1124);
\draw[mapOutline] (432,1126) -- (432,1130);
\draw[mapOutline] (789,1129) -- (791,1137) -- (786,1138);
\draw[mapOutline] (2379,1129) -- (2378,1133);
\draw[mapOutline] (2377,1136) -- (2375,1143);
\draw[mapOutline] (431,1137) -- (431,1141);
\draw[mapOutline] (783,1140) -- (780,1142) -- (777,1142);
\draw[mapOutline] (1165,1141) -- (1165,1146);
\draw[mapOutline] (775,1145) -- (775,1147) -- (773,1149);
\draw[mapOutline] (1168,1149) -- (1170,1151) -- (1170,1153);
\draw[mapOutline] (763,1152) -- (761,1156);
\draw[mapOutline] (1478,1153) -- (1476,1161);
\draw[mapOutline] (2371,1159) -- (2374,1161);
\draw[mapOutline] (429,1161) -- (430,1165);
\draw[mapOutline] (1475,1164) -- (1470,1165);
\draw[mapOutline] (1169,1166) -- (1172,1167) -- (1173,1169);
\draw[mapOutline] (749,1168) -- (745,1169);
\draw[mapOutline] (429,1171) -- (429,1176);
\draw[mapOutline] (740,1171) -- (738,1172) -- (740,1175);
\draw[mapOutline] (2380,1171) -- (2382,1173) -- (2382,1176);
\draw[mapOutline] (1174,1179) -- (1178,1181);
\draw[mapOutline] (2384,1179) -- (2382,1182) -- (2384,1185);
\draw[mapOutline] (428,1184) -- (428,1188);
\draw[mapOutline] (1177,1184) -- (1176,1185) -- (1169,1185);
\draw[mapOutline] (2388,1185) -- (2391,1189);
\draw[mapOutline] (1167,1186) -- (1163,1190);
\draw[mapOutline] (1149,1190) -- (1145,1192) -- (1142,1192);
\draw[mapOutline] (1155,1191) -- (1160,1191);
\draw[mapOutline] (2390,1191) -- (2391,1192) -- (2389,1196);
\draw[mapOutline] (428,1194) -- (427,1198);
\draw[mapOutline] (752,1194) -- (756,1197);
\draw[mapOutline] (1141,1194) -- (1137,1196);
\draw[mapOutline] (1135,1198) -- (1130,1201);
\draw[mapOutline] (844,1200) -- (847,1200) -- (850,1202);
\draw[mapOutline] (2388,1200) -- (2387,1203) -- (2389,1205);
\draw[mapOutline] (753,1201) -- (757,1204);
\draw[mapOutline] (833,1202) -- (829,1204);
\draw[mapOutline] (1127,1203) -- (1121,1206);
\draw[mapOutline] (823,1204) -- (819,1205);
\draw[mapOutline] (853,1204) -- (854,1206) -- (859,1207);
\draw[mapOutline] (247,1205) -- (247,1213);
\draw[mapOutline] (758,1207) -- (754,1208) -- (752,1211) -- (755,1215) -- (762,1215);
\draw[mapOutline] (861,1208) -- (864,1212);
\draw[mapOutline] (1117,1210) -- (1113,1212);
\draw[mapOutline] (807,1214) -- (805,1216) -- (802,1216);
\draw[mapOutline] (907,1214) -- (912,1212) -- (916,1214);
\draw[mapOutline] (767,1215) -- (771,1215);
\draw[mapOutline] (776,1215) -- (778,1216) -- (784,1216) -- (785,1215) -- (789,1216);
\draw[mapOutline] (905,1215) -- (901,1216) -- (900,1218);
\draw[mapOutline] (794,1216) -- (798,1216);
\draw[mapOutline] (1109,1216) -- (1103,1220);
\draw[mapOutline] (425,1217) -- (425,1222);
\draw[mapOutline] (917,1217) -- (920,1221) -- (926,1222) -- (929,1224);
\draw[mapOutline] (873,1219) -- (878,1225) -- (881,1225);
\draw[mapOutline] (899,1221) -- (901,1225) -- (899,1229) -- (893,1229);
\draw[mapOutline] (1102,1222) -- (1099,1225);
\draw[mapOutline] (273,1224) -- (274,1225) -- (277,1224);
\draw[mapOutline] (932,1225) -- (937,1230);
\draw[mapOutline] (1091,1225) -- (1097,1225);
\draw[mapOutline] (883,1226) -- (885,1228) -- (891,1228);
\draw[mapOutline] (1086,1226) -- (1086,1230);
\draw[mapOutline] (259,1229) -- (253,1230) -- (247,1236);
\draw[mapOutline] (937,1232) -- (941,1233);
\draw[mapOutline] (1087,1233) -- (1087,1240);
\draw[mapOutline] (355,1236) -- (361,1237);
\draw[mapOutline] (246,1238) -- (240,1241);
\draw[mapOutline] (943,1238) -- (944,1242);
\draw[mapOutline] (242,1244) -- (245,1246);
\draw[mapOutline] (407,1244) -- (412,1247) -- (413,1246);
\draw[mapOutline] (945,1246) -- (945,1250);
\draw[mapOutline] (2422,1246) -- (2428,1251);
\draw[mapOutline] (248,1250) -- (248,1255) -- (246,1257);
\draw[mapOutline] (1080,1252) -- (1078,1255);
\draw[mapOutline] (2430,1254) -- (2436,1258);
\draw[mapOutline] (957,1257) -- (961,1258);
\draw[mapOutline] (244,1259) -- (245,1264);
\draw[mapOutline] (961,1260) -- (960,1261) -- (961,1265);
\draw[mapOutline] (1075,1260) -- (1076,1266);
\draw[mapOutline] (961,1268) -- (958,1271) -- (959,1275);
\draw[mapOutline] (245,1270) -- (245,1274) -- (243,1275);
\draw[mapOutline] (1076,1270) -- (1078,1275);
\draw[mapOutline] (2436,1273) -- (2438,1277);
\draw[mapOutline] (242,1278) -- (257,1290) -- (262,1292);
\draw[mapOutline] (2440,1278) -- (2447,1282);
\draw[mapOutline] (965,1279) -- (959,1285);
\draw[mapOutline] (2448,1287) -- (2450,1291);
\draw[mapOutline] (2475,1291) -- (2482,1292) -- (2484,1297);
\draw[mapOutline] (2475,1291) -- (2471,1294) -- (2467,1294) -- (2466,1298);
\draw[mapOutline] (959,1293) -- (956,1300);
\draw[mapOutline] (266,1296) -- (270,1299);
\draw[mapOutline] (2451,1297) -- (2451,1301);
\draw[mapOutline] (2484,1300) -- (2484,1303) -- (2481,1309);
\draw[mapOutline] (244,1301) -- (245,1302) -- (250,1302);
\draw[mapOutline] (953,1302) -- (946,1308) -- (943,1307);
\draw[mapOutline] (2465,1303) -- (2463,1305) -- (2459,1306) -- (2456,1304);
\draw[mapOutline] (254,1304) -- (260,1306);
\draw[mapOutline] (269,1304) -- (265,1305);
\draw[mapOutline] (223,1306) -- (228,1308);
\draw[mapOutline] (230,1308) -- (231,1309) -- (237,1308);
\draw[mapOutline] (940,1310) -- (936,1314) -- (937,1318);
\draw[mapOutline] (2479,1311) -- (2474,1316) -- (2472,1322);
\draw[mapOutline] (226,1320) -- (228,1324);
\draw[mapOutline] (1083,1321) -- (1091,1326) -- (1095,1325);
\draw[mapOutline] (1099,1323) -- (1098,1328);
\draw[mapOutline] (2468,1324) -- (2466,1326) -- (2467,1329);
\draw[mapOutline] (230,1331) -- (230,1335);
\draw[mapOutline] (2466,1333) -- (2464,1338) -- (2462,1339);
\draw[mapOutline] (2462,1339) -- (2462,1342) -- (2464,1346);
\draw[mapOutline] (229,1343) -- (229,1347);
\draw[mapOutline] (1101,1343) -- (1106,1346);
\draw[mapOutline] (2465,1349) -- (2468,1355);
\draw[mapOutline] (2414,1353) -- (2405,1364) -- (2403,1364);
\draw[mapOutline] (2416,1353) -- (2421,1353) -- (2422,1354);
\draw[mapOutline] (1089,1355) -- (1083,1360) -- (1082,1365);
\draw[mapOutline] (2425,1355) -- (2435,1365) -- (2438,1365);
\draw[mapOutline] (1092,1356) -- (1096,1360);
\draw[mapOutline] (2468,1359) -- (2469,1364);
\draw[mapOutline] (225,1360) -- (230,1364);
\draw[mapOutline] (2401,1364) -- (2392,1375);
\draw[mapOutline] (2443,1366) -- (2446,1371);
\draw[mapOutline] (2462,1367) -- (2466,1369);
\draw[mapOutline] (1264,1368) -- (1261,1370);
\draw[mapOutline] (1270,1368) -- (1273,1371);
\draw[mapOutline] (2472,1370) -- (2474,1373);
\draw[mapOutline] (938,1371) -- (945,1375) -- (948,1379);
\draw[mapOutline] (2453,1371) -- (2451,1370) -- (2449,1372);
\draw[mapOutline] (1255,1372) -- (1251,1374);
\draw[mapOutline] (1282,1374) -- (1275,1375);
\draw[mapOutline] (2465,1375) -- (2466,1380);
\draw[mapOutline] (1189,1377) -- (1183,1382);
\draw[mapOutline] (1246,1377) -- (1241,1378);
\draw[mapOutline] (2477,1377) -- (2479,1379) -- (2479,1381) -- (2481,1382);
\draw[mapOutline] (963,1378) -- (965,1378) -- (966,1381);
\draw[mapOutline] (1284,1378) -- (1282,1382);
\draw[mapOutline] (1190,1379) -- (1192,1382) -- (1198,1382) -- (1205,1386) -- (1210,1391) -- (1213,1392);
\draw[mapOutline] (960,1381) -- (958,1383) -- (956,1383);
\draw[mapOutline] (1080,1381) -- (1078,1385);
\draw[mapOutline] (1235,1381) -- (1229,1383) -- (1228,1382);
\draw[mapOutline] (949,1382) -- (953,1382);
\draw[mapOutline] (967,1383) -- (970,1383) -- (973,1381) -- (983,1390) -- (985,1388) -- (987,1388) -- (992,1393);
\draw[mapOutline] (1225,1383) -- (1221,1385);
\draw[mapOutline] (1180,1384) -- (1178,1388) -- (1175,1390);
\draw[mapOutline] (2466,1384) -- (2467,1386) -- (2470,1386);
\draw[mapOutline] (2483,1385) -- (2482,1386) -- (2486,1392);
\draw[mapOutline] (1282,1387) -- (1285,1391);
\draw[mapOutline] (1076,1388) -- (1073,1391) -- (1071,1391);
\draw[mapOutline] (1174,1388) -- (1169,1389);
\draw[mapOutline] (267,1389) -- (272,1389);
\draw[mapOutline] (2469,1389) -- (2469,1395);
\draw[mapOutline] (294,1392) -- (302,1393);
\draw[mapOutline] (996,1393) -- (1002,1399);
\draw[mapOutline] (1166,1393) -- (1162,1396);
\draw[mapOutline] (1285,1393) -- (1292,1395);
\draw[mapOutline] (2386,1395) -- (2376,1405);
\draw[mapOutline] (374,1397) -- (371,1401);
\draw[mapOutline] (1011,1397) -- (1007,1398) -- (1005,1400);
\draw[mapOutline] (1055,1397) -- (1050,1400);
\draw[mapOutline] (2487,1397) -- (2492,1399);
\draw[mapOutline] (317,1398) -- (311,1399);
\draw[mapOutline] (319,1399) -- (325,1403);
\draw[mapOutline] (1021,1399) -- (1025,1403) -- (1031,1402);
\draw[mapOutline] (1291,1399) -- (1289,1403);
\draw[mapOutline] (1475,1399) -- (1477,1398) -- (1482,1402);
\draw[mapOutline] (707,1400) -- (709,1399) -- (713,1401);
\draw[mapOutline] (1472,1400) -- (1468,1406) -- (1469,1407);
\draw[mapOutline] (1496,1400) -- (1493,1404) -- (1485,1403);
\draw[mapOutline] (1643,1400) -- (1649,1400);
\draw[mapOutline] (705,1401) -- (703,1404);
\draw[mapOutline] (1047,1401) -- (1043,1404);
\draw[mapOutline] (1157,1401) -- (1162,1402);
\draw[mapOutline] (1635,1401) -- (1631,1404);
\draw[mapOutline] (1653,1401) -- (1653,1404) -- (1655,1407);
\draw[mapOutline] (389,1403) -- (390,1407) -- (392,1409);
\draw[mapOutline] (715,1403) -- (719,1407) -- (724,1409);
\draw[mapOutline] (1034,1403) -- (1038,1404) -- (1040,1403);
\draw[mapOutline] (1500,1403) -- (1502,1406);
\draw[mapOutline] (2469,1403) -- (2469,1405) -- (2466,1408) -- (2463,1407);
\draw[mapOutline] (2475,1403) -- (2478,1401) -- (2481,1401) -- (2483,1404) -- (2486,1404) -- (2491,1408);
\draw[mapOutline] (1613,1408) -- (1608,1411);
\draw[mapOutline] (2461,1408) -- (2460,1413) -- (2463,1416) -- (2461,1420) -- (2457,1423) -- (2458,1426);
\draw[mapOutline] (1289,1409) -- (1291,1412);
\draw[mapOutline] (2374,1409) -- (2368,1414);
\draw[mapOutline] (340,1410) -- (344,1411);
\draw[mapOutline] (699,1410) -- (695,1413);
\draw[mapOutline] (726,1410) -- (730,1412);
\draw[mapOutline] (1462,1410) -- (1457,1414);
\draw[mapOutline] (361,1411) -- (357,1416) -- (349,1414);
\draw[mapOutline] (1605,1413) -- (1600,1415);
\draw[mapOutline] (1160,1414) -- (1165,1414);
\draw[mapOutline] (1657,1414) -- (1656,1415) -- (1657,1420);
\draw[mapOutline] (737,1415) -- (743,1416);
\draw[mapOutline] (1597,1415) -- (1594,1418) -- (1591,1424);
\draw[mapOutline] (746,1417) -- (751,1419);
\draw[mapOutline] (1512,1417) -- (1511,1418) -- (1512,1421);
\draw[mapOutline] (392,1421) -- (393,1426);
\draw[mapOutline] (1659,1421) -- (1663,1424);
\draw[mapOutline] (2364,1423) -- (2360,1429) -- (2362,1432);
\draw[mapOutline] (1074,1424) -- (1071,1427) -- (1072,1432);
\draw[mapOutline] (1512,1424) -- (1508,1427) -- (1506,1427);
\draw[mapOutline] (751,1425) -- (751,1429);
\draw[mapOutline] (1665,1425) -- (1671,1432);
\draw[mapOutline] (1590,1426) -- (1590,1430);
\draw[mapOutline] (691,1428) -- (687,1433);
\draw[mapOutline] (1095,1428) -- (1091,1429);
\draw[mapOutline] (1158,1428) -- (1157,1431) -- (1154,1433);
\draw[mapOutline] (1450,1429) -- (1446,1433);
\draw[mapOutline] (397,1430) -- (400,1435);
\draw[mapOutline] (1097,1431) -- (1103,1436);
\draw[mapOutline] (2393,1432) -- (2394,1431) -- (2402,1436);
\draw[mapOutline] (687,1433) -- (689,1435) -- (693,1435);
\draw[mapOutline] (1669,1433) -- (1663,1438) -- (1664,1442);
\draw[mapOutline] (1132,1434) -- (1131,1435) -- (1129,1434) -- (1126,1436);
\draw[mapOutline] (1591,1434) -- (1586,1439);
\draw[mapOutline] (751,1435) -- (750,1440);
\draw[mapOutline] (1073,1435) -- (1076,1437);
\draw[mapOutline] (1152,1435) -- (1148,1438);
\draw[mapOutline] (698,1437) -- (702,1437);
\draw[mapOutline] (1146,1437) -- (1143,1436) -- (1142,1439);
\draw[mapOutline] (2401,1438) -- (2399,1442) -- (2401,1444);
\draw[mapOutline] (403,1439) -- (403,1444) -- (406,1447);
\draw[mapOutline] (1107,1439) -- (1111,1438) -- (1116,1440) -- (1121,1439);
\draw[mapOutline] (2462,1440) -- (2467,1439) -- (2472,1443);
\draw[mapOutline] (1075,1441) -- (1075,1445);
\draw[mapOutline] (2632,1442) -- (2636,1442);
\draw[mapOutline] (1437,1443) -- (1434,1448) -- (1437,1452);
\draw[mapOutline] (2661,1443) -- (2657,1447);
\draw[mapOutline] (2646,1444) -- (2647,1443) -- (2648,1444) -- (2647,1446);
\draw[mapOutline] (2666,1444) -- (2670,1444);
\draw[mapOutline] (1497,1445) -- (1501,1447);
\draw[mapOutline] (2474,1445) -- (2473,1446) -- (2475,1448) -- (2474,1453);
\draw[mapOutline] (2631,1445) -- (2635,1446);
\draw[mapOutline] (2386,1446) -- (2386,1451);
\draw[mapOutline] (408,1447) -- (413,1448);
\draw[mapOutline] (1587,1447) -- (1585,1449) -- (1583,1449);
\draw[mapOutline] (2637,1447) -- (2640,1451) -- (2638,1453);
\draw[mapOutline] (750,1448) -- (753,1454) -- (758,1457);
\draw[mapOutline] (703,1449) -- (695,1454);
\draw[mapOutline] (1664,1449) -- (1664,1451) -- (1663,1452) -- (1662,1451);
\draw[mapOutline] (2367,1449) -- (2368,1450) -- (2373,1450) -- (2374,1452);
\draw[mapOutline] (2400,1449) -- (2401,1453) -- (2397,1460) -- (2400,1463);
\draw[mapOutline] (2647,1449) -- (2653,1450);
\draw[mapOutline] (131,1450) -- (135,1450);
\draw[mapOutline] (1080,1450) -- (1083,1456) -- (1079,1462) -- (1081,1467);
\draw[mapOutline] (1578,1450) -- (1574,1452);
\draw[mapOutline] (413,1451) -- (415,1458);
\draw[mapOutline] (1506,1451) -- (1507,1455) -- (1503,1457) -- (1503,1459) -- (1505,1460);
\draw[mapOutline] (2386,1451) -- (2377,1452);
\draw[mapOutline] (2386,1451) -- (2388,1453) -- (2390,1453);
\draw[mapOutline] (1437,1452) -- (1436,1456);
\draw[mapOutline] (1569,1452) -- (1562,1457);
\draw[mapOutline] (1659,1452) -- (1656,1453) -- (1654,1456);
\draw[mapOutline] (2632,1453) -- (2633,1454) -- (2636,1453);
\draw[mapOutline] (2680,1453) -- (2679,1454) -- (2683,1459);
\draw[mapOutline] (2691,1454) -- (2697,1460) -- (2694,1464) -- (2700,1470) -- (2700,1472);
\draw[mapOutline] (2689,1455) -- (2686,1458);
\draw[mapOutline] (694,1456) -- (690,1459) -- (690,1461);
\draw[mapOutline] (1295,1457) -- (1292,1460);
\draw[mapOutline] (1425,1460) -- (1422,1465);
\draw[mapOutline] (2483,1461) -- (2486,1464);
\draw[mapOutline] (763,1462) -- (765,1464) -- (767,1464);
\draw[mapOutline] (1506,1462) -- (1508,1464) -- (1510,1463);
\draw[mapOutline] (1510,1463) -- (1512,1464) -- (1512,1468);
\draw[mapOutline] (2400,1463) -- (2402,1462) -- (2403,1463) -- (2402,1465);
\draw[mapOutline] (690,1464) -- (692,1467);
\draw[mapOutline] (1289,1464) -- (1288,1469) -- (1291,1473) -- (1291,1476);
\draw[mapOutline] (2389,1464) -- (2389,1466) -- (2387,1468) -- (2388,1471);
\draw[mapOutline] (1602,1465) -- (1598,1466);
\draw[mapOutline] (1631,1465) -- (1626,1470) -- (1622,1472) -- (1621,1476);
\draw[mapOutline] (2402,1465) -- (2404,1470) -- (2401,1471);
\draw[mapOutline] (1607,1466) -- (1609,1469);
\draw[mapOutline] (2627,1466) -- (2629,1469);
\draw[mapOutline] (772,1467) -- (777,1469);
\draw[mapOutline] (1593,1467) -- (1588,1469);
\draw[mapOutline] (2486,1467) -- (2491,1472) -- (2493,1478);
\draw[mapOutline] (691,1469) -- (686,1476) -- (686,1480) -- (681,1484) -- (681,1486);
\draw[mapOutline] (1555,1469) -- (1550,1471);
\draw[mapOutline] (779,1470) -- (781,1473) -- (785,1475);
\draw[mapOutline] (1512,1470) -- (1508,1477);
\draw[mapOutline] (1420,1471) -- (1418,1472) -- (1418,1475);
\draw[mapOutline] (1612,1471) -- (1615,1473) -- (1615,1475);
\draw[mapOutline] (1082,1472) -- (1087,1476);
\draw[mapOutline] (1547,1472) -- (1542,1474) -- (1536,1473) -- (1534,1476);
\draw[mapOutline] (1578,1473) -- (1574,1475) -- (1571,1474);
\draw[mapOutline] (2435,1473) -- (2433,1475) -- (2433,1482) -- (2432,1483) -- (2429,1482);
\draw[mapOutline] (2437,1473) -- (2440,1475) -- (2443,1475) -- (2447,1473) -- (2448,1474);
\draw[mapOutline] (2696,1474) -- (2695,1478) -- (2699,1483);
\draw[mapOutline] (1418,1475) -- (1411,1479) -- (1409,1479);
\draw[mapOutline] (1092,1476) -- (1095,1478);
\draw[mapOutline] (1626,1476) -- (1630,1476);
\draw[mapOutline] (789,1478) -- (792,1481) -- (794,1481);
\draw[mapOutline] (1402,1478) -- (1407,1479);
\draw[mapOutline] (1392,1479) -- (1389,1481);
\draw[mapOutline] (1392,1479) -- (1392,1481) -- (1394,1483) -- (1396,1483) -- (1400,1479);
\draw[mapOutline] (1532,1479) -- (1527,1484);
\draw[mapOutline] (1631,1479) -- (1629,1482);
\draw[mapOutline] (2422,1480) -- (2426,1480);
\draw[mapOutline] (1505,1481) -- (1505,1483) -- (1502,1486) -- (1496,1488);
\draw[mapOutline] (2420,1481) -- (2416,1483) -- (2417,1487) -- (2413,1489);
\draw[mapOutline] (2497,1481) -- (2500,1483);
\draw[mapOutline] (796,1483) -- (801,1487) -- (803,1487);
\draw[mapOutline] (1328,1484) -- (1324,1485);
\draw[mapOutline] (1331,1484) -- (1338,1484);
\draw[mapOutline] (1343,1485) -- (1358,1488);
\draw[mapOutline] (1295,1486) -- (1300,1482) -- (1303,1484) -- (1303,1488) -- (1306,1490);
\draw[mapOutline] (1666,1486) -- (1668,1489);
\draw[mapOutline] (1628,1487) -- (1628,1490) -- (1631,1493) -- (1633,1491);
\draw[mapOutline] (2396,1487) -- (2399,1490) -- (2400,1493) -- (2402,1494);
\draw[mapOutline] (2702,1487) -- (2706,1488);
\draw[mapOutline] (1368,1488) -- (1363,1489);
\draw[mapOutline] (843,1489) -- (848,1489);
\draw[mapOutline] (1313,1489) -- (1308,1492);
\draw[mapOutline] (2413,1489) -- (2411,1490) -- (2408,1495) -- (2406,1493) -- (2402,1494);
\draw[mapOutline] (2502,1489) -- (2505,1492);
\draw[mapOutline] (2708,1489) -- (2711,1493);
\draw[mapOutline] (663,1490) -- (657,1495) -- (651,1496) -- (641,1505);
\draw[mapOutline] (807,1490) -- (811,1493);
\draw[mapOutline] (1673,1491) -- (1677,1491);
\draw[mapOutline] (2631,1491) -- (2632,1495);
\draw[mapOutline] (2720,1491) -- (2717,1491) -- (2715,1493);
\draw[mapOutline] (2723,1491) -- (2732,1494);
\draw[mapOutline] (1680,1492) -- (1684,1497);
\draw[mapOutline] (2181,1492) -- (2181,1496);
\draw[mapOutline] (839,1493) -- (837,1495) -- (835,1495);
\draw[mapOutline] (814,1494) -- (820,1496);
\draw[mapOutline] (2402,1494) -- (2408,1501);
\draw[mapOutline] (830,1496) -- (826,1497);
\draw[mapOutline] (1125,1496) -- (1131,1501);
\draw[mapOutline] (412,1497) -- (410,1501) -- (408,1502);
\draw[mapOutline] (2513,1498) -- (2517,1499) -- (2526,1509) -- (2529,1509);
\draw[mapOutline] (2738,1498) -- (2740,1503);
\draw[mapOutline] (871,1499) -- (878,1502);
\draw[mapOutline] (2633,1500) -- (2633,1504) -- (2641,1512);
\draw[mapOutline] (2410,1501) -- (2415,1505);
\draw[mapOutline] (1493,1502) -- (1495,1505) -- (1490,1508);
\draw[mapOutline] (1597,1502) -- (1600,1505);
\draw[mapOutline] (1683,1502) -- (1683,1506);
\draw[mapOutline] (1131,1503) -- (1137,1507);
\draw[mapOutline] (2730,1504) -- (2726,1506) -- (2725,1510) -- (2723,1510) -- (2721,1507);
\draw[mapOutline] (2327,1505) -- (2335,1506) -- (2337,1509);
\draw[mapOutline] (1509,1506) -- (1503,1507);
\draw[mapOutline] (2531,1506) -- (2534,1506) -- (2535,1508);
\draw[mapOutline] (638,1507) -- (635,1509);
\draw[mapOutline] (1139,1508) -- (1149,1512);
\draw[mapOutline] (1488,1509) -- (1485,1514) -- (1485,1517);
\draw[mapOutline] (2340,1509) -- (2350,1512);
\draw[mapOutline] (2676,1509) -- (2669,1512);
\draw[mapOutline] (2535,1510) -- (2535,1514);
\draw[mapOutline] (2428,1511) -- (2430,1512) -- (2434,1511);
\draw[mapOutline] (2677,1511) -- (2676,1518);
\draw[mapOutline] (2709,1511) -- (2709,1514) -- (2713,1517) -- (2711,1522);
\draw[mapOutline] (620,1512) -- (616,1514);
\draw[mapOutline] (879,1512) -- (880,1518);
\draw[mapOutline] (1683,1513) -- (1680,1516) -- (1674,1518);
\draw[mapOutline] (2427,1513) -- (2425,1517);
\draw[mapOutline] (2667,1513) -- (2662,1513) -- (2661,1514);
\draw[mapOutline] (2195,1514) -- (2194,1515) -- (2192,1514) -- (2189,1515) -- (2191,1519);
\draw[mapOutline] (2736,1514) -- (2734,1516) -- (2732,1516);
\draw[mapOutline] (403,1515) -- (403,1519) -- (401,1521);
\draw[mapOutline] (1150,1515) -- (1152,1519);
\draw[mapOutline] (2133,1517) -- (2129,1514) -- (2126,1515) -- (2124,1517) -- (2124,1522) -- (2117,1526);
\draw[mapOutline] (2136,1517) -- (2137,1516) -- (2139,1517) -- (2149,1517);
\draw[mapOutline] (2174,1517) -- (2174,1526);
\draw[mapOutline] (1602,1520) -- (1602,1524);
\draw[mapOutline] (2656,1520) -- (2662,1523);
\draw[mapOutline] (2676,1521) -- (2676,1523) -- (2686,1534);
\draw[mapOutline] (879,1522) -- (881,1525);
\draw[mapOutline] (1153,1522) -- (1152,1523) -- (1155,1529);
\draw[mapOutline] (2420,1522) -- (2425,1523) -- (2429,1527) -- (2435,1529) -- (2439,1533);
\draw[mapOutline] (1479,1523) -- (1477,1529);
\draw[mapOutline] (2148,1523) -- (2147,1528) -- (2149,1529);
\draw[mapOutline] (2662,1523) -- (2666,1527);
\draw[mapOutline] (2353,1524) -- (2357,1525);
\draw[mapOutline] (2360,1525) -- (2367,1527);
\draw[mapOutline] (2708,1525) -- (2704,1526);
\draw[mapOutline] (2242,1526) -- (2237,1527) -- (2232,1533) -- (2227,1530);
\draw[mapOutline] (2244,1526) -- (2248,1527) -- (2252,1531) -- (2256,1531);
\draw[mapOutline] (1455,1527) -- (1448,1530) -- (1442,1536) -- (1441,1541);
\draw[mapOutline] (1458,1527) -- (1466,1529);
\draw[mapOutline] (2295,1527) -- (2293,1535);
\draw[mapOutline] (2522,1527) -- (2518,1534);
\draw[mapOutline] (1603,1528) -- (1602,1529) -- (1605,1533);
\draw[mapOutline] (2190,1528) -- (2197,1530);
\draw[mapOutline] (1469,1529) -- (1475,1529);
\draw[mapOutline] (2283,1529) -- (2279,1530);
\draw[mapOutline] (395,1530) -- (393,1535);
\draw[mapOutline] (883,1530) -- (886,1533);
\draw[mapOutline] (2112,1530) -- (2117,1531);
\draw[mapOutline] (2667,1530) -- (2667,1535) -- (2672,1539);
\draw[mapOutline] (2702,1530) -- (2697,1536);
\draw[mapOutline] (2261,1531) -- (2261,1533) -- (2259,1535) -- (2256,1535);
\draw[mapOutline] (2277,1531) -- (2273,1533) -- (2265,1541);
\draw[mapOutline] (2150,1533) -- (2149,1535) -- (2147,1535);
\draw[mapOutline] (2370,1533) -- (2372,1536);
\draw[mapOutline] (1156,1534) -- (1157,1539);
\draw[mapOutline] (2515,1534) -- (2513,1534) -- (2509,1531) -- (2505,1535) -- (2499,1536) -- (2496,1534) -- (2493,1535) -- (2492,1537);
\draw[mapOutline] (2652,1534) -- (2659,1539);
\draw[mapOutline] (887,1536) -- (894,1542);
\draw[mapOutline] (2186,1536) -- (2186,1538) -- (2188,1539);
\draw[mapOutline] (2443,1536) -- (2448,1539);
\draw[mapOutline] (2694,1536) -- (2692,1538) -- (2688,1537);
\draw[mapOutline] (1606,1537) -- (1603,1541);
\draw[mapOutline] (2157,1537) -- (2151,1542) -- (2147,1542) -- (2145,1538);
\draw[mapOutline] (2254,1539) -- (2253,1545) -- (2252,1546) -- (2249,1545) -- (2248,1549);
\draw[mapOutline] (1663,1540) -- (1663,1542) -- (1661,1544);
\draw[mapOutline] (2493,1540) -- (2498,1541) -- (2499,1546);
\draw[mapOutline] (2378,1541) -- (2381,1543) -- (2386,1543) -- (2388,1541) -- (2392,1541) -- (2395,1544) -- (2395,1547);
\draw[mapOutline] (2449,1542) -- (2448,1544) -- (2449,1546);
\draw[mapOutline] (2674,1542) -- (2669,1544);
\draw[mapOutline] (392,1543) -- (393,1547);
\draw[mapOutline] (2187,1544) -- (2188,1547) -- (2194,1553) -- (2197,1554);
\draw[mapOutline] (2187,1544) -- (2186,1543) -- (2187,1542) -- (2188,1543) -- cycle;
\draw[mapOutline] (1156,1545) -- (1158,1550);
\draw[mapOutline] (2264,1545) -- (2262,1547) -- (2259,1547);
\draw[mapOutline] (2103,1546) -- (2104,1551) -- (2102,1552);
\draw[mapOutline] (1595,1547) -- (1594,1549) -- (1595,1550);
\draw[mapOutline] (1440,1548) -- (1442,1551);
\draw[mapOutline] (2168,1549) -- (2165,1557) -- (2165,1565);
\draw[mapOutline] (2207,1549) -- (2212,1552);
\draw[mapOutline] (2257,1549) -- (2250,1552);
\draw[mapOutline] (2396,1549) -- (2399,1553);
\draw[mapOutline] (2450,1549) -- (2449,1550) -- (2450,1551) -- (2456,1552);
\draw[mapOutline] (2456,1549) -- (2459,1548) -- (2465,1555) -- (2463,1559) -- (2465,1561);
\draw[mapOutline] (1160,1550) -- (1161,1554);
\draw[mapOutline] (1660,1550) -- (1660,1554);
\draw[mapOutline] (2204,1550) -- (2202,1553) -- (2202,1559);
\draw[mapOutline] (877,1551) -- (871,1557) -- (870,1565) -- (876,1569) -- (877,1568);
\draw[mapOutline] (893,1551) -- (900,1554) -- (902,1553);
\draw[mapOutline] (2135,1551) -- (2131,1552);
\draw[mapOutline] (394,1552) -- (399,1558);
\draw[mapOutline] (909,1552) -- (907,1552) -- (905,1554);
\draw[mapOutline] (2097,1553) -- (2098,1555) -- (2100,1555);
\draw[mapOutline] (2129,1553) -- (2126,1555) -- (2120,1556);
\draw[mapOutline] (2145,1553) -- (2145,1557) -- (2147,1558);
\draw[mapOutline] (1445,1554) -- (1444,1555) -- (1444,1561);
\draw[mapOutline] (2102,1555) -- (2105,1556) -- (2108,1559) -- (2112,1559) -- (2117,1557);
\draw[mapOutline] (2505,1555) -- (2510,1559) -- (2509,1562) -- (2510,1565);
\draw[mapOutline] (1161,1556) -- (1154,1561);
\draw[mapOutline] (2211,1557) -- (2210,1559) -- (2213,1562);
\draw[mapOutline] (2198,1558) -- (2199,1560) -- (2201,1560);
\draw[mapOutline] (2225,1558) -- (2225,1561) -- (2219,1566) -- (2215,1564);
\draw[mapOutline] (399,1560) -- (400,1564);
\draw[mapOutline] (2147,1560) -- (2147,1563) -- (2149,1567);
\draw[mapOutline] (1157,1563) -- (1157,1567);
\draw[mapOutline] (1589,1563) -- (1583,1567) -- (1578,1575) -- (1585,1583) -- (1585,1587) -- (1582,1591) -- (1579,1592);
\draw[mapOutline] (2468,1564) -- (2473,1565);
\draw[mapOutline] (459,1565) -- (455,1566);
\draw[mapOutline] (464,1566) -- (470,1566);
\draw[mapOutline] (399,1567) -- (404,1575);
\draw[mapOutline] (544,1567) -- (540,1570);
\draw[mapOutline] (2510,1567) -- (2512,1569) -- (2512,1571);
\draw[mapOutline] (485,1568) -- (490,1569);
\draw[mapOutline] (1443,1568) -- (1442,1572) -- (1435,1579) -- (1434,1585);
\draw[mapOutline] (2475,1568) -- (2482,1577) -- (2486,1579);
\draw[mapOutline] (496,1570) -- (501,1571);
\draw[mapOutline] (879,1570) -- (885,1576) -- (886,1584);
\draw[mapOutline] (1665,1570) -- (1664,1574);
\draw[mapOutline] (507,1572) -- (511,1572);
\draw[mapOutline] (518,1572) -- (523,1573);
\draw[mapOutline] (534,1572) -- (529,1573);
\draw[mapOutline] (1546,1572) -- (1542,1573);
\draw[mapOutline] (1540,1573) -- (1535,1578);
\draw[mapOutline] (1548,1573) -- (1549,1574) -- (1547,1578);
\draw[mapOutline] (2401,1573) -- (2407,1577);
\draw[mapOutline] (2148,1574) -- (2145,1575) -- (2142,1578) -- (2142,1580);
\draw[mapOutline] (2511,1575) -- (2514,1577) -- (2514,1579);
\draw[mapOutline] (1641,1578) -- (1635,1576) -- (1632,1579) -- (1634,1581);
\draw[mapOutline] (1652,1578) -- (1647,1579);
\draw[mapOutline] (1657,1578) -- (1661,1578);
\draw[mapOutline] (2509,1578) -- (2503,1580);
\draw[mapOutline] (2408,1579) -- (2417,1588) -- (2424,1583) -- (2428,1584);
\draw[mapOutline] (2482,1581) -- (2477,1582) -- (2470,1588);
\draw[mapOutline] (413,1582) -- (421,1584);
\draw[mapOutline] (1636,1583) -- (1632,1585);
\draw[mapOutline] (409,1584) -- (408,1588) -- (406,1589);
\draw[mapOutline] (426,1584) -- (431,1584);
\draw[mapOutline] (2434,1584) -- (2440,1584) -- (2441,1585);
\draw[mapOutline] (2457,1584) -- (2454,1587) -- (2452,1587);
\draw[mapOutline] (2466,1584) -- (2468,1587);
\draw[mapOutline] (874,1585) -- (871,1586) -- (869,1591);
\draw[mapOutline] (2448,1585) -- (2443,1587);
\draw[mapOutline] (879,1586) -- (883,1586);
\draw[mapOutline] (2429,1586) -- (2433,1586);
\draw[mapOutline] (1432,1589) -- (1431,1592) -- (1432,1593);
\draw[mapOutline] (2502,1589) -- (2504,1594);
\draw[mapOutline] (2136,1590) -- (2133,1593);
\draw[mapOutline] (405,1592) -- (402,1594);
\draw[mapOutline] (1513,1592) -- (1512,1595) -- (1510,1596);
\draw[mapOutline] (868,1594) -- (867,1598);
\draw[mapOutline] (1566,1594) -- (1561,1595);
\draw[mapOutline] (2081,1594) -- (2077,1598);
\draw[mapOutline] (1558,1595) -- (1552,1596);
\draw[mapOutline] (380,1596) -- (384,1595) -- (388,1596);
\draw[mapOutline] (399,1596) -- (392,1599) -- (391,1598);
\draw[mapOutline] (1431,1596) -- (1434,1599) -- (1437,1600);
\draw[mapOutline] (1544,1596) -- (1545,1598) -- (1549,1597);
\draw[mapOutline] (378,1597) -- (375,1599);
\draw[mapOutline] (1507,1598) -- (1505,1603) -- (1496,1604) -- (1494,1608);
\draw[mapOutline] (2508,1601) -- (2509,1603) -- (2517,1601);
\draw[mapOutline] (1439,1602) -- (1443,1604) -- (1446,1604);
\draw[mapOutline] (2155,1602) -- (2154,1604) -- (2154,1613);
\draw[mapOutline] (2085,1603) -- (2086,1610);
\draw[mapOutline] (371,1604) -- (365,1618);
\draw[mapOutline] (868,1604) -- (868,1608);
\draw[mapOutline] (2516,1604) -- (2514,1610) -- (2517,1612) -- (2521,1609);
\draw[mapOutline] (1157,1606) -- (1158,1610);
\draw[mapOutline] (1450,1607) -- (1453,1609);
\draw[mapOutline] (2129,1610) -- (2126,1616) -- (2131,1622);
\draw[mapOutline] (2522,1612) -- (2521,1616);
\draw[mapOutline] (2534,1612) -- (2535,1613) -- (2541,1612) -- (2543,1614);
\draw[mapOutline] (1158,1613) -- (1150,1632);
\draw[mapOutline] (2085,1614) -- (2080,1615);
\draw[mapOutline] (2531,1614) -- (2525,1622) -- (2522,1620);
\draw[mapOutline] (2543,1614) -- (2540,1617) -- (2542,1621) -- (2542,1624) -- (2539,1627) -- (2542,1631) -- (2542,1634) -- (2539,1639);
\draw[mapOutline] (868,1615) -- (867,1616) -- (867,1621) -- (864,1624) -- (860,1625);
\draw[mapOutline] (2152,1620) -- (2152,1622) -- (2150,1624) -- (2150,1630);
\draw[mapOutline] (1508,1628) -- (1510,1629) -- (1510,1631);
\draw[mapOutline] (1150,1632) -- (1147,1636);
\draw[mapOutline] (1606,1633) -- (1607,1634) -- (1606,1638);
\draw[mapOutline] (2536,1640) -- (2531,1639) -- (2522,1642);
\draw[mapOutline] (1151,1641) -- (1153,1647) -- (1159,1654);
\draw[mapOutline] (1612,1644) -- (1619,1650) -- (1620,1654) -- (1623,1653);
\draw[mapOutline] (855,1645) -- (854,1647) -- (858,1649) -- (859,1651);
\draw[mapOutline] (1495,1648) -- (1496,1651) -- (1498,1653) -- (1500,1653);
\draw[mapOutline] (1515,1652) -- (1512,1655);
\draw[mapOutline] (1449,1653) -- (1455,1655) -- (1460,1659);
\draw[mapOutline] (1502,1653) -- (1509,1654);
\draw[mapOutline] (860,1654) -- (865,1660) -- (872,1662);
\draw[mapOutline] (1625,1655) -- (1629,1658);
\draw[mapOutline] (2111,1657) -- (2116,1654) -- (2121,1655) -- (2123,1659);
\draw[mapOutline] (1171,1659) -- (1174,1662) -- (1174,1664);
\draw[mapOutline] (1433,1659) -- (1437,1662);
\draw[mapOutline] (1631,1662) -- (1634,1662) -- (1637,1665);
\draw[mapOutline] (905,1664) -- (901,1665) -- (895,1670) -- (892,1670) -- (884,1664) -- (880,1664) -- (875,1666);
\draw[mapOutline] (910,1665) -- (913,1667);
\draw[mapOutline] (1412,1666) -- (1409,1668);
\draw[mapOutline] (1440,1666) -- (1443,1672) -- (1446,1673) -- (1449,1676) -- (1448,1680);
\draw[mapOutline] (1173,1667) -- (1171,1671) -- (1169,1672);
\draw[mapOutline] (1639,1668) -- (1646,1670);
\draw[mapOutline] (2138,1670) -- (2138,1683);
\draw[mapOutline] (1405,1671) -- (1401,1674);
\draw[mapOutline] (2506,1671) -- (2507,1675) -- (2509,1676);
\draw[mapOutline] (920,1675) -- (918,1678) -- (915,1676) -- (912,1678) -- (909,1678) -- (901,1686);
\draw[mapOutline] (1398,1676) -- (1391,1683) -- (1391,1686) -- (1395,1690);
\draw[mapOutline] (1646,1676) -- (1647,1680) -- (1643,1687);
\draw[mapOutline] (1158,1678) -- (1156,1680) -- (1156,1682);
\draw[mapOutline] (2509,1679) -- (2509,1682) -- (2511,1683);
\draw[mapOutline] (2521,1683) -- (2515,1691) -- (2512,1688);
\draw[mapOutline] (2528,1683) -- (2528,1694) -- (2529,1696) -- (2527,1697);
\draw[mapOutline] (2077,1684) -- (2076,1689);
\draw[mapOutline] (1155,1685) -- (1153,1688) -- (1146,1691);
\draw[mapOutline] (2137,1686) -- (2133,1698) -- (2132,1710);
\draw[mapOutline] (901,1688) -- (904,1688) -- (907,1690);
\draw[mapOutline] (1144,1690) -- (1140,1690) -- (1136,1692);
\draw[mapOutline] (1395,1690) -- (1391,1695);
\draw[mapOutline] (1642,1690) -- (1637,1695) -- (1634,1695);
\draw[mapOutline] (1134,1693) -- (1132,1696) -- (1122,1703);
\draw[mapOutline] (905,1694) -- (904,1695) -- (905,1699);
\draw[mapOutline] (2527,1700) -- (2524,1704);
\draw[mapOutline] (1389,1701) -- (1390,1706) -- (1385,1712) -- (1386,1717);
\draw[mapOutline] (1590,1702) -- (1590,1704) -- (1592,1706) -- (1594,1705);
\draw[mapOutline] (2068,1703) -- (2063,1709) -- (2065,1711) -- (2065,1715);
\draw[mapOutline] (1052,1704) -- (1045,1707);
\draw[mapOutline] (1576,1704) -- (1568,1713) -- (1568,1716) -- (1564,1721);
\draw[mapOutline] (1632,1706) -- (1627,1710);
\draw[mapOutline] (1614,1707) -- (1609,1711);
\draw[mapOutline] (2524,1707) -- (2522,1709) -- (2520,1715);
\draw[mapOutline] (1119,1709) -- (1119,1713);
\draw[mapOutline] (1482,1709) -- (1486,1713) -- (1495,1714);
\draw[mapOutline] (1625,1709) -- (1619,1713) -- (1616,1710);
\draw[mapOutline] (892,1710) -- (897,1711);
\draw[mapOutline] (1045,1710) -- (1040,1713);
\draw[mapOutline] (1510,1710) -- (1504,1714) -- (1497,1715);
\draw[mapOutline] (905,1711) -- (903,1713) -- (899,1712);
\draw[mapOutline] (887,1712) -- (884,1714) -- (879,1721);
\draw[mapOutline] (1471,1712) -- (1472,1711) -- (1478,1712);
\draw[mapOutline] (1057,1713) -- (1062,1715);
\draw[mapOutline] (1081,1713) -- (1085,1712) -- (1089,1716) -- (1092,1715);
\draw[mapOutline] (2131,1713) -- (2130,1714) -- (2130,1719);
\draw[mapOutline] (1038,1714) -- (1032,1718);
\draw[mapOutline] (1077,1714) -- (1071,1720) -- (1066,1720);
\draw[mapOutline] (1468,1714) -- (1464,1715);
\draw[mapOutline] (1596,1714) -- (1598,1715) -- (1602,1714) -- (1606,1718) -- (1609,1716);
\draw[mapOutline] (1094,1715) -- (1101,1718);
\draw[mapOutline] (1104,1719) -- (1103,1724) -- (1107,1728);
\draw[mapOutline] (1387,1720) -- (1392,1722);
\draw[mapOutline] (1522,1723) -- (1524,1726);
\draw[mapOutline] (2516,1723) -- (2513,1724) -- (2509,1723) -- (2503,1729);
\draw[mapOutline] (873,1724) -- (864,1726) -- (859,1729);
\draw[mapOutline] (1563,1724) -- (1558,1728);
\draw[mapOutline] (1296,1726) -- (1303,1730) -- (1311,1730) -- (1316,1734) -- (1321,1730) -- (1325,1730) -- (1329,1733) -- (1332,1733) -- (1336,1729) -- (1339,1730);
\draw[mapOutline] (1392,1726) -- (1389,1730);
\draw[mapOutline] (1027,1727) -- (1023,1728) -- (1020,1731) -- (1018,1737);
\draw[mapOutline] (1114,1727) -- (1109,1728);
\draw[mapOutline] (1293,1728) -- (1291,1733) -- (1289,1734);
\draw[mapOutline] (2073,1728) -- (2077,1735) -- (2082,1734);
\draw[mapOutline] (2245,1728) -- (2250,1728);
\draw[mapOutline] (1555,1730) -- (1553,1730) -- (1551,1727) -- (1546,1729) -- (1542,1729) -- (1538,1726) -- (1531,1729) -- (1526,1735) -- (1531,1744);
\draw[mapOutline] (2240,1730) -- (2235,1731);
\draw[mapOutline] (853,1731) -- (846,1735);
\draw[mapOutline] (1701,1731) -- (1698,1734) -- (1695,1743);
\draw[mapOutline] (1707,1731) -- (1709,1733) -- (1709,1735) -- (1707,1737) -- (1709,1740);
\draw[mapOutline] (2501,1731) -- (2498,1732) -- (2494,1736);
\draw[mapOutline] (1521,1733) -- (1517,1737);
\draw[mapOutline] (2084,1734) -- (2084,1739) -- (2085,1740);
\draw[mapOutline] (1384,1735) -- (1382,1739);
\draw[mapOutline] (1287,1736) -- (1284,1738);
\draw[mapOutline] (2491,1737) -- (2487,1738);
\draw[mapOutline] (1017,1740) -- (1015,1744);
\draw[mapOutline] (1271,1740) -- (1269,1742) -- (1269,1745);
\draw[mapOutline] (1274,1740) -- (1277,1745) -- (1280,1743);
\draw[mapOutline] (1346,1740) -- (1343,1742) -- (1341,1741);
\draw[mapOutline] (2481,1740) -- (2476,1742) -- (2474,1741);
\draw[mapOutline] (1379,1742) -- (1374,1747);
\draw[mapOutline] (1351,1743) -- (1353,1746);
\draw[mapOutline] (2471,1743) -- (2467,1744);
\draw[mapOutline] (2086,1744) -- (2089,1750);
\draw[mapOutline] (849,1745) -- (852,1749);
\draw[mapOutline] (1528,1745) -- (1521,1750);
\draw[mapOutline] (1367,1746) -- (1372,1749);
\draw[mapOutline] (1695,1746) -- (1692,1749) -- (1691,1754) -- (1689,1756) -- (1685,1756) -- (1680,1762);
\draw[mapOutline] (1365,1747) -- (1358,1750);
\draw[mapOutline] (1720,1748) -- (1719,1752) -- (1721,1755);
\draw[mapOutline] (2468,1748) -- (2477,1754) -- (2481,1755);
\draw[mapOutline] (1010,1750) -- (1007,1754);
\draw[mapOutline] (1505,1750) -- (1504,1752) -- (1504,1757) -- (1505,1758) -- (1507,1757) -- (1511,1758) -- (1517,1751) -- (1521,1750);
\draw[mapOutline] (1521,1750) -- (1524,1754) -- (1527,1752) -- (1530,1752) -- (1535,1756);
\draw[mapOutline] (2106,1750) -- (2108,1748) -- (2112,1748) -- (2113,1750);
\draw[mapOutline] (926,1754) -- (922,1756);
\draw[mapOutline] (944,1754) -- (949,1755);
\draw[mapOutline] (928,1755) -- (934,1759);
\draw[mapOutline] (856,1756) -- (860,1761) -- (858,1766) -- (864,1771) -- (866,1771);
\draw[mapOutline] (913,1757) -- (917,1761) -- (919,1761) -- (921,1759);
\draw[mapOutline] (980,1757) -- (995,1759);
\draw[mapOutline] (1267,1757) -- (1267,1759) -- (1263,1763);
\draw[mapOutline] (911,1758) -- (908,1761);
\draw[mapOutline] (950,1758) -- (949,1759) -- (954,1761);
\draw[mapOutline] (973,1758) -- (972,1762);
\draw[mapOutline] (1539,1758) -- (1541,1762);
\draw[mapOutline] (1720,1758) -- (1720,1764) -- (1721,1765);
\draw[mapOutline] (2483,1758) -- (2486,1762);
\draw[mapOutline] (1000,1760) -- (1005,1760);
\draw[mapOutline] (958,1763) -- (961,1764) -- (965,1768) -- (967,1768) -- (971,1764);
\draw[mapOutline] (1539,1763) -- (1536,1766) -- (1537,1770) -- (1534,1774) -- (1535,1777) -- (1534,1780);
\draw[mapOutline] (903,1765) -- (891,1772);
\draw[mapOutline] (1683,1766) -- (1688,1769);
\draw[mapOutline] (2486,1766) -- (2486,1771);
\draw[mapOutline] (1261,1767) -- (1261,1772) -- (1264,1777);
\draw[mapOutline] (2354,1769) -- (2356,1780);
\draw[mapOutline] (1688,1771) -- (1686,1776) -- (1686,1782);
\draw[mapOutline] (2371,1774) -- (2370,1778);
\draw[mapOutline] (870,1775) -- (877,1779) -- (883,1780);
\draw[mapOutline] (2484,1775) -- (2482,1780);
\draw[mapOutline] (893,1779) -- (895,1782);
\draw[mapOutline] (887,1782) -- (895,1785);
\draw[mapOutline] (1266,1782) -- (1262,1785);
\draw[mapOutline] (2039,1782) -- (2037,1783) -- (2037,1787) -- (2039,1789) -- (2041,1789);
\draw[mapOutline] (2042,1783) -- (2046,1785);
\draw[mapOutline] (2440,1783) -- (2435,1783) -- (2431,1786) -- (2430,1785);
\draw[mapOutline] (1686,1784) -- (1691,1786) -- (1698,1792);
\draw[mapOutline] (2280,1784) -- (2274,1787);
\draw[mapOutline] (2286,1784) -- (2287,1788);
\draw[mapOutline] (2463,1784) -- (2459,1785);
\draw[mapOutline] (2466,1784) -- (2470,1788);
\draw[mapOutline] (1733,1785) -- (1735,1790);
\draw[mapOutline] (2442,1785) -- (2454,1786);
\draw[mapOutline] (1260,1786) -- (1256,1791) -- (1257,1794);
\draw[mapOutline] (2482,1786) -- (2486,1787);
\draw[mapOutline] (2488,1786) -- (2492,1788) -- (2493,1790) -- (2492,1793);
\draw[mapOutline] (1531,1787) -- (1530,1790) -- (1532,1793);
\draw[mapOutline] (1967,1789) -- (1961,1791);
\draw[mapOutline] (2266,1791) -- (2267,1795);
\draw[mapOutline] (2274,1791) -- (2276,1797);
\draw[mapOutline] (2352,1791) -- (2351,1792) -- (2342,1792) -- (2341,1791) -- (2334,1792);
\draw[mapOutline] (2366,1791) -- (2364,1792) -- (2364,1797) -- (2366,1798);
\draw[mapOutline] (2420,1791) -- (2416,1792);
\draw[mapOutline] (1737,1792) -- (1745,1802);
\draw[mapOutline] (1701,1793) -- (1706,1797);
\draw[mapOutline] (2332,1793) -- (2326,1797);
\draw[mapOutline] (2317,1794) -- (2311,1799) -- (2306,1799);
\draw[mapOutline] (2471,1794) -- (2469,1805);
\draw[mapOutline] (1254,1795) -- (1249,1799);
\draw[mapOutline] (2320,1795) -- (2324,1798) -- (2326,1797);
\draw[mapOutline] (1971,1796) -- (1970,1798) -- (1972,1799);
\draw[mapOutline] (2054,1797) -- (2061,1802) -- (2064,1802) -- (2068,1798) -- (2071,1798);
\draw[mapOutline] (2374,1797) -- (2380,1803);
\draw[mapOutline] (2368,1798) -- (2372,1798);
\draw[mapOutline] (1982,1800) -- (1976,1806) -- (1973,1806) -- (1972,1802);
\draw[mapOutline] (1989,1800) -- (1984,1801);
\draw[mapOutline] (1248,1801) -- (1251,1806);
\draw[mapOutline] (1992,1801) -- (1994,1806);
\draw[mapOutline] (1715,1802) -- (1717,1804) -- (1715,1813) -- (1716,1814);
\draw[mapOutline] (1746,1805) -- (1745,1809) -- (1748,1814);
\draw[mapOutline] (2401,1805) -- (2408,1810) -- (2408,1813) -- (2404,1818);
\draw[mapOutline] (2380,1807) -- (2380,1809) -- (2373,1815);
\draw[mapOutline] (2033,1808) -- (2039,1815) -- (2042,1816) -- (2044,1821);
\draw[mapOutline] (1997,1809) -- (1991,1814) -- (1985,1815) -- (1983,1819) -- (1980,1819);
\draw[mapOutline] (2470,1809) -- (2471,1813) -- (2475,1817) -- (2480,1817);
\draw[mapOutline] (1962,1810) -- (1960,1812) -- (1960,1814);
\draw[mapOutline] (2265,1811) -- (2265,1816);
\draw[mapOutline] (1708,1813) -- (1713,1814);
\draw[mapOutline] (1258,1814) -- (1261,1820);
\draw[mapOutline] (1695,1816) -- (1689,1822);
\draw[mapOutline] (2361,1816) -- (2354,1821) -- (2347,1821) -- (2345,1823) -- (2346,1826) -- (2349,1828);
\draw[mapOutline] (2363,1816) -- (2364,1817) -- (2371,1817);
\draw[mapOutline] (1544,1817) -- (1534,1821);
\draw[mapOutline] (1566,1821) -- (1570,1819) -- (1571,1820) -- (1583,1821);
\draw[mapOutline] (1585,1822) -- (1593,1824) -- (1598,1830) -- (1599,1835) -- (1607,1835) -- (1608,1834) -- (1611,1836) -- (1611,1841) -- (1616,1843) -- (1619,1842) -- (1620,1843) -- (1616,1849) -- (1618,1856);
\draw[mapOutline] (1964,1822) -- (1963,1823) -- (1964,1825);
\draw[mapOutline] (1971,1822) -- (1966,1824);
\draw[mapOutline] (2266,1823) -- (2267,1827);
\draw[mapOutline] (1563,1824) -- (1557,1828) -- (1552,1827) -- (1550,1829);
\draw[mapOutline] (1975,1824) -- (1979,1825);
\draw[mapOutline] (2046,1824) -- (2045,1827) -- (2049,1834);
\draw[mapOutline] (2116,1825) -- (2115,1820) -- (2112,1818) -- (2108,1822) -- (2099,1826);
\draw[mapOutline] (2099,1826) -- (2087,1835) -- (2083,1835) -- (2076,1829) -- (2074,1829);
\draw[mapOutline] (2114,1827) -- (2112,1829) -- (2115,1831);
\draw[mapOutline] (1538,1828) -- (1539,1830) -- (1542,1830) -- (1545,1833) -- (1547,1832);
\draw[mapOutline] (2351,1828) -- (2354,1829) -- (2356,1832);
\draw[mapOutline] (1686,1829) -- (1680,1837);
\draw[mapOutline] (2121,1831) -- (2117,1832);
\draw[mapOutline] (2126,1832) -- (2128,1835) -- (2135,1832);
\draw[mapOutline] (2403,1832) -- (2405,1835) -- (2403,1837) -- (2399,1837);
\draw[mapOutline] (1273,1833) -- (1273,1837) -- (1270,1839);
\draw[mapOutline] (2268,1833) -- (2273,1835);
\draw[mapOutline] (2359,1833) -- (2355,1834);
\draw[mapOutline] (2362,1834) -- (2366,1833) -- (2372,1838) -- (2377,1836);
\draw[mapOutline] (2387,1834) -- (2382,1837) -- (2379,1836);
\draw[mapOutline] (2390,1834) -- (2395,1839) -- (2397,1838);
\draw[mapOutline] (2048,1837) -- (2049,1841);
\draw[mapOutline] (2164,1837) -- (2166,1833) -- (2163,1832) -- (2161,1834) -- (2159,1840) -- (2156,1843) -- (2153,1844) -- (2147,1839) -- (2144,1839) -- (2138,1841) -- (2137,1845);
\draw[mapOutline] (2134,1838) -- (2132,1841) -- (2128,1841) -- (2126,1845) -- (2123,1847);
\draw[mapOutline] (2275,1838) -- (2277,1839) -- (2279,1843);
\draw[mapOutline] (2165,1840) -- (2164,1845);
\draw[mapOutline] (2041,1843) -- (2039,1843) -- (2038,1846) -- (2033,1849) -- (2033,1853) -- (2038,1855) -- (2038,1859);
\draw[mapOutline] (2048,1844) -- (2044,1845);
\draw[mapOutline] (1854,1845) -- (1856,1848);
\draw[mapOutline] (2140,1845) -- (2145,1845) -- (2157,1855) -- (2163,1851);
\draw[mapOutline] (2282,1847) -- (2282,1853);
\draw[mapOutline] (1570,1848) -- (1577,1852) -- (1584,1851);
\draw[mapOutline] (1564,1849) -- (1559,1850);
\draw[mapOutline] (1259,1850) -- (1257,1853);
\draw[mapOutline] (1630,1850) -- (1631,1856) -- (1633,1857);
\draw[mapOutline] (1661,1850) -- (1665,1850) -- (1668,1852) -- (1671,1852) -- (1673,1850);
\draw[mapOutline] (1713,1850) -- (1706,1852);
\draw[mapOutline] (1815,1850) -- (1815,1856) -- (1813,1860);
\draw[mapOutline] (1660,1853) -- (1656,1856) -- (1653,1856) -- (1648,1853) -- (1644,1855);
\draw[mapOutline] (1704,1854) -- (1705,1855) -- (1702,1860);
\draw[mapOutline] (1769,1854) -- (1776,1859) -- (1780,1865) -- (1788,1863) -- (1797,1871) -- (1800,1871) -- (1806,1865) -- (1812,1865) -- (1813,1864);
\draw[mapOutline] (1588,1855) -- (1592,1862) -- (1592,1865);
\draw[mapOutline] (1532,1857) -- (1532,1862);
\draw[mapOutline] (2283,1857) -- (2287,1858);
\draw[mapOutline] (1251,1860) -- (1246,1861);
\draw[mapOutline] (1257,1860) -- (1256,1862) -- (1253,1861);
\draw[mapOutline] (2114,1860) -- (2123,1861);
\draw[mapOutline] (2114,1860) -- (2113,1862) -- (2114,1865);
\draw[mapOutline] (1724,1862) -- (1725,1863) -- (1725,1868) -- (1721,1871);
\draw[mapOutline] (1701,1863) -- (1698,1869);
\draw[mapOutline] (2312,1864) -- (2305,1870);
\draw[mapOutline] (1592,1865) -- (1588,1868);
\draw[mapOutline] (1240,1866) -- (1236,1870);
\draw[mapOutline] (1541,1867) -- (1540,1871);
\draw[mapOutline] (1534,1868) -- (1534,1872) -- (1540,1877) -- (1541,1880);
\draw[mapOutline] (1588,1868) -- (1590,1870) -- (1590,1874);
\draw[mapOutline] (2114,1868) -- (2115,1873);
\draw[mapOutline] (1707,1872) -- (1710,1872) -- (1711,1876);
\draw[mapOutline] (2122,1872) -- (2128,1871) -- (2131,1873);
\draw[mapOutline] (2122,1872) -- (2118,1875);
\draw[mapOutline] (2039,1873) -- (2032,1879);
\draw[mapOutline] (1704,1874) -- (1700,1876);
\draw[mapOutline] (1721,1874) -- (1722,1878);
\draw[mapOutline] (1695,1875) -- (1690,1876);
\draw[mapOutline] (2131,1876) -- (2135,1880) -- (2140,1875) -- (2142,1875) -- (2149,1879) -- (2152,1879) -- (2154,1877) -- (2156,1878) -- (2158,1877) -- (2165,1882);
\draw[mapOutline] (1684,1880) -- (1684,1883) -- (1682,1886) -- (1685,1889);
\draw[mapOutline] (1983,1881) -- (1985,1879) -- (1987,1881);
\draw[mapOutline] (2032,1881) -- (2034,1889);
\draw[mapOutline] (1702,1882) -- (1704,1885) -- (1703,1888);
\draw[mapOutline] (1721,1884) -- (1724,1889);
\draw[mapOutline] (1979,1884) -- (1975,1887) -- (1973,1887);
\draw[mapOutline] (1621,1885) -- (1618,1885) -- (1613,1888);
\draw[mapOutline] (1991,1885) -- (1994,1888);
\draw[mapOutline] (2166,1886) -- (2166,1891);
\draw[mapOutline] (1689,1889) -- (1691,1893) -- (1694,1895);
\draw[mapOutline] (1710,1889) -- (1711,1895);
\draw[mapOutline] (1703,1891) -- (1702,1895);
\draw[mapOutline] (1255,1893) -- (1251,1894);
\draw[mapOutline] (1632,1893) -- (1633,1894) -- (1639,1894);
\draw[mapOutline] (1260,1894) -- (1265,1901);
\draw[mapOutline] (1724,1894) -- (1724,1896) -- (1722,1897);
\draw[mapOutline] (1974,1894) -- (1974,1899) -- (1968,1905) -- (1974,1911);
\draw[mapOutline] (1702,1895) -- (1701,1896) -- (1696,1896);
\draw[mapOutline] (1594,1896) -- (1590,1897);
\draw[mapOutline] (1608,1896) -- (1610,1899) -- (1610,1902) -- (1612,1904);
\draw[mapOutline] (1663,1897) -- (1675,1907);
\draw[mapOutline] (2175,1897) -- (2179,1902) -- (2182,1903) -- (2185,1906) -- (2177,1914) -- (2178,1916) -- (2181,1916);
\draw[mapOutline] (2175,1897) -- (2174,1896) -- (2175,1895) -- (2176,1896) -- cycle;
\draw[mapOutline] (1532,1898) -- (1539,1899);
\draw[mapOutline] (1541,1898) -- (1543,1900) -- (1547,1901);
\draw[mapOutline] (1563,1899) -- (1557,1900) -- (1554,1905) -- (1551,1904);
\draw[mapOutline] (2265,1899) -- (2285,1900) -- (2286,1901);
\draw[mapOutline] (2260,1900) -- (2252,1902) -- (2249,1904) -- (2242,1904);
\draw[mapOutline] (1878,1901) -- (1880,1902) -- (1884,1898) -- (1886,1898) -- (1888,1901);
\draw[mapOutline] (2007,1901) -- (2011,1905);
\draw[mapOutline] (2035,1902) -- (2035,1906);
\draw[mapOutline] (1614,1903) -- (1617,1909);
\draw[mapOutline] (2234,1903) -- (2239,1905) -- (2242,1904);
\draw[mapOutline] (2103,1904) -- (2110,1901) -- (2120,1910) -- (2123,1915);
\draw[mapOutline] (1888,1905) -- (1889,1908) -- (1892,1909);
\draw[mapOutline] (1509,1906) -- (1506,1908);
\draw[mapOutline] (2303,1906) -- (2308,1910) -- (2310,1916);
\draw[mapOutline] (1677,1907) -- (1678,1910) -- (1682,1912) -- (1683,1916);
\draw[mapOutline] (1267,1908) -- (1273,1915) -- (1280,1920);
\draw[mapOutline] (2102,1908) -- (2101,1911) -- (2102,1914);
\draw[mapOutline] (1498,1910) -- (1495,1914);
\draw[mapOutline] (2037,1910) -- (2039,1917) -- (2041,1919) -- (2044,1919) -- (2046,1922);
\draw[mapOutline] (1507,1912) -- (1510,1913) -- (1509,1918);
\draw[mapOutline] (1903,1914) -- (1908,1913) -- (1910,1915);
\draw[mapOutline] (1974,1914) -- (1970,1916);
\draw[mapOutline] (2181,1916) -- (2182,1920);
\draw[mapOutline] (1920,1917) -- (1915,1922) -- (1911,1918);
\draw[mapOutline] (2220,1917) -- (2218,1921) -- (2216,1920) -- (2213,1921);
\draw[mapOutline] (1965,1918) -- (1961,1918) -- (1958,1920) -- (1954,1926);
\draw[mapOutline] (2100,1918) -- (2101,1922) -- (2099,1923);
\draw[mapOutline] (1684,1919) -- (1686,1921) -- (1687,1927);
\draw[mapOutline] (2311,1919) -- (2313,1925);
\draw[mapOutline] (1943,1920) -- (1945,1924);
\draw[mapOutline] (2123,1920) -- (2124,1924) -- (2123,1925);
\draw[mapOutline] (2179,1921) -- (2175,1925);
\draw[mapOutline] (2203,1921) -- (2198,1924);
\draw[mapOutline] (2018,1922) -- (2022,1923);
\draw[mapOutline] (1282,1924) -- (1285,1927);
\draw[mapOutline] (2098,1926) -- (2095,1927) -- (2096,1929);
\draw[mapOutline] (2120,1927) -- (2123,1930);
\draw[mapOutline] (2181,1927) -- (2175,1930);
\draw[mapOutline] (1487,1928) -- (1483,1930);
\draw[mapOutline] (2313,1928) -- (2309,1933);
\draw[mapOutline] (2018,1929) -- (2019,1933);
\draw[mapOutline] (2189,1929) -- (2190,1930) -- (2193,1929);
\draw[mapOutline] (1686,1930) -- (1681,1934) -- (1681,1936) -- (1685,1938);
\draw[mapOutline] (1507,1932) -- (1500,1940) -- (1500,1943) -- (1503,1946);
\draw[mapOutline] (1479,1935) -- (1475,1938);
\draw[mapOutline] (2020,1936) -- (2020,1938) -- (2022,1940) -- (2024,1940);
\draw[mapOutline] (2308,1936) -- (2302,1941);
\draw[mapOutline] (1472,1938) -- (1464,1941);
\draw[mapOutline] (1462,1941) -- (1454,1943);
\draw[mapOutline] (1445,1942) -- (1448,1941) -- (1450,1943) -- (1452,1943);
\draw[mapOutline] (1510,1943) -- (1506,1946);
\draw[mapOutline] (1515,1944) -- (1518,1944) -- (1522,1947);
\draw[mapOutline] (1830,1945) -- (1827,1945) -- (1821,1942) -- (1814,1943) -- (1811,1946) -- (1813,1951) -- (1810,1954);
\draw[mapOutline] (2261,1945) -- (2261,1947) -- (2264,1949);
\draw[mapOutline] (1832,1946) -- (1834,1953) -- (1837,1956) -- (1840,1955);
\draw[mapOutline] (2439,1947) -- (2438,1948) -- (2437,1947) -- (2435,1950);
\draw[mapOutline] (1558,1949) -- (1562,1949);
\draw[mapOutline] (2196,1950) -- (2199,1949) -- (2201,1951) -- (2202,1958) -- (2204,1961);
\draw[mapOutline] (1522,1951) -- (1524,1954) -- (1528,1953);
\draw[mapOutline] (1444,1952) -- (1443,1951) -- (1444,1950) -- (1445,1951) -- cycle;
\draw[mapOutline] (1298,1953) -- (1301,1957);
\draw[mapOutline] (2282,1954) -- (2278,1958) -- (2278,1960);
\draw[mapOutline] (2435,1954) -- (2432,1962) -- (2427,1967) -- (2422,1968) -- (2421,1970);
\draw[mapOutline] (1842,1955) -- (1846,1957) -- (1849,1957);
\draw[mapOutline] (2235,1955) -- (2231,1958);
\draw[mapOutline] (2237,1955) -- (2239,1956) -- (2238,1958);
\draw[mapOutline] (1681,1956) -- (1679,1957) -- (1676,1964);
\draw[mapOutline] (1438,1957) -- (1433,1961) -- (1432,1964) -- (1425,1967) -- (1424,1966);
\draw[mapOutline] (1807,1957) -- (1802,1959);
\draw[mapOutline] (2442,1957) -- (2440,1959) -- (2441,1965);
\draw[mapOutline] (1659,1958) -- (1657,1960) -- (1655,1960);
\draw[mapOutline] (2222,1961) -- (2214,1961) -- (2210,1963);
\draw[mapOutline] (1303,1964) -- (1304,1968);
\draw[mapOutline] (2000,1964) -- (1997,1966);
\draw[mapOutline] (2006,1964) -- (2012,1966);
\draw[mapOutline] (2018,1964) -- (2014,1965);
\draw[mapOutline] (2093,1965) -- (2103,1970);
\draw[mapOutline] (2027,1966) -- (2027,1969) -- (2023,1972);
\draw[mapOutline] (2455,1966) -- (2460,1969);
\draw[mapOutline] (2455,1966) -- (2448,1972) -- (2441,1972) -- (2440,1971);
\draw[mapOutline] (1905,1967) -- (1908,1969);
\draw[mapOutline] (2481,1967) -- (2486,1968);
\draw[mapOutline] (1417,1968) -- (1422,1968);
\draw[mapOutline] (1902,1968) -- (1900,1970) -- (1898,1970);
\draw[mapOutline] (2477,1968) -- (2475,1972) -- (2472,1971) -- (2468,1974) -- (2467,1973);
\draw[mapOutline] (1997,1969) -- (2004,1970);
\draw[mapOutline] (1412,1970) -- (1407,1972);
\draw[mapOutline] (2517,1970) -- (2520,1973) -- (2519,1975);
\draw[mapOutline] (1914,1971) -- (1918,1972) -- (1924,1979);
\draw[mapOutline] (2489,1972) -- (2493,1972);
\draw[mapOutline] (1803,1973) -- (1798,1979);
\draw[mapOutline] (2003,1973) -- (2004,1974) -- (2002,1977) -- (2005,1979);
\draw[mapOutline] (2460,1973) -- (2458,1976) -- (2459,1978) -- (2465,1974);
\draw[mapOutline] (2512,1973) -- (2510,1975) -- (2510,1977);
\draw[mapOutline] (1305,1974) -- (1305,1978);
\draw[mapOutline] (1406,1974) -- (1403,1976) -- (1401,1980);
\draw[mapOutline] (2496,1974) -- (2498,1976) -- (2502,1975);
\draw[mapOutline] (2523,1975) -- (2524,1976) -- (2523,1977) -- (2519,1977);
\draw[mapOutline] (2526,1976) -- (2530,1978);
\draw[mapOutline] (2505,1977) -- (2508,1979);
\draw[mapOutline] (1783,1980) -- (1779,1982) -- (1775,1987);
\draw[mapOutline] (1792,1980) -- (1787,1984);
\draw[mapOutline] (2039,1981) -- (2046,1985);
\draw[mapOutline] (2246,1981) -- (2249,1985) -- (2251,1984);
\draw[mapOutline] (2253,1981) -- (2261,1983) -- (2265,1989);
\draw[mapOutline] (2525,1981) -- (2521,1983) -- (2520,1985);
\draw[mapOutline] (1401,1983) -- (1396,1984) -- (1393,1986) -- (1395,1990);
\draw[mapOutline] (1306,1984) -- (1308,1990) -- (1310,1991);
\draw[mapOutline] (1670,1984) -- (1669,1988);
\draw[mapOutline] (2085,1988) -- (2087,1988) -- (2091,1994);
\draw[mapOutline] (2675,1988) -- (2671,1990) -- (2667,1987) -- (2665,1989);
\draw[mapOutline] (1667,1992) -- (1668,1999) -- (1671,2002);
\draw[mapOutline] (2664,1993) -- (2660,1995) -- (2657,1995);
\draw[mapOutline] (1312,1994) -- (1323,2003) -- (1330,1997);
\draw[mapOutline] (1332,1995) -- (1334,1995) -- (1338,1998);
\draw[mapOutline] (1381,1995) -- (1384,1995) -- (1388,1998);
\draw[mapOutline] (1770,1995) -- (1765,1997);
\draw[mapOutline] (1379,1996) -- (1374,2002) -- (1372,2001);
\draw[mapOutline] (1393,1996) -- (1390,1999);
\draw[mapOutline] (2655,1997) -- (2649,2001);
\draw[mapOutline] (2094,1999) -- (2087,2002);
\draw[mapOutline] (1341,2000) -- (1343,2002) -- (1345,2001) -- (1348,2002);
\draw[mapOutline] (1370,2001) -- (1366,2002);
\draw[mapOutline] (2515,2004) -- (2513,2010) -- (2516,2011) -- (2518,2014);
\draw[mapOutline] (2648,2004) -- (2647,2010) -- (2640,2015);
\draw[mapOutline] (1751,2006) -- (1746,2011);
\draw[mapOutline] (2054,2007) -- (2051,2010);
\draw[mapOutline] (2102,2009) -- (2096,2015);
\draw[mapOutline] (2410,2009) -- (2407,2012);
\draw[mapOutline] (2107,2010) -- (2110,2015);
\draw[mapOutline] (2391,2013) -- (2388,2015);
\draw[mapOutline] (2410,2013) -- (2413,2018);
\draw[mapOutline] (1967,2014) -- (1969,2018);
\draw[mapOutline] (2051,2014) -- (2057,2016) -- (2060,2019) -- (2060,2023) -- (2062,2028) -- (2068,2034) -- (2071,2035) -- (2071,2038);
\draw[mapOutline] (2110,2015) -- (2105,2017) -- (2106,2018);
\draw[mapOutline] (2383,2015) -- (2380,2017);
\draw[mapOutline] (2371,2020) -- (2375,2020);
\draw[mapOutline] (2635,2020) -- (2639,2021);
\draw[mapOutline] (1740,2021) -- (1733,2025);
\draw[mapOutline] (2413,2021) -- (2413,2025);
\draw[mapOutline] (1943,2022) -- (1946,2022) -- (1948,2025);
\draw[mapOutline] (2368,2023) -- (2365,2027);
\draw[mapOutline] (1969,2024) -- (1966,2028) -- (1963,2028);
\draw[mapOutline] (1948,2025) -- (1947,2027) -- (1955,2033) -- (1958,2033) -- (1960,2030);
\draw[mapOutline] (2516,2026) -- (2518,2031);
\draw[mapOutline] (2338,2027) -- (2339,2028) -- (2344,2028);
\draw[mapOutline] (1731,2029) -- (1728,2033);
\draw[mapOutline] (2360,2029) -- (2357,2031);
\draw[mapOutline] (2637,2030) -- (2641,2032);
\draw[mapOutline] (2339,2031) -- (2339,2033) -- (2342,2035);
\draw[mapOutline] (2411,2031) -- (2413,2035);
\draw[mapOutline] (1726,2035) -- (1721,2038);
\draw[mapOutline] (2508,2035) -- (2507,2039);
\draw[mapOutline] (2417,2037) -- (2418,2041);
\draw[mapOutline] (2607,2037) -- (2613,2038);
\draw[mapOutline] (2345,2038) -- (2345,2042) -- (2347,2044) -- (2347,2047);
\draw[mapOutline] (2604,2038) -- (2602,2040) -- (2603,2044);
\draw[mapOutline] (2615,2038) -- (2620,2038);
\draw[mapOutline] (2625,2038) -- (2629,2038);
\draw[mapOutline] (2631,2038) -- (2638,2038);
\draw[mapOutline] (1654,2039) -- (1651,2042);
\draw[mapOutline] (2093,2039) -- (2098,2040);
\draw[mapOutline] (2685,2040) -- (2689,2040);
\draw[mapOutline] (1642,2041) -- (1645,2043) -- (1648,2042);
\draw[mapOutline] (2680,2041) -- (2677,2043) -- (2678,2045);
\draw[mapOutline] (1717,2042) -- (1712,2047);
\draw[mapOutline] (2509,2042) -- (2511,2047);
\draw[mapOutline] (2071,2044) -- (2076,2046);
\draw[mapOutline] (2089,2044) -- (2094,2044);
\draw[mapOutline] (2337,2044) -- (2341,2046);
\draw[mapOutline] (2689,2044) -- (2693,2044) -- (2697,2046);
\draw[mapOutline] (1639,2045) -- (1637,2048);
\draw[mapOutline] (2078,2045) -- (2081,2044) -- (2086,2045);
\draw[mapOutline] (2304,2045) -- (2306,2046) -- (2307,2049);
\draw[mapOutline] (2402,2045) -- (2407,2046);
\draw[mapOutline] (2681,2046) -- (2685,2046);
\draw[mapOutline] (2709,2046) -- (2717,2047);
\draw[mapOutline] (2511,2047) -- (2509,2050) -- (2510,2055);
\draw[mapOutline] (2707,2047) -- (2706,2052) -- (2704,2052) -- (2699,2048);
\draw[mapOutline] (2581,2048) -- (2576,2052);
\draw[mapOutline] (2584,2048) -- (2585,2047) -- (2595,2048);
\draw[mapOutline] (2397,2049) -- (2398,2054) -- (2400,2054) -- (2403,2051);
\draw[mapOutline] (2597,2049) -- (2598,2050) -- (2603,2049);
\draw[mapOutline] (2688,2051) -- (2685,2052) -- (2679,2051) -- (2675,2054);
\draw[mapOutline] (2307,2053) -- (2308,2058);
\draw[mapOutline] (2406,2053) -- (2405,2057);
\draw[mapOutline] (2573,2054) -- (2572,2060);
\draw[mapOutline] (2506,2056) -- (2504,2064);
\draw[mapOutline] (2675,2057) -- (2674,2062);
\draw[mapOutline] (2405,2060) -- (2397,2065) -- (2394,2064);
\draw[mapOutline] (2390,2064) -- (2387,2067) -- (2386,2070);
\draw[mapOutline] (2504,2064) -- (2501,2068) -- (2499,2068);
\draw[mapOutline] (2670,2064) -- (2664,2069);
\draw[mapOutline] (2573,2065) -- (2577,2068);
\draw[mapOutline] (2487,2067) -- (2486,2068) -- (2487,2075);
\draw[mapOutline] (2490,2067) -- (2494,2067);
\draw[mapOutline] (1707,2068) -- (1700,2071);
\draw[mapOutline] (2309,2069) -- (2311,2072);
\draw[mapOutline] (2662,2069) -- (2658,2070);
\draw[mapOutline] (1645,2071) -- (1643,2075);
\draw[mapOutline] (2578,2073) -- (2574,2077) -- (2574,2079);
\draw[mapOutline] (2375,2074) -- (2380,2075) -- (2381,2074);
\draw[mapOutline] (2488,2078) -- (2488,2082);
\draw[mapOutline] (2656,2082) -- (2652,2083);
\draw[mapOutline] (2344,2084) -- (2351,2081) -- (2353,2081) -- (2357,2084);
\draw[mapOutline] (2559,2084) -- (2556,2087) -- (2555,2091);
\draw[mapOutline] (2650,2084) -- (2644,2091);
\draw[mapOutline] (1683,2085) -- (1681,2088) -- (1678,2089);
\draw[mapOutline] (2489,2088) -- (2490,2094);
\draw[mapOutline] (1649,2091) -- (1644,2104) -- (1642,2105);
\draw[mapOutline] (1676,2092) -- (1670,2097);
\draw[mapOutline] (2553,2094) -- (2551,2099);
\draw[mapOutline] (2643,2094) -- (2641,2097);
\draw[mapOutline] (2549,2098) -- (2546,2098) -- (2542,2100);
\draw[mapOutline] (1669,2100) -- (1661,2108) -- (1650,2107) -- (1646,2110);
\draw[mapOutline] (2639,2100) -- (2633,2104);
\draw[mapOutline] (2500,2101) -- (2501,2105);
\draw[mapOutline] (2541,2103) -- (2539,2104) -- (2539,2106);
\draw[mapOutline] (2503,2104) -- (2507,2102) -- (2511,2105);
\draw[mapOutline] (2628,2104) -- (2623,2105);
\draw[mapOutline] (2511,2105) -- (2513,2104) -- (2515,2106);
\draw[mapOutline] (2530,2106) -- (2527,2112);
\draw[mapOutline] (2167,2108) -- (2168,2110) -- (2167,2112);
\draw[mapOutline] (2516,2110) -- (2515,2111) -- (2510,2111) -- (2503,2115);
\draw[mapOutline] (1644,2115) -- (1639,2119);
\draw[mapOutline] (2501,2115) -- (2496,2118);
\draw[mapOutline] (2527,2115) -- (2522,2117);
\draw[mapOutline] (2493,2117) -- (2491,2119) -- (2492,2121);
\draw[mapOutline] (2363,2119) -- (2366,2120) -- (2373,2119);
\draw[mapOutline] (2482,2119) -- (2481,2120) -- (2481,2125);
\draw[mapOutline] (2388,2120) -- (2389,2121) -- (2388,2124);
\draw[mapOutline] (2519,2122) -- (2514,2126);
\draw[mapOutline] (2490,2123) -- (2492,2125) -- (2500,2125) -- (2503,2127);
\draw[mapOutline] (2618,2123) -- (2618,2129) -- (2619,2130) -- (2623,2129);
\draw[mapOutline] (2388,2127) -- (2382,2131);
\draw[mapOutline] (2327,2130) -- (2330,2131) -- (2336,2138);
\draw[mapOutline] (1633,2131) -- (1637,2134) -- (1645,2136);
\draw[mapOutline] (2626,2132) -- (2626,2134) -- (2629,2135);
\draw[mapOutline] (2692,2133) -- (2687,2134);
\draw[mapOutline] (1686,2134) -- (1691,2134);
\draw[mapOutline] (1696,2134) -- (1699,2134) -- (1701,2138);
\draw[mapOutline] (2697,2134) -- (2702,2135);
\draw[mapOutline] (1668,2136) -- (1663,2137);
\draw[mapOutline] (2631,2136) -- (2633,2138) -- (2633,2142) -- (2630,2143);
\draw[mapOutline] (1651,2137) -- (1657,2138);
\draw[mapOutline] (1671,2137) -- (1675,2140);
\draw[mapOutline] (2464,2137) -- (2468,2138);
\draw[mapOutline] (2464,2137) -- (2462,2139) -- (2462,2141);
\draw[mapOutline] (2683,2137) -- (2678,2139);
\draw[mapOutline] (1675,2140) -- (1670,2148) -- (1663,2154);
\draw[mapOutline] (2676,2140) -- (2668,2146);
\draw[mapOutline] (1729,2146) -- (1723,2149) -- (1716,2156) -- (1713,2157);
\draw[mapOutline] (1698,2147) -- (1697,2154);
\draw[mapOutline] (2697,2148) -- (2691,2155) -- (2686,2156);
\draw[mapOutline] (2659,2149) -- (2653,2154) -- (2651,2154);
\draw[mapOutline] (2636,2150) -- (2634,2152) -- (2630,2151);
\draw[mapOutline] (2642,2150) -- (2645,2153);
\draw[mapOutline] (2081,2151) -- (2080,2150) -- (2079,2152) -- (2082,2156);
\draw[mapOutline] (1884,2152) -- (1882,2151) -- (1883,2150) -- cycle;
\draw[mapOutline] (1884,2152) -- (1886,2153) -- (1884,2155);
\draw[mapOutline] (1657,2155) -- (1651,2158);
\draw[mapOutline] (1698,2157) -- (1702,2159);
\draw[mapOutline] (2680,2157) -- (2675,2158);
\draw[mapOutline] (1711,2158) -- (1710,2159) -- (1704,2159);
\draw[mapOutline] (2605,2158) -- (2609,2159);
\draw[mapOutline] (1757,2159) -- (1763,2159);
\draw[mapOutline] (2084,2159) -- (2092,2172) -- (2094,2173);
\draw[mapOutline] (2405,2159) -- (2405,2163) -- (2406,2164) -- (2409,2163);
\draw[mapOutline] (2413,2159) -- (2411,2162);
\draw[mapOutline] (2462,2159) -- (2463,2160) -- (2468,2159);
\draw[mapOutline] (2669,2159) -- (2665,2160);
\draw[mapOutline] (2326,2161) -- (2323,2161) -- (2320,2164) -- (2321,2167) -- (2325,2170);
\draw[mapOutline] (2643,2161) -- (2642,2162) -- (2636,2162);
\draw[mapOutline] (2645,2161) -- (2653,2162);
\draw[mapOutline] (2660,2161) -- (2656,2162);
\draw[mapOutline] (2594,2163) -- (2592,2165) -- (2591,2169);
\draw[mapOutline] (2626,2163) -- (2631,2163);
\draw[mapOutline] (2608,2164) -- (2613,2165);
\draw[mapOutline] (2317,2165) -- (2313,2167);
\draw[mapOutline] (2451,2165) -- (2451,2168) -- (2449,2169);
\draw[mapOutline] (2615,2165) -- (2616,2166) -- (2622,2165);
\draw[mapOutline] (2439,2167) -- (2442,2170) -- (2444,2170);
\draw[mapOutline] (1781,2168) -- (1780,2171) -- (1782,2177);
\draw[mapOutline] (2259,2170) -- (2262,2173) -- (2264,2173);
\draw[mapOutline] (2584,2171) -- (2581,2169) -- (2578,2169) -- (2566,2173);
\draw[mapOutline] (2592,2171) -- (2590,2173) -- (2587,2172);
\draw[mapOutline] (2326,2172) -- (2329,2178);
\draw[mapOutline] (2556,2172) -- (2560,2172);
\draw[mapOutline] (2449,2173) -- (2444,2183);
\draw[mapOutline] (2330,2180) -- (2329,2183) -- (2332,2185);
\draw[mapOutline] (2315,2181) -- (2313,2184);
\draw[mapOutline] (2310,2186) -- (2309,2187) -- (2311,2191) -- (2313,2192);
\draw[mapOutline] (2250,2188) -- (2252,2191);
\draw[mapOutline] (2333,2189) -- (2335,2194);
\draw[mapOutline] (2439,2190) -- (2435,2193);
\draw[mapOutline] (1777,2199) -- (1782,2201);
\draw[mapOutline] (2209,2199) -- (2214,2201);
\draw[mapOutline] (2253,2203) -- (2254,2206) -- (2252,2209);
\draw[mapOutline] (2283,2203) -- (2287,2204);
\draw[mapOutline] (1770,2206) -- (1772,2210);
\draw[mapOutline] (2287,2206) -- (2289,2208) -- (2294,2207);
\draw[mapOutline] (2296,2207) -- (2300,2208);
\draw[mapOutline] (2282,2208) -- (2285,2212) -- (2282,2215);
\draw[mapOutline] (2243,2211) -- (2246,2211) -- (2249,2213) -- (2251,2212);
\draw[mapOutline] (2241,2212) -- (2235,2217);
\draw[mapOutline] (2352,2212) -- (2357,2215);
\draw[mapOutline] (1774,2213) -- (1776,2217);
\draw[mapOutline] (2360,2213) -- (2362,2215) -- (2362,2217) -- (2361,2218) -- (2358,2217);
\draw[mapOutline] (2367,2213) -- (2370,2215) -- (2371,2218);
\draw[mapOutline] (2227,2219) -- (2225,2222) -- (2220,2225);
\draw[mapOutline] (2309,2219) -- (2310,2221) -- (2314,2221) -- (2318,2223) -- (2322,2222);
\draw[mapOutline] (1779,2220) -- (1778,2224) -- (1774,2224);
\draw[mapOutline] (2234,2221) -- (2233,2224) -- (2235,2227);
\draw[mapOutline] (1604,2222) -- (1602,2224) -- (1601,2230) -- (1598,2231);
\draw[mapOutline] (1760,2223) -- (1758,2225) -- (1759,2229);
\draw[mapOutline] (1763,2223) -- (1766,2227) -- (1772,2225);
\draw[mapOutline] (2375,2228) -- (2378,2231);
\draw[mapOutline] (2431,2230) -- (2435,2236);
\draw[mapOutline] (1760,2232) -- (1759,2233) -- (1762,2238);
\draw[mapOutline] (2160,2233) -- (2160,2235) -- (2158,2237);
\draw[mapOutline] (2381,2236) -- (2382,2240);
\draw[mapOutline] (2436,2238) -- (2440,2241);
\draw[mapOutline] (2448,2239) -- (2443,2241);
\draw[mapOutline] (1597,2240) -- (1598,2244);
\draw[mapOutline] (1763,2241) -- (1764,2242) -- (1760,2248);
\draw[mapOutline] (2172,2241) -- (2172,2245);
\draw[mapOutline] (2450,2242) -- (2450,2244) -- (2452,2246);
\draw[mapOutline] (2413,2245) -- (2408,2246) -- (2403,2251) -- (2401,2250) -- (2402,2247) -- (2401,2246) -- (2393,2247) -- (2388,2251);
\draw[mapOutline] (2415,2246) -- (2419,2248);
\draw[mapOutline] (2382,2247) -- (2385,2250);
\draw[mapOutline] (1595,2248) -- (1592,2252);
\draw[mapOutline] (2452,2248) -- (2450,2255) -- (2458,2263) -- (2463,2264);
\draw[mapOutline] (1753,2249) -- (1757,2249);
\draw[mapOutline] (1748,2250) -- (1746,2253);
\draw[mapOutline] (2418,2253) -- (2417,2258);
\draw[mapOutline] (1589,2254) -- (1585,2254) -- (1579,2257);
\draw[mapOutline] (1420,2259) -- (1424,2259);
\draw[mapOutline] (1746,2259) -- (1748,2262);
\draw[mapOutline] (1414,2261) -- (1409,2262);
\draw[mapOutline] (1443,2261) -- (1446,2264);
\draw[mapOutline] (1473,2262) -- (1466,2269) -- (1462,2268);
\draw[mapOutline] (1476,2263) -- (1479,2268) -- (1479,2272);
\draw[mapOutline] (2410,2264) -- (2411,2269);
\draw[mapOutline] (1400,2265) -- (1404,2265);
\draw[mapOutline] (1453,2265) -- (1457,2265) -- (1460,2268);
\draw[mapOutline] (1395,2266) -- (1388,2268);
\draw[mapOutline] (2464,2266) -- (2465,2267) -- (2464,2271);
\draw[mapOutline] (1563,2268) -- (1562,2272);
\draw[mapOutline] (1384,2269) -- (1380,2271);
\draw[mapOutline] (2412,2272) -- (2421,2281) -- (2425,2282);
\draw[mapOutline] (2474,2272) -- (2477,2272) -- (2480,2275);
\draw[mapOutline] (1373,2273) -- (1372,2278);
\draw[mapOutline] (2464,2274) -- (2465,2275) -- (2470,2274);
\draw[mapOutline] (2480,2275) -- (2479,2279) -- (2485,2286);
\draw[mapOutline] (1478,2276) -- (1478,2281);
\draw[mapOutline] (1549,2281) -- (1551,2284);
\draw[mapOutline] (2433,2282) -- (2429,2283);
\draw[mapOutline] (2436,2284) -- (2442,2286);
\draw[mapOutline] (1371,2286) -- (1372,2290);
\draw[mapOutline] (1551,2288) -- (1551,2293) -- (1552,2294);
\draw[mapOutline] (2488,2288) -- (2498,2296);
\draw[mapOutline] (2445,2289) -- (2451,2291);
\draw[mapOutline] (1766,2292) -- (1770,2297);
\draw[mapOutline] (2454,2292) -- (2461,2296);
\draw[mapOutline] (1535,2295) -- (1542,2299);
\draw[mapOutline] (2065,2295) -- (2062,2296) -- (2064,2301);
\draw[mapOutline] (2080,2296) -- (2080,2301) -- (2083,2302);
\draw[mapOutline] (2463,2297) -- (2466,2299);
\draw[mapOutline] (2087,2298) -- (2089,2298) -- (2091,2301) -- (2093,2301) -- (2095,2299);
\draw[mapOutline] (1544,2300) -- (1550,2300);
\draw[mapOutline] (1527,2301) -- (1519,2305);
\draw[mapOutline] (2495,2301) -- (2492,2301) -- (2490,2303);
\draw[mapOutline] (2472,2302) -- (2478,2305);
\draw[mapOutline] (2064,2304) -- (2058,2306);
\draw[mapOutline] (1473,2305) -- (1474,2309);
\draw[mapOutline] (2489,2306) -- (2487,2310) -- (2480,2307);
\draw[mapOutline] (1516,2308) -- (1517,2309) -- (1516,2310) -- (1516,2315);
\draw[mapOutline] (2061,2310) -- (2063,2312) -- (2063,2316) -- (2057,2321);
\draw[mapOutline] (1474,2313) -- (1474,2316) -- (1476,2320);
\draw[mapOutline] (1516,2317) -- (1514,2319) -- (1512,2319);
\draw[mapOutline] (1498,2318) -- (1506,2319);
\draw[mapOutline] (1476,2320) -- (1480,2320);
\draw[mapOutline] (2047,2320) -- (2052,2321);
\draw[mapOutline] (1464,2321) -- (1461,2323) -- (1463,2328);
\draw[mapOutline] (1486,2321) -- (1492,2321);
\draw[mapOutline] (2041,2322) -- (2041,2324) -- (2045,2329) -- (2045,2332) -- (2042,2333);
\draw[mapOutline] (1463,2330) -- (1463,2335);
\draw[mapOutline] (2033,2330) -- (2039,2334);
\draw[mapOutline] (1402,2331) -- (1403,2335) -- (1407,2339) -- (1405,2341);
\draw[mapOutline] (1462,2337) -- (1465,2339) -- (1469,2339);
\draw[mapOutline] (1471,2338) -- (1473,2338) -- (1475,2341);
\draw[mapOutline] (1399,2340) -- (1396,2342);
\draw[mapOutline] (1477,2345) -- (1481,2349) -- (1480,2351);
\draw[mapOutline] (1765,2346) -- (1764,2348) -- (1766,2349);
\draw[mapOutline] (1857,2349) -- (1853,2352);
\draw[mapOutline] (1897,2351) -- (1897,2353) -- (1899,2355);
\draw[mapOutline] (1965,2351) -- (1959,2356) -- (1957,2361);
\draw[mapOutline] (1764,2352) -- (1763,2353) -- (1757,2352) -- (1754,2358);
\draw[mapOutline] (1476,2354) -- (1477,2353) -- (1480,2354);
\draw[mapOutline] (1918,2356) -- (1921,2358);
\draw[mapOutline] (1377,2357) -- (1383,2357);
\draw[mapOutline] (1390,2357) -- (1391,2356) -- (1395,2357);
\draw[mapOutline] (1842,2357) -- (1841,2356) -- (1840,2358);
\draw[mapOutline] (1371,2358) -- (1365,2359);
\draw[mapOutline] (1475,2358) -- (1476,2363) -- (1479,2364);
\draw[mapOutline] (1746,2358) -- (1749,2361) -- (1751,2360);
\draw[mapOutline] (1840,2358) -- (1838,2359) -- (1837,2358) -- (1833,2360);
\draw[mapOutline] (1833,2360) -- (1832,2362) -- (1831,2361);
\draw[mapOutline] (1954,2361) -- (1949,2363);
\draw[mapOutline] (1862,2363) -- (1860,2370) -- (1857,2370) -- (1855,2367);
\draw[mapOutline] (1924,2363) -- (1927,2366);
\draw[mapOutline] (1360,2364) -- (1351,2371) -- (1343,2370) -- (1342,2365);
\draw[mapOutline] (1481,2365) -- (1489,2370) -- (1492,2368) -- (1493,2365) -- (1495,2365) -- (1497,2368);
\draw[mapOutline] (1741,2365) -- (1734,2371);
\draw[mapOutline] (1944,2366) -- (1938,2373) -- (1937,2378);
\draw[mapOutline] (1497,2371) -- (1498,2373) -- (1494,2377) -- (1496,2379);
\draw[mapOutline] (1840,2372) -- (1841,2373) -- (1844,2372);
\draw[mapOutline] (1835,2375) -- (1834,2383);
\draw[mapOutline] (1496,2379) -- (1496,2381) -- (1494,2383) -- (1495,2384);
\draw[mapOutline] (1821,2383) -- (1824,2385) -- (1827,2385) -- (1831,2383);
\draw[mapOutline] (1925,2390) -- (1924,2391) -- (1925,2394) -- (1923,2396);
\draw[mapOutline] (1486,2397) -- (1497,2407);
\draw[mapOutline] (1497,2412) -- (1491,2417);
\draw[mapOutline] (1489,2418) -- (1487,2422) -- (1486,2433);
\draw[mapOutline] (1837,2421) -- (1841,2421) -- (1843,2424);
\draw[mapOutline] (1914,2425) -- (1917,2428);
\draw[mapOutline] (1841,2429) -- (1842,2434);
\draw[mapOutline] (1971,2429) -- (1973,2428) -- (1979,2429);
\draw[mapOutline] (1812,491) -- (1814,489) -- (1817,490) -- (1824,490) -- (1825,489) -- (1830,490) -- (1833,489) -- (1834,494) -- (1830,496) -- (1817,496) -- cycle;

\draw[mapOutline] (1205,158) -- (1201,153) -- (1203,146) -- (1206,142) --
  (1210,140) -- (1214,142) -- (1217,146) -- (1220,152) -- (1218,157);
\draw[mapOutline] (1206,149) -- (1210,147) -- (1214,149) -- (1218,152);
% Geographical strokes visible at survey-line and lettering crossings.
\draw[mapOutline] (783,266) -- (783,269) -- (781,270) -- (781,272) -- (777,274) -- (777,276);
\draw[mapOutline] (1642,270) -- (1644,274) -- (1641,277);
\draw[mapOutline] (1642,279) -- (1647,279) -- (1650,285);
\draw[mapOutline] (1668,257) -- (1669,259) -- (1669,266) -- (1670,267) -- (1670,270) -- (1672,272) -- (1671,277);
\draw[mapOutline] (646,339) -- (636,340);
\draw[mapOutline] (661,339) -- (663,339);
\draw[mapOutline] (666,339) -- (667,339) -- (668,341) -- (669,340) -- (673,341);
\draw[mapOutline] (684,345) -- (685,345) -- (686,347) -- (689,347) -- (695,350);
\draw[mapOutline] (712,347) -- (711,349) -- (708,349) -- (701,353);
\draw[mapOutline] (736,343) -- (723,347);
\draw[mapOutline] (780,333) -- (777,333) -- (775,335) -- (767,335);
\draw[mapOutline] (852,329) -- (832,332);
\draw[mapOutline] (884,329) -- (894,330);
\draw[mapOutline] (1163,410) -- (1164,411) -- (1164,414) -- (1165,414) -- (1164,418) -- (1166,418) -- (1167,423);
\draw[mapOutline] (1710,373) -- (1711,375) -- (1714,377) -- (1714,380) -- (1715,380) -- (1716,382) -- (1716,386) -- (1719,387);
\draw[mapOutline] (1709,357) -- (1713,360) -- (1713,362) -- (1714,362) -- (1714,363) -- (1716,364) -- (1720,369);
\draw[mapOutline] (1712,338) -- (1713,338) -- (1714,340) -- (1717,340) -- (1722,342);
\draw[mapOutline] (265,467) -- (267,469) -- (269,469) -- (269,471) -- (275,474) -- (276,476);
\draw[mapOutline] (277,478) -- (279,480);
\draw[mapOutline] (343,535) -- (344,535) -- (344,537) -- (346,538) -- (350,538) -- (351,539) -- (352,538) -- (353,539) -- (355,539) -- (356,541) -- (358,542);
\draw[mapOutline] (355,531) -- (357,532) -- (358,534);
\draw[mapOutline] (357,537) -- (356,538) -- (356,540) -- (358,542);
\draw[mapOutline] (437,526) -- (437,528) -- (438,529) -- (437,531) -- (438,533) -- (438,539);
\draw[mapOutline] (869,427) -- (869,430) -- (866,435) -- (866,439) -- (868,443) -- (868,448) -- (867,450);
\draw[mapOutline] (896,500) -- (900,501) -- (903,503) -- (905,502) -- (906,504) -- (907,504);
\draw[mapOutline] (993,433) -- (993,436) -- (994,438) -- (996,438) -- (997,442) -- (1001,449);
\draw[mapOutline] (1064,500) -- (1074,502);
\draw[mapOutline] (1109,544) -- (1108,553) -- (1107,554) -- (1107,556) -- (1105,558) -- (1106,562);
\draw[mapOutline] (1212,485) -- (1223,491);
\draw[mapOutline] (1226,476) -- (1227,476) -- (1228,478) -- (1235,477) -- (1237,479) -- (1237,481) -- (1238,481);
\draw[mapOutline] (1782,449) -- (1783,451) -- (1787,451) -- (1793,455) -- (1793,458) -- (1794,459) -- (1792,462) -- (1793,463);
\draw[mapOutline] (1810,467) -- (1811,467) -- (1811,469) -- (1814,471) -- (1815,470) -- (1820,470) -- (1826,474) -- (1830,474) -- (1831,476);
\draw[mapOutline] (1911,525) -- (1913,529) -- (1912,532) -- (1914,533) -- (1914,535) -- (1912,538) -- (1917,538) -- (1917,539) -- (1919,539) -- (1920,540) -- (1920,544) -- (1917,548);
\draw[mapOutline] (1919,550) -- (1920,552) -- (1922,552) -- (1924,555) -- (1927,555) -- (1927,556);
\draw[mapOutline] (1997,544) -- (1996,544) -- (1996,546) -- (1995,546) -- (1996,548) -- (1995,549) -- (1993,549) -- (1992,551) -- (1990,551) -- (1991,556) -- (1988,556);
\draw[mapOutline] (2025,508) -- (2025,511) -- (2021,511) -- (2016,509) -- (2016,512);
\draw[mapOutline] (2014,514) -- (2013,515) -- (2012,514) -- (2010,514) -- (2010,519);
\draw[mapOutline] (2033,509) -- (2033,513) -- (2035,514) -- (2035,515) -- (2037,515) -- (2038,517);
\draw[mapOutline] (2038,520) -- (2038,521) -- (2044,526) -- (2047,531);
\draw[mapOutline] (2099,530) -- (2097,533) -- (2097,535) -- (2095,536) -- (2093,539) -- (2091,538) -- (2091,539) -- (2087,539) -- (2085,538);
\draw[mapOutline] (2083,540) -- (2084,542) -- (2083,544);
\draw[mapOutline] (433,602) -- (434,605) -- (436,606) -- (436,608) -- (438,612);
\draw[mapOutline] (493,656) -- (498,656);
\draw[mapOutline] (865,676) -- (863,678) -- (863,680) -- (861,680) -- (858,684);
\draw[mapOutline] (908,571) -- (906,572) -- (906,573) -- (900,576) -- (900,577);
\draw[mapOutline] (1022,574) -- (1023,576) -- (1028,577) -- (1029,579) -- (1034,583);
\draw[mapOutline] (1991,563) -- (1993,562) -- (1994,564) -- (1996,564) -- (1997,566) -- (1999,565) -- (2002,568) -- (2003,566) -- (2006,567);
\draw[mapOutline] (2343,567) -- (2341,570) -- (2340,576);
\draw[mapOutline] (2340,578) -- (2341,579) -- (2340,581);
\draw[mapOutline] (2346,554) -- (2344,556) -- (2345,558) -- (2345,563) -- (2344,563) -- (2343,568) -- (2342,568) -- (2340,572);
\draw[mapOutline] (448,713) -- (454,724);
\draw[mapOutline] (480,740) -- (484,740) -- (486,738) -- (489,738) -- (493,735) -- (499,735) -- (501,734);
\draw[mapOutline] (828,717) -- (827,718);
\draw[mapOutline] (827,722) -- (825,721) -- (824,723) -- (821,724) -- (820,725) -- (821,727) -- (817,730);
\draw[mapOutline] (1329,813) -- (1336,815) -- (1337,816) -- (1337,818) -- (1340,816) -- (1342,816) -- (1344,817) -- (1344,818) -- (1346,817) -- (1349,813) -- (1352,813) -- (1355,816);
\draw[mapOutline] (1365,819) -- (1367,820) -- (1372,819);
\draw[mapOutline] (1374,820) -- (1374,827) -- (1376,828) -- (1376,830) -- (1377,831);
\draw[mapOutline] (766,918) -- (768,923) -- (771,923);
\draw[mapOutline] (1363,939) -- (1366,940) -- (1369,938) -- (1371,938);
\draw[mapOutline] (859,1207) -- (859,1208) -- (862,1208) -- (862,1210) -- (863,1210) -- (864,1212) -- (867,1212) -- (867,1214) -- (871,1217);
\draw[mapOutline] (228,1350) -- (227,1350) -- (228,1352) -- (226,1355) -- (225,1360);
\draw[mapOutline] (277,1390) -- (283,1390) -- (284,1391) -- (290,1391) -- (292,1392);
\draw[mapOutline] (1004,1401) -- (1006,1399) -- (1008,1399) -- (1009,1398) -- (1016,1398);
\draw[mapOutline] (1018,1398) -- (1019,1399) -- (1021,1399);
\draw[mapOutline] (1069,1393) -- (1068,1393) -- (1067,1395) -- (1057,1396);
\draw[mapOutline] (1098,1328) -- (1097,1332) -- (1099,1336) -- (1098,1339) -- (1099,1340);
\draw[mapOutline] (371,1401) -- (370,1402);
\draw[mapOutline] (368,1404) -- (363,1409) -- (363,1410) -- (361,1411);
\draw[mapOutline] (401,1521) -- (400,1523) -- (399,1523) -- (398,1526) -- (395,1528) -- (396,1529) -- (395,1530);
\draw[mapOutline] (1628,1405) -- (1624,1406) -- (1624,1407) -- (1621,1406) -- (1620,1407) -- (1616,1407) -- (1616,1408);
\draw[mapOutline] (1614,1408) -- (1611,1410) -- (1607,1411) -- (1607,1412);
\draw[mapOutline] (1652,1457) -- (1650,1457) -- (1646,1454);
\draw[mapOutline] (1644,1456) -- (1644,1459) -- (1643,1459);
\draw[mapOutline] (1637,1457) -- (1636,1463) -- (1634,1465) -- (1633,1464);
\draw[mapOutline] (1674,1518) -- (1673,1522) -- (1670,1523) -- (1670,1527) -- (1669,1528) -- (1669,1531) -- (1668,1533) -- (1667,1533) -- (1667,1535);
\draw[mapOutline] (1665,1538) -- (1665,1539) -- (1663,1540);
\draw[mapOutline] (2167,1525) -- (2166,1525) -- (2165,1527) -- (2165,1532) -- (2163,1535) -- (2161,1533) -- (2160,1535) -- (2157,1537);
\draw[mapOutline] (860,1628) -- (860,1629) -- (857,1631) -- (858,1635) -- (857,1635);
\draw[mapOutline] (856,1637) -- (855,1638) -- (855,1645);
\draw[mapOutline] (879,1551) -- (881,1550) -- (888,1550) -- (890,1551);
\draw[mapOutline] (1169,1672) -- (1165,1675) -- (1164,1677) -- (1158,1678) -- (1156,1683);
\draw[mapOutline] (1159,1654) -- (1166,1658);
\draw[mapOutline] (1168,1658) -- (1171,1659);
\draw[mapOutline] (1433,1659) -- (1432,1660) -- (1432,1664);
\draw[mapOutline] (1425,1664) -- (1425,1668) -- (1418,1666) -- (1418,1667) -- (1415,1667) -- (1415,1668) -- (1413,1667) -- (1408,1668);
\draw[mapOutline] (1406,1669) -- (1406,1671) -- (1405,1671);
\draw[mapOutline] (2124,1662) -- (2127,1664) -- (2130,1668) -- (2131,1668) -- (2132,1671) -- (2136,1671) -- (2138,1670);
\draw[mapOutline] (1040,1713) -- (1040,1714) -- (1038,1714) -- (1035,1717) -- (1032,1718) -- (1032,1719);
\draw[mapOutline] (1031,1721) -- (1030,1722) -- (1031,1724) -- (1027,1727);
\draw[mapOutline] (1706,1815) -- (1705,1814) -- (1702,1814) -- (1701,1815) -- (1695,1816);
\draw[mapOutline] (1706,1797) -- (1707,1798) -- (1709,1798) -- (1710,1800) -- (1712,1800) -- (1715,1802);
\draw[mapOutline] (2042,1793) -- (2051,1790) -- (2052,1794);
\draw[mapOutline] (2065,1715) -- (2066,1716) -- (2065,1717) -- (2066,1718) -- (2065,1720) -- (2068,1720) -- (2069,1722) -- (2072,1724) -- (2073,1726) -- (2072,1727) -- (2073,1728);
\draw[mapOutline] (2074,1692) -- (2069,1692) -- (2068,1693) -- (2069,1698) -- (2067,1700);
\draw[mapOutline] (2287,1788) -- (2288,1788) -- (2288,1791) -- (2289,1791) -- (2290,1793) -- (2296,1793) -- (2296,1799) -- (2300,1799) -- (2301,1800) -- (2306,1799);
\draw[mapOutline] (1493,1920) -- (1491,1922) -- (1491,1925);
\draw[mapOutline] (1490,1927) -- (1489,1928) -- (1487,1928);
\draw[mapOutline] (1509,1920) -- (1508,1922) -- (1508,1929);
\draw[mapOutline] (1550,1829) -- (1548,1830) -- (1548,1832) -- (1547,1833) -- (1545,1833) -- (1544,1834);
\draw[mapOutline] (1539,1836) -- (1538,1838);
\draw[mapOutline] (1536,1843) -- (1538,1848) -- (1537,1849) -- (1538,1851);
\draw[mapOutline] (1559,1850) -- (1559,1851) -- (1556,1851);
\draw[mapOutline] (1553,1850) -- (1551,1851);
\draw[mapOutline] (1545,1856) -- (1545,1859) -- (1542,1862) -- (1542,1864) -- (1541,1864);
\draw[mapOutline] (1590,1897) -- (1586,1902) -- (1583,1900) -- (1581,1900);
\draw[mapOutline] (1579,1899) -- (1575,1896) -- (1573,1900) -- (1568,1901) -- (1566,1904) -- (1563,1899);
\draw[mapOutline] (1623,1886) -- (1625,1886) -- (1626,1888) -- (1627,1888) -- (1627,1890) -- (1630,1891) -- (1630,1892) -- (1632,1893);
\draw[mapOutline] (1685,1930) -- (1685,1931) -- (1681,1934) -- (1685,1939);
\draw[mapOutline] (1684,1941) -- (1683,1947) -- (1685,1949) -- (1684,1950) -- (1682,1950) -- (1682,1952) -- (1681,1953) -- (1678,1953) -- (1681,1956);
\draw[mapOutline] (1892,1909) -- (1893,1910) -- (1893,1913) -- (1896,1911) -- (1901,1913);
\draw[mapOutline] (2286,1901) -- (2288,1902);
\draw[mapOutline] (2291,1902) -- (2298,1903) -- (2298,1905) -- (2301,1905);
\draw[mapOutline] (2307,1936) -- (2307,1937) -- (2305,1938) -- (2302,1942) -- (2300,1943) -- (2297,1942) -- (2296,1943) -- (2296,1947) -- (2294,1947) -- (2293,1950) -- (2292,1950) -- (2292,1952) -- (2286,1952) -- (2284,1954);
\draw[mapOutline] (1700,2071) -- (1699,2071) -- (1699,2073) -- (1695,2079) -- (1693,2080);
\draw[mapOutline] (1689,2080) -- (1689,2083) -- (1687,2083) -- (1683,2085);
\draw[mapOutline] (2106,2018) -- (2109,2024) -- (2108,2026);
\draw[mapOutline] (2105,2030) -- (2102,2031) -- (2100,2035) -- (2101,2036);
\draw[mapOutline] (2101,2039) -- (2096,2039) -- (2096,2040) -- (2093,2041) -- (2094,2044);
\draw[mapOutline] (2385,2072) -- (2383,2075) -- (2381,2074) -- (2381,2075) -- (2379,2074) -- (2377,2074) -- (2376,2075) -- (2372,2074) -- (2371,2077) -- (2368,2080) -- (2364,2082);
\draw[mapOutline] (2362,2084) -- (2357,2084);
\draw[mapOutline] (2302,2210) -- (2307,2214) -- (2307,2217) -- (2308,2217) -- (2309,2219);
\draw[mapOutline] (2490,2123) -- (2490,2124) -- (2488,2124) -- (2488,2125) -- (2486,2124);
\draw[mapOutline] (2482,2122) -- (2479,2128) -- (2477,2129) -- (2477,2130);
\draw[mapOutline] (2476,2132) -- (2475,2134) -- (2474,2134) -- (2474,2136) -- (2472,2139) -- (2468,2138);
\draw[mapOutline] (2605,2158) -- (2604,2159) -- (2599,2160) -- (2597,2163);
\draw[mapOutline] (2594,2163) -- (2591,2167) -- (2591,2173) -- (2584,2171);
\draw[mapOutline] (1579,2257) -- (1576,2257) -- (1576,2259) -- (1574,2260) -- (1574,2262) -- (1571,2266) -- (1569,2265) -- (1563,2268);
\draw[mapOutline] (1852,2367) -- (1850,2370) -- (1848,2368) -- (1846,2370);
\draw[mapOutline] (1923,2396) -- (1922,2399);
\draw[mapOutline] (1919,2400) -- (1916,2404);
\draw[mapOutline] (1915,2406) -- (1915,2408) -- (1918,2412) -- (1917,2414) -- (1919,2417);
\draw[mapOutline] (1913,2418) -- (1914,2425);
\draw[mapOutline] (496,403) -- (497,403) -- (500,399) -- (502,399) -- (502,398) -- (507,394) -- (508,389) -- (509,389) -- (513,384) -- (519,384) -- (522,382) -- (522,380) -- (524,380) -- (523,379) -- (525,380) -- (524,378) -- (526,376) -- (526,366) -- (530,363) -- (530,362) -- (536,361) -- (542,353) -- (552,352) -- (552,351) -- (554,351) -- (554,350) -- (556,350) -- (556,349) -- (559,350) -- (558,348) -- (559,350) -- (559,348) -- (562,347);
\draw[mapOutline] (517,450) -- (520,449) -- (521,443) -- (524,441) -- (525,437) -- (527,436) -- (527,437) -- (528,437) -- (528,445) -- (532,447) -- (529,448) -- (530,449) -- (529,450) -- (528,456);
\draw[mapOutline] (514,475) -- (513,474) -- (513,475) -- (510,474) -- (512,477) -- (511,478);
\draw[mapOutline] (511,475) -- (511,478) -- (513,478);
\draw[mapOutline] (481,493) -- (485,493) -- (491,498);
\draw[mapOutline] (1052,513) -- (1055,517) -- (1054,519) -- (1058,523) -- (1062,517) -- (1063,522) -- (1065,519) -- (1063,522) -- (1067,526) -- (1071,524) -- (1071,525);
\draw[mapOutline] (1074,528) -- (1074,525) -- (1076,525) -- (1077,527) -- (1078,526);
\draw[mapOutline] (1072,521) -- (1064,519) -- (1063,518) -- (1063,515) -- (1061,516) -- (1061,515) -- (1058,515) -- (1055,513) -- (1052,514);
\draw[mapOutline] (1064,514) -- (1064,513) -- (1067,513) -- (1068,512) -- (1068,513) -- (1069,512) -- (1069,508);
\draw[mapOutline] (267,1233) -- (272,1234) -- (273,1231) -- (274,1231) -- (273,1227);
\draw[mapOutline] (278,1223) -- (280,1222) -- (285,1222) -- (285,1220) -- (288,1220) -- (289,1223) -- (292,1223) -- (293,1221) -- (296,1221) -- (296,1224) -- (297,1224) -- (298,1222) -- (302,1222) -- (302,1221) -- (305,1223) -- (309,1223) -- (309,1224) -- (310,1224);
\draw[mapOutline] (311,1228) -- (313,1228) -- (313,1229) -- (319,1228) -- (320,1227) -- (321,1228) -- (321,1227) -- (324,1226) -- (330,1227) -- (331,1224);
\draw[mapOutline] (336,1222) -- (336,1221) -- (339,1222);
\draw[mapOutline] (348,1229) -- (351,1230);
\draw[mapOutline] (355,1230) -- (355,1231) -- (358,1232);
\draw[mapOutline] (363,1237) -- (375,1235);
\draw[mapOutline] (377,1235) -- (377,1234) -- (381,1235) -- (381,1233) -- (382,1235) -- (382,1233) -- (384,1235) -- (389,1235);
\draw[mapOutline] (396,1236) -- (402,1238);
\draw[mapOutline] (406,1239) -- (412,1239) -- (414,1240);
\draw[mapOutline] (415,1244) -- (419,1241);
\draw[mapOutline] (431,1243) -- (434,1243) -- (434,1242) -- (437,1244) -- (437,1243) -- (438,1243);
\draw[mapOutline] (1272,948) -- (1274,949) -- (1276,948) -- (1275,949) -- (1278,951) -- (1279,950) -- (1278,947) -- (1281,949) -- (1282,948) -- (1280,946) -- (1281,946) -- (1281,943) -- (1283,940) -- (1288,941) -- (1295,938) -- (1297,939) -- (1297,938) -- (1299,938) -- (1299,937) -- (1304,937) -- (1306,938) -- (1305,938) -- (1307,939) -- (1307,941) -- (1309,943);
\draw[mapOutline] (1326,952) -- (1327,949) -- (1326,952) -- (1329,950) -- (1328,954) -- (1329,954) -- (1329,955);
\draw[mapOutline] (1746,421) -- (1745,422) -- (1746,428) -- (1751,429) -- (1751,426) -- (1753,428) -- (1751,426) -- (1752,427) -- (1752,426) -- (1749,425) -- (1748,422);
\draw[mapOutline] (1754,431) -- (1753,431) -- (1753,434) -- (1752,434) -- (1755,435) -- (1754,436) -- (1753,435) -- (1754,436) -- (1756,436) -- (1756,435) -- (1756,437) -- (1760,437) -- (1757,439) -- (1758,440) -- (1757,439) -- (1759,438) -- (1760,436) -- (1759,433) -- (1755,431);
\draw[mapOutline] (1766,433) -- (1766,434) -- (1763,433) -- (1764,441) -- (1767,442) -- (1767,443) -- (1771,442) -- (1770,444) -- (1771,444) -- (1771,445) -- (1770,444) -- (1771,442);
\draw[mapOutline] (1770,439) -- (1769,436) -- (1765,433);
\draw[mapOutline] (1753,447) -- (1754,447) -- (1758,452) -- (1757,455) -- (1762,455) -- (1763,456) -- (1761,459) -- (1762,459) -- (1766,464) -- (1771,465) -- (1771,466) -- (1773,466) -- (1776,468) -- (1780,468) -- (1783,470) -- (1785,470);
\draw[mapOutline] (2115,561) -- (2116,564) -- (2117,564) -- (2117,566) -- (2115,567) -- (2116,569) -- (2114,569) -- (2116,569) -- (2122,573);
\draw[mapOutline] (2122,573) -- (2128,572) -- (2131,575) -- (2131,576) -- (2129,577) -- (2131,576) -- (2131,579) -- (2132,579) -- (2132,580) -- (2136,579) -- (2137,577);
\draw[mapOutline] (2137,577) -- (2139,574) -- (2141,574) -- (2142,576) -- (2143,576);
\draw[mapOutline] (2148,576) -- (2145,580) -- (2149,580);
\draw[mapOutline] (2156,596) -- (2156,598) -- (2158,599) -- (2163,597);
\draw[mapOutline] (2165,597) -- (2167,597) -- (2168,595) -- (2169,598) -- (2170,598) -- (2171,596) -- (2172,596);
\draw[mapOutline] (2105,627) -- (2108,624);
\draw[mapOutline] (2110,622) -- (2110,621) -- (2108,620) -- (2111,620) -- (2111,619);
\draw[mapOutline] (2124,604) -- (2125,604) -- (2128,600);
\draw[mapOutline] (2167,593) -- (2168,587) -- (2167,586) -- (2166,576);
\draw[mapOutline] (2162,577) -- (2163,575) -- (2160,573);
\draw[mapOutline] (2178,591) -- (2178,590) -- (2181,588);
\draw[mapOutline] (2627,848) -- (2632,850) -- (2633,852) -- (2636,853);
\draw[mapOutline] (2638,858) -- (2637,863) -- (2641,865) -- (2638,866) -- (2639,878) -- (2640,881) -- (2641,881) -- (2643,884) -- (2645,882);
\draw[mapOutline] (2663,896) -- (2670,906) -- (2673,908) -- (2679,907) -- (2679,908);
\draw[mapOutline] (2708,895) -- (2711,896) -- (2711,892) -- (2715,893);
\draw[mapOutline] (2720,893) -- (2725,893);
\draw[mapOutline] (2727,893) -- (2735,893);
\draw[mapOutline] (2737,893) -- (2740,894) -- (2740,895) -- (2742,895) -- (2742,896) -- (2744,897);
\draw[mapOutline] (2774,952) -- (2774,953) -- (2772,953) -- (2772,954) -- (2771,953) -- (2770,954) -- (2765,954) -- (2764,955) -- (2764,958);
\draw[mapOutline] (2764,966) -- (2763,966) -- (2762,973);
\draw[mapOutline] (2762,975) -- (2761,981) -- (2760,981) -- (2760,982);
\draw[mapOutline] (2759,984) -- (2759,985) -- (2758,985) -- (2760,986) -- (2758,989) -- (2757,995);
\draw[mapOutline] (2777,953) -- (2775,958) -- (2777,960) -- (2779,960) -- (2780,958) -- (2783,959) -- (2782,956) -- (2783,959) -- (2784,957) -- (2783,959) -- (2784,962) -- (2786,962);
\draw[mapOutline] (2378,1105) -- (2381,1105);
\draw[mapOutline] (2383,1105) -- (2384,1105) -- (2383,1102);
\draw[mapOutline] (2380,1126) -- (2380,1127) -- (2379,1127) -- (2380,1128);
\draw[mapOutline] (2403,1098) -- (2405,1100) -- (2411,1099) -- (2410,1098) -- (2410,1094);
\draw[mapOutline] (2430,1094) -- (2428,1096) -- (2430,1095) -- (2430,1096);
\draw[mapOutline] (570,1555) -- (570,1554) -- (572,1554) -- (573,1551) -- (574,1551) -- (575,1549);
\draw[mapOutline] (578,1543) -- (579,1541) -- (582,1540) -- (582,1539) -- (584,1539) -- (584,1538) -- (589,1538) -- (589,1536) -- (590,1536) -- (592,1539) -- (591,1539) -- (593,1540) -- (595,1537) -- (594,1540) -- (595,1537) -- (596,1537);
\draw[mapOutline] (600,1533) -- (599,1531) -- (601,1532) -- (602,1530) -- (603,1532) -- (609,1530);
\draw[mapOutline] (626,1511) -- (632,1510) -- (633,1509);
\draw[mapOutline] (231,1365) -- (232,1367) -- (233,1366) -- (233,1369) -- (235,1369) -- (237,1371);
\draw[mapOutline] (241,1373) -- (248,1378);
\draw[mapOutline] (249,1379) -- (250,1379) -- (251,1381) -- (254,1382);
\draw[mapOutline] (258,1385) -- (258,1386) -- (263,1388) -- (264,1387) -- (264,1388) -- (265,1388);
\draw[mapOutline] (437,1583) -- (438,1577) -- (439,1577) -- (444,1571) -- (444,1568) -- (443,1568) -- (444,1568) -- (444,1566) -- (445,1565) -- (449,1565);
\draw[mapOutline] (460,1566) -- (463,1565);
\draw[mapOutline] (475,1567) -- (481,1567) -- (483,1568);
\draw[mapOutline] (2143,1748) -- (2144,1748) -- (2144,1751) -- (2146,1753) -- (2151,1752) -- (2150,1752) -- (2151,1752) -- (2154,1756) -- (2157,1756) -- (2157,1757) -- (2159,1758) -- (2159,1760) -- (2164,1761) -- (2166,1764) -- (2168,1764) -- (2171,1766) -- (2171,1767) -- (2183,1765) -- (2185,1764) -- (2185,1763) -- (2191,1762) -- (2191,1761) -- (2193,1762) -- (2195,1761) -- (2196,1762) -- (2196,1761) -- (2199,1759);
\draw[mapOutline] (2218,1751) -- (2224,1749);
\draw[mapOutline] (2252,1727) -- (2253,1728) -- (2255,1727) -- (2257,1728) -- (2257,1725) -- (2259,1725);
\draw[mapOutline] (2261,1724) -- (2265,1722) -- (2267,1725) -- (2267,1723) -- (2268,1725) -- (2268,1722) -- (2269,1724);
\draw[mapOutline] (2268,1726) -- (2268,1729);
\draw[mapOutline] (2269,1731) -- (2270,1745) -- (2269,1745) -- (2269,1750) -- (2267,1751) -- (2269,1753) -- (2270,1757) -- (2271,1757);
\draw[mapOutline] (2271,1759) -- (2271,1760) -- (2270,1760) -- (2270,1763) -- (2269,1763) -- (2270,1766) -- (2269,1766) -- (2271,1770);
\draw[mapOutline] (2269,1771) -- (2269,1772) -- (2268,1772) -- (2268,1776);
\draw[mapOutline] (2266,1780) -- (2265,1786);
\draw[mapOutline] (2269,1798) -- (2271,1801) -- (2275,1800);
\draw[mapOutline] (2268,1801) -- (2266,1806) -- (2269,1807) -- (2266,1808) -- (2266,1809);
\draw[mapOutline] (2629,1453) -- (2623,1454) -- (2625,1454) -- (2624,1457) -- (2625,1458) -- (2627,1457) -- (2627,1458) -- (2626,1458) -- (2628,1459) -- (2627,1463) -- (2628,1463);
\draw[mapOutline] (2630,1471) -- (2630,1473) -- (2624,1474) -- (2623,1475) -- (2624,1476) -- (2623,1480);
\draw[mapOutline] (2624,1482) -- (2623,1484) -- (2624,1483) -- (2624,1486) -- (2626,1486) -- (2627,1488);
\draw[mapOutline] (2642,1513) -- (2642,1516);
\draw[mapOutline] (2643,1518) -- (2645,1522);
\draw[mapOutline] (2646,1524) -- (2646,1526) -- (2648,1527) -- (2651,1533);
\draw[mapOutline] (2660,1540) -- (2662,1540) -- (2662,1541) -- (2665,1543) -- (2666,1546) -- (2666,1545) -- (2667,1545);
\draw[mapOutline] (2684,1504) -- (2682,1501) -- (2682,1498) -- (2685,1496) -- (2687,1496) -- (2687,1495) -- (2689,1495) -- (2689,1494) -- (2694,1493) -- (2694,1492) -- (2699,1494);
\draw[mapOutline] (2703,1492) -- (2699,1494) -- (2699,1495) -- (2701,1496);
\draw[mapOutline] (2701,1499) -- (2698,1498) -- (2697,1499) -- (2699,1505) -- (2698,1506) -- (2699,1504) -- (2698,1506) -- (2695,1506) -- (2694,1505) -- (2690,1507) -- (2690,1509) -- (2688,1510) -- (2687,1509) -- (2687,1510) -- (2685,1510) -- (2684,1507) -- (2685,1507) -- (2684,1504) -- (2685,1505) -- (2684,1505);
\draw[mapOutline] (2710,1510) -- (2712,1509) -- (2710,1507) -- (2711,1507) -- (2713,1510) -- (2716,1508);
\draw[mapOutline] (2730,1509) -- (2731,1511) -- (2731,1508);
\draw[mapOutline] (2737,1506) -- (2738,1506) -- (2737,1513);
\draw[mapOutline] (2366,1420) -- (2366,1416);
\draw[mapOutline] (2388,1395) -- (2386,1393) -- (2388,1392) -- (2389,1389) -- (2392,1388) -- (2391,1387) -- (2392,1385) -- (2390,1381);
\draw[mapOutline] (2386,1443) -- (2386,1438) -- (2387,1438) -- (2388,1435) -- (2390,1434) -- (2390,1433);
\draw[mapOutline] (2390,1473) -- (2390,1476);
\draw[mapOutline] (2393,1481) -- (2395,1484) -- (2394,1482) -- (2395,1482);
\draw[mapOutline] (2445,1492) -- (2448,1492) -- (2450,1491) -- (2451,1489) -- (2453,1491);
\draw[mapOutline] (2459,1431) -- (2459,1439);
\draw[mapOutline] (2476,1455) -- (2477,1458) -- (2478,1458) -- (2479,1460) -- (2482,1460);
\draw[mapOutline] (2500,1485) -- (2502,1486) -- (2501,1488);
\draw[mapOutline] (2436,1492) -- (2439,1494) -- (2437,1491) -- (2438,1489) -- (2442,1490) -- (2442,1491) -- (2446,1491);
\draw[mapOutline] (2449,1488) -- (2449,1487) -- (2450,1489);
\draw[mapOutline] (2236,1960) -- (2236,1962) -- (2234,1965) -- (2235,1968) -- (2234,1969);
\draw[mapOutline] (2235,1971) -- (2235,1974) -- (2238,1973) -- (2237,1975) -- (2238,1973) -- (2241,1972) -- (2238,1971) -- (2240,1970) -- (2241,1972);
\draw[mapOutline] (2243,1972) -- (2244,1972) -- (2244,1974) -- (2246,1975) -- (2247,1979);
\draw[mapOutline] (2267,1988) -- (2269,1987) -- (2271,1988);
\draw[mapOutline] (2278,1966) -- (2278,1968) -- (2279,1970) -- (2280,1970);
\draw[mapOutline] (2281,1974) -- (2281,1977);
\draw[mapOutline] (2267,2034) -- (2268,2035) -- (2268,2039) -- (2269,2041) -- (2270,2041) -- (2270,2044) -- (2267,2046);
\draw[mapOutline] (2271,2045) -- (2273,2050);
\draw[mapOutline] (2274,2052) -- (2274,2053) -- (2275,2053) -- (2274,2055) -- (2276,2055) -- (2273,2057) -- (2273,2059);
\draw[mapOutline] (2281,2066) -- (2281,2069);
\draw[mapOutline] (2401,2006) -- (2400,2006) -- (2400,2007) -- (2397,2006) -- (2394,2010);
\draw[mapOutline] (2409,2007) -- (2407,2004) -- (2410,2002);
\draw[mapOutline] (2412,1995) -- (2416,1996);
\draw[mapOutline] (2412,1990) -- (2411,1987) -- (2410,1989) -- (2410,1986) -- (2409,1987) -- (2410,1987) -- (2410,1985);
\draw[mapOutline] (2027,1964) -- (2027,1961) -- (2029,1961);
\draw[mapOutline] (2023,1978) -- (2023,1982) -- (2027,1981) -- (2026,1982) -- (2026,1981) -- (2028,1981) -- (2029,1979) -- (2033,1980) -- (2034,1982) -- (2037,1981);
\draw[mapOutline] (2047,1997) -- (2042,1997) -- (2044,1999) -- (2043,1999) -- (2043,2004) -- (2044,2006) -- (2046,2006) -- (2046,2007) -- (2046,2006) -- (2052,2005);
\draw[mapOutline] (2087,2005) -- (2087,2009) -- (2090,2011);
\draw[mapOutline] (1732,2374) -- (1732,2379) -- (1733,2381);
\draw[mapOutline] (1735,2382) -- (1740,2382);
\draw[mapOutline] (1742,2381) -- (1742,2382) -- (1739,2382);
\draw[mapOutline] (1742,2385) -- (1739,2387) -- (1739,2388) -- (1737,2388) -- (1737,2389) -- (1731,2388) -- (1732,2389) -- (1729,2389) -- (1729,2391) -- (1728,2391) -- (1729,2404) -- (1730,2404) -- (1728,2406) -- (1729,2408) -- (1733,2406);
\draw[mapOutline] (1731,2410) -- (1733,2412) -- (1737,2411) -- (1737,2413) -- (1739,2411) -- (1741,2411) -- (1745,2414);
\draw[mapOutline] (1747,2412) -- (1748,2409) -- (1747,2412) -- (1748,2412) -- (1750,2418) -- (1757,2421) -- (1757,2422) -- (1758,2422) -- (1759,2424) -- (1761,2424) -- (1762,2426) -- (1764,2425) -- (1764,2427);
\draw[mapOutline] (1766,2429) -- (1769,2430);
\draw[mapOutline] (1770,2426) -- (1773,2430) -- (1774,2430) -- (1775,2428);
\draw[mapOutline] (1780,2434) -- (1783,2436) -- (1783,2437) -- (1787,2438) -- (1787,2439) -- (1795,2440);
\draw[mapOutline] (1797,2441) -- (1802,2441) -- (1806,2444);
\draw[mapOutline] (1823,2451) -- (1824,2450) -- (1823,2452) -- (1825,2450);
\draw[mapOutline] (1841,2448) -- (1840,2445) -- (1842,2447);
\draw[mapOutline] (1807,2391) -- (1808,2390) -- (1810,2391) -- (1811,2389) -- (1815,2388) -- (1815,2387) -- (1818,2386) -- (1819,2384);
\draw[mapOutline] (1875,2365) -- (1875,2364) -- (1873,2363) -- (1873,2362) -- (1874,2362);
\draw[mapOutline] (1879,2358) -- (1882,2362) -- (1885,2360) -- (1886,2362) -- (1889,2363) -- (1885,2361);
\draw[mapOutline] (1885,2366) -- (1882,2363);
\draw[mapOutline] (1883,2365) -- (1879,2365) -- (1880,2374) -- (1878,2377) -- (1879,2380) -- (1880,2380) -- (1881,2382);
\draw[mapOutline] (1883,2383) -- (1884,2383) -- (1886,2386);
\draw[mapOutline] (1889,2383) -- (1889,2386) -- (1896,2386) -- (1897,2385) -- (1901,2388);
\draw[mapOutline] (1916,2351) -- (1915,2352) -- (1915,2355) -- (1917,2355);
\draw[mapOutline] (1906,2415) -- (1906,2419);
\draw[mapOutline] (1920,2418) -- (1921,2418) -- (1921,2420) -- (1923,2421) -- (1923,2423) -- (1925,2426) -- (1929,2429) -- (1930,2432) -- (1931,2432) -- (1931,2434) -- (1937,2439) -- (1938,2438);
\draw[mapOutline] (1941,2437) -- (1942,2437) -- (1942,2440) -- (1944,2437) -- (1944,2438) -- (1945,2438);

% The ten colored triangulation systems.
\draw[mapBand,draw=HannoverHaupt] (jever) -- (langwarden) -- (bremerlehe) -- (brillit) -- (zeven) -- (litberg);
\draw[mapBand,draw=HannoverHaupt] (hamburg) -- (hohenhorn) -- (lauenburg);
\draw[mapBand,draw=HannoverHaupt] (lueneburg) -- (nindorf) -- (timpenberg) -- (wulfsode) -- (2075,1057) -- (breithorn) -- (scharnhorst) -- (garssen) -- (lichtenberg);
\draw[mapBand,draw=HannoverHaupt] (brocken) -- (2240,2440);
\draw[mapBand,draw=HannoverHaupt] (hamburg) -- (wilsede) -- (falkenberg) -- (deister-i) -- (hils) -- (hohehagen) -- (1940,2440);
\draw[mapBand,draw=HannoverHaupt] (falkenberg) -- (elmhorst) -- (1738,906) -- (steinberg-n);
\draw[mapBand,draw=HannoverHaupt] (bremen) -- (garlste) -- (varel);
\draw[mapBand,draw=HannoverWest] (varel) -- (westerstede) -- (krapendorf) -- (mordkuhlenberg) -- (twistringen) -- (bremen);
\draw[mapBand,draw=HannoverWest] (steinberg-n) -- (asendorf) -- (knickberg);
\draw[mapBand,draw=HannoverWest] (nonnenstein) -- (doerenberg) -- (bentheim) -- (kirchhesepe) -- (windberg) -- (onstwedde) -- (emden) -- (aurich) -- (dornum);
\draw[mapBand,draw=HannoverOst,line width=.72mm] (borkum) -- (juist) -- (norderney) -- (baltrum) -- (langeoog) -- (spiekeroog) -- (jever);
\draw[mapBand,draw=HannoverOst,line width=.72mm] (dornum) -- (hage) -- (pilsum) -- (borkum);
\draw[mapBand,draw=HannoverBremen] (bremerlehe) -- (holssel) -- (1302,354) -- (altenwalde) -- (silberberg) -- (crempe) -- (hamburg);
\draw[mapBand,draw=HannoverBremen] (litberg) -- (stade) -- (basdahl) -- (bremerlehe);
\draw[mapBand,draw=HannoverLueneburg] (lauenburg) -- (bretze) -- (glienitz) -- (hoe-beck) -- (pugelatz) -- (hankensbuettel) -- (wieren) -- (holxerberg) -- (taetendorf) -- (lueneburg);
\draw[mapBand,draw=HannoverLueneburg] (lichtenberg) -- (wohlenberg) -- (festberg) -- (brocken);
\draw[mapBand,draw=HannoverMittel] (knickberg) -- (osterberg-w) -- (deister-ii) -- (osterwald) -- (hils) -- (koeterberg) -- (wittekindstein) -- (nonnenstein);
\draw[mapBand,draw=HannoverEichs] (tockenburg) -- (wulften) -- (sonnenstein) -- (bazenberg) -- (hauseberg) -- (tockenburg);
\draw[mapBand,draw=HannoverHarz] (fuhrenberg) -- (bocksberg) -- (brocken) -- (schalliethe) -- (sonnenstein) -- (haspelkopf) -- (fuhrenberg);
\draw[mapBand,draw=HannoverOber,line width=.61mm] (fahlenstollen) -- (kleper) -- (gieseberg) -- (steinberg-s) -- (fahlenstollen);
\draw[draw=HannoverOber!70!black,line width=.10pt] (1781.2,2107.7) -- (1790.9,2105.5);
\draw[draw=HannoverOber!70!black,line width=.10pt] (1784.2,2110.3) -- (1794.0,2108.1);
\draw[draw=HannoverOber!70!black,line width=.10pt] (1787.2,2112.9) -- (1797.0,2110.7);
\draw[draw=HannoverOber!70!black,line width=.10pt] (1790.3,2115.5) -- (1800.0,2113.3);
\draw[draw=HannoverOber!70!black,line width=.10pt] (1793.3,2118.1) -- (1803.1,2115.9);
\draw[draw=HannoverOber!70!black,line width=.10pt] (1796.4,2120.7) -- (1806.1,2118.5);
\draw[draw=HannoverOber!70!black,line width=.10pt] (1799.4,2123.3) -- (1809.2,2121.1);
\draw[draw=HannoverOber!70!black,line width=.10pt] (1802.4,2125.9) -- (1812.2,2123.7);
\draw[draw=HannoverOber!70!black,line width=.10pt] (1805.5,2128.5) -- (1815.2,2126.3);
\draw[draw=HannoverOber!70!black,line width=.10pt] (1808.5,2131.1) -- (1818.3,2128.9);
\draw[draw=HannoverOber!70!black,line width=.10pt] (1811.6,2133.7) -- (1821.3,2131.5);
\draw[draw=HannoverOber!70!black,line width=.10pt] (1814.6,2136.3) -- (1824.4,2134.1);
\draw[draw=HannoverOber!70!black,line width=.10pt] (1817.6,2138.9) -- (1827.4,2136.7);
\draw[draw=HannoverOber!70!black,line width=.10pt] (1820.7,2141.5) -- (1830.4,2139.3);
\draw[draw=HannoverOber!70!black,line width=.10pt] (1823.7,2144.1) -- (1833.5,2141.9);
\draw[draw=HannoverOber!70!black,line width=.10pt] (1826.8,2146.7) -- (1836.5,2144.5);
\draw[draw=HannoverOber!70!black,line width=.10pt] (1829.8,2149.3) -- (1839.6,2147.1);
\draw[draw=HannoverOber!70!black,line width=.10pt] (1832.8,2151.9) -- (1842.6,2149.7);
\draw[draw=HannoverOber!70!black,line width=.10pt] (1835.9,2154.5) -- (1845.6,2152.3);
\draw[draw=HannoverOber!70!black,line width=.10pt] (1838.9,2157.1) -- (1848.7,2154.9);
\draw[draw=HannoverOber!70!black,line width=.10pt] (1842.0,2159.7) -- (1851.7,2157.5);
\draw[draw=HannoverOber!70!black,line width=.10pt] (1845.0,2162.3) -- (1854.8,2160.1);
\draw[draw=HannoverOber!70!black,line width=.10pt] (1848.0,2164.9) -- (1857.8,2162.7);
\draw[draw=HannoverOber!70!black,line width=.10pt] (1851.1,2167.5) -- (1860.8,2165.3);
\draw[draw=HannoverOber!70!black,line width=.10pt] (1854.1,2170.1) -- (1863.9,2167.9);
\draw[draw=HannoverOber!70!black,line width=.10pt] (1857.2,2172.7) -- (1866.9,2170.5);
\draw[draw=HannoverOber!70!black,line width=.10pt] (1860.2,2175.3) -- (1869.9,2173.1);
\draw[draw=HannoverOber!70!black,line width=.10pt] (1863.2,2177.9) -- (1873.0,2175.7);
\draw[draw=HannoverOber!70!black,line width=.10pt] (1866.3,2180.5) -- (1876.0,2178.3);
\draw[draw=HannoverOber!70!black,line width=.10pt] (1869.3,2183.1) -- (1879.1,2180.9);
\draw[draw=HannoverOber!70!black,line width=.10pt] (1872.3,2185.7) -- (1882.1,2183.5);
\draw[draw=HannoverOber!70!black,line width=.10pt] (1875.4,2188.3) -- (1885.1,2186.1);
\draw[draw=HannoverOber!70!black,line width=.10pt] (1878.4,2190.9) -- (1888.2,2188.7);
\draw[draw=HannoverOber!70!black,line width=.10pt] (1881.5,2193.5) -- (1891.2,2191.3);
\draw[draw=HannoverOber!70!black,line width=.10pt] (1884.5,2196.1) -- (1894.3,2193.9);
\draw[draw=HannoverOber!70!black,line width=.10pt] (1887.5,2198.7) -- (1897.3,2196.5);
\draw[draw=HannoverOber!70!black,line width=.10pt] (1890.6,2201.3) -- (1900.3,2199.1);
\draw[draw=HannoverOber!70!black,line width=.10pt] (1893.6,2203.9) -- (1903.4,2201.7);
\draw[draw=HannoverOber!70!black,line width=.10pt] (1896.7,2206.5) -- (1906.4,2204.3);
\draw[draw=HannoverOber!70!black,line width=.10pt] (1899.7,2209.1) -- (1909.5,2206.9);
\draw[draw=HannoverOber!70!black,line width=.10pt] (1902.7,2211.7) -- (1912.5,2209.5);
\draw[draw=HannoverOber!70!black,line width=.10pt] (1905.8,2214.3) -- (1915.5,2212.1);
\draw[draw=HannoverOber!70!black,line width=.10pt] (1908.8,2216.9) -- (1918.6,2214.7);
\draw[draw=HannoverOber!70!black,line width=.10pt] (1911.9,2219.5) -- (1921.6,2217.3);
\draw[draw=HannoverOber!70!black,line width=.10pt] (1914.9,2222.1) -- (1924.7,2219.9);
\draw[draw=HannoverOber!70!black,line width=.10pt] (1917.9,2224.7) -- (1927.7,2222.5);
\draw[draw=HannoverOber!70!black,line width=.10pt] (1921.0,2227.3) -- (1930.7,2225.1);
\draw[draw=HannoverOber!70!black,line width=.10pt] (1924.0,2229.9) -- (1933.8,2227.7);
\draw[draw=HannoverOber!70!black,line width=.10pt] (1927.1,2232.5) -- (1936.8,2230.3);
\draw[draw=HannoverOber!70!black,line width=.10pt] (1930.1,2235.1) -- (1939.9,2232.9);
\draw[draw=HannoverOber!70!black,line width=.10pt] (1933.1,2237.7) -- (1942.9,2235.5);
\draw[draw=HannoverOber!70!black,line width=.10pt] (1936.2,2240.3) -- (1945.9,2238.1);
\draw[draw=HannoverOber!70!black,line width=.10pt] (1939.2,2242.9) -- (1949.0,2240.7);
\draw[draw=HannoverOber!70!black,line width=.10pt] (1945.1,2244.8) -- (1948.8,2254.1);
\draw[draw=HannoverOber!70!black,line width=.10pt] (1943.0,2248.2) -- (1946.7,2257.5);
\draw[draw=HannoverOber!70!black,line width=.10pt] (1940.9,2251.6) -- (1944.7,2260.9);
\draw[draw=HannoverOber!70!black,line width=.10pt] (1938.9,2255.1) -- (1942.6,2264.3);
\draw[draw=HannoverOber!70!black,line width=.10pt] (1936.8,2258.5) -- (1940.5,2267.8);
\draw[draw=HannoverOber!70!black,line width=.10pt] (1934.7,2261.9) -- (1938.5,2271.2);
\draw[draw=HannoverOber!70!black,line width=.10pt] (1932.6,2265.3) -- (1936.4,2274.6);
\draw[draw=HannoverOber!70!black,line width=.10pt] (1930.6,2268.8) -- (1934.3,2278.0);
\draw[draw=HannoverOber!70!black,line width=.10pt] (1928.5,2272.2) -- (1932.3,2281.4);
\draw[draw=HannoverOber!70!black,line width=.10pt] (1926.4,2275.6) -- (1930.2,2284.9);
\draw[draw=HannoverOber!70!black,line width=.10pt] (1924.4,2279.0) -- (1928.1,2288.3);
\draw[draw=HannoverOber!70!black,line width=.10pt] (1922.3,2282.4) -- (1926.1,2291.7);
\draw[draw=HannoverOber!70!black,line width=.10pt] (1920.2,2285.9) -- (1924.0,2295.1);
\draw[draw=HannoverOber!70!black,line width=.10pt] (1918.2,2289.3) -- (1921.9,2298.6);
\draw[draw=HannoverOber!70!black,line width=.10pt] (1916.1,2292.7) -- (1919.8,2302.0);
\draw[draw=HannoverOber!70!black,line width=.10pt] (1914.0,2296.1) -- (1917.8,2305.4);
\draw[draw=HannoverOber!70!black,line width=.10pt] (1912.0,2299.6) -- (1915.7,2308.8);
\draw[draw=HannoverOber!70!black,line width=.10pt] (1909.9,2303.0) -- (1913.6,2312.3);
\draw[draw=HannoverOber!70!black,line width=.10pt] (1907.8,2306.4) -- (1911.6,2315.7);
\draw[draw=HannoverOber!70!black,line width=.10pt] (1905.8,2309.8) -- (1909.5,2319.1);
\draw[draw=HannoverOber!70!black,line width=.10pt] (1903.7,2313.3) -- (1907.4,2322.5);
\draw[draw=HannoverOber!70!black,line width=.10pt] (1901.6,2316.7) -- (1905.4,2326.0);
\draw[draw=HannoverOber!70!black,line width=.10pt] (1899.6,2320.1) -- (1903.3,2329.4);
\draw[draw=HannoverOber!70!black,line width=.10pt] (1897.5,2323.5) -- (1901.2,2332.8);
\draw[draw=HannoverOber!70!black,line width=.10pt] (1895.4,2327.0) -- (1899.2,2336.2);
\draw[draw=HannoverOber!70!black,line width=.10pt] (1893.3,2330.4) -- (1897.1,2339.7);
\draw[draw=HannoverOber!70!black,line width=.10pt] (1891.3,2333.8) -- (1895.0,2343.1);
\draw[draw=HannoverOber!70!black,line width=.10pt] (1888.5,2338.7) -- (1886.1,2348.4);
\draw[draw=HannoverOber!70!black,line width=.10pt] (1884.8,2340.3) -- (1882.4,2350.0);
\draw[draw=HannoverOber!70!black,line width=.10pt] (1881.2,2341.8) -- (1878.8,2351.6);
\draw[draw=HannoverOber!70!black,line width=.10pt] (1877.5,2343.4) -- (1875.1,2353.1);
\draw[draw=HannoverOber!70!black,line width=.10pt] (1873.8,2345.0) -- (1871.4,2354.7);
\draw[draw=HannoverOber!70!black,line width=.10pt] (1870.1,2346.5) -- (1867.7,2356.3);
\draw[draw=HannoverOber!70!black,line width=.10pt] (1866.4,2348.1) -- (1864.0,2357.8);
\draw[draw=HannoverOber!70!black,line width=.10pt] (1862.7,2349.7) -- (1860.4,2359.4);
\draw[draw=HannoverOber!70!black,line width=.10pt] (1859.1,2351.2) -- (1856.7,2361.0);
\draw[draw=HannoverOber!70!black,line width=.10pt] (1855.4,2352.8) -- (1853.0,2362.5);
\draw[draw=HannoverOber!70!black,line width=.10pt] (1851.7,2354.4) -- (1849.3,2364.1);
\draw[draw=HannoverOber!70!black,line width=.10pt] (1848.0,2355.9) -- (1845.6,2365.6);
\draw[draw=HannoverOber!70!black,line width=.10pt] (1844.3,2357.5) -- (1842.0,2367.2);
\draw[draw=HannoverOber!70!black,line width=.10pt] (1840.7,2359.1) -- (1838.3,2368.8);
\draw[draw=HannoverOber!70!black,line width=.10pt] (1837.0,2360.6) -- (1834.6,2370.3);
\draw[draw=HannoverOber!70!black,line width=.10pt] (1833.3,2362.2) -- (1830.9,2371.9);
\draw[draw=HannoverOber!70!black,line width=.10pt] (1829.6,2363.8) -- (1827.2,2373.5);
\draw[draw=HannoverOber!70!black,line width=.10pt] (1825.9,2365.3) -- (1823.6,2375.0);
\draw[draw=HannoverOber!70!black,line width=.10pt] (1822.3,2366.9) -- (1819.9,2376.6);
\draw[draw=HannoverOber!70!black,line width=.10pt] (1818.6,2368.5) -- (1816.2,2378.2);
\draw[draw=HannoverOber!70!black,line width=.10pt] (1814.9,2370.0) -- (1812.5,2379.7);
\draw[draw=HannoverOber!70!black,line width=.10pt] (1811.2,2371.6) -- (1808.8,2381.3);
\draw[draw=HannoverOber!70!black,line width=.10pt] (1807.5,2373.2) -- (1805.2,2382.9);
\draw[draw=HannoverOber!70!black,line width=.10pt] (1803.9,2374.7) -- (1801.5,2384.4);
\draw[draw=HannoverOber!70!black,line width=.10pt] (1800.9,2380.8) -- (1792.7,2375.2);
\draw[draw=HannoverOber!70!black,line width=.10pt] (1800.7,2376.8) -- (1792.5,2371.2);
\draw[draw=HannoverOber!70!black,line width=.10pt] (1800.5,2372.8) -- (1792.3,2367.2);
\draw[draw=HannoverOber!70!black,line width=.10pt] (1800.3,2368.8) -- (1792.0,2363.2);
\draw[draw=HannoverOber!70!black,line width=.10pt] (1800.1,2364.8) -- (1791.8,2359.2);
\draw[draw=HannoverOber!70!black,line width=.10pt] (1799.9,2360.8) -- (1791.6,2355.2);
\draw[draw=HannoverOber!70!black,line width=.10pt] (1799.7,2356.8) -- (1791.4,2351.2);
\draw[draw=HannoverOber!70!black,line width=.10pt] (1799.5,2352.8) -- (1791.2,2347.2);
\draw[draw=HannoverOber!70!black,line width=.10pt] (1799.3,2348.8) -- (1791.0,2343.3);
\draw[draw=HannoverOber!70!black,line width=.10pt] (1799.1,2344.8) -- (1790.8,2339.3);
\draw[draw=HannoverOber!70!black,line width=.10pt] (1798.9,2340.9) -- (1790.6,2335.3);
\draw[draw=HannoverOber!70!black,line width=.10pt] (1798.7,2336.9) -- (1790.4,2331.3);
\draw[draw=HannoverOber!70!black,line width=.10pt] (1798.5,2332.9) -- (1790.2,2327.3);
\draw[draw=HannoverOber!70!black,line width=.10pt] (1798.3,2328.9) -- (1790.0,2323.3);
\draw[draw=HannoverOber!70!black,line width=.10pt] (1798.1,2324.9) -- (1789.8,2319.3);
\draw[draw=HannoverOber!70!black,line width=.10pt] (1797.9,2320.9) -- (1789.6,2315.3);
\draw[draw=HannoverOber!70!black,line width=.10pt] (1797.7,2316.9) -- (1789.4,2311.3);
\draw[draw=HannoverOber!70!black,line width=.10pt] (1797.5,2312.9) -- (1789.2,2307.3);
\draw[draw=HannoverOber!70!black,line width=.10pt] (1797.3,2308.9) -- (1789.0,2303.3);
\draw[draw=HannoverOber!70!black,line width=.10pt] (1797.1,2304.9) -- (1788.8,2299.3);
\draw[draw=HannoverOber!70!black,line width=.10pt] (1796.9,2300.9) -- (1788.6,2295.3);
\draw[draw=HannoverOber!70!black,line width=.10pt] (1796.7,2296.9) -- (1788.4,2291.3);
\draw[draw=HannoverOber!70!black,line width=.10pt] (1796.5,2292.9) -- (1788.2,2287.3);
\draw[draw=HannoverOber!70!black,line width=.10pt] (1796.3,2288.9) -- (1788.0,2283.3);
\draw[draw=HannoverOber!70!black,line width=.10pt] (1796.1,2284.9) -- (1787.8,2279.3);
\draw[draw=HannoverOber!70!black,line width=.10pt] (1795.9,2280.9) -- (1787.6,2275.3);
\draw[draw=HannoverOber!70!black,line width=.10pt] (1795.7,2276.9) -- (1787.4,2271.3);
\draw[draw=HannoverOber!70!black,line width=.10pt] (1795.5,2272.9) -- (1787.2,2267.3);
\draw[draw=HannoverOber!70!black,line width=.10pt] (1795.3,2268.9) -- (1787.0,2263.4);
\draw[draw=HannoverOber!70!black,line width=.10pt] (1795.1,2264.9) -- (1786.8,2259.4);
\draw[draw=HannoverOber!70!black,line width=.10pt] (1794.9,2261.0) -- (1786.6,2255.4);
\draw[draw=HannoverOber!70!black,line width=.10pt] (1794.7,2257.0) -- (1786.4,2251.4);
\draw[draw=HannoverOber!70!black,line width=.10pt] (1794.5,2253.0) -- (1786.2,2247.4);
\draw[draw=HannoverOber!70!black,line width=.10pt] (1794.3,2249.0) -- (1786.0,2243.4);
\draw[draw=HannoverOber!70!black,line width=.10pt] (1794.1,2245.0) -- (1785.8,2239.4);
\draw[draw=HannoverOber!70!black,line width=.10pt] (1793.9,2241.0) -- (1785.6,2235.4);
\draw[draw=HannoverOber!70!black,line width=.10pt] (1793.7,2237.0) -- (1785.4,2231.4);
\draw[draw=HannoverOber!70!black,line width=.10pt] (1793.5,2233.0) -- (1785.2,2227.4);
\draw[draw=HannoverOber!70!black,line width=.10pt] (1793.3,2229.0) -- (1785.0,2223.4);
\draw[draw=HannoverOber!70!black,line width=.10pt] (1793.1,2225.0) -- (1784.8,2219.4);
\draw[draw=HannoverOber!70!black,line width=.10pt] (1792.9,2221.0) -- (1784.6,2215.4);
\draw[draw=HannoverOber!70!black,line width=.10pt] (1792.7,2217.0) -- (1784.4,2211.4);
\draw[draw=HannoverOber!70!black,line width=.10pt] (1792.5,2213.0) -- (1784.2,2207.4);
\draw[draw=HannoverOber!70!black,line width=.10pt] (1792.3,2209.0) -- (1784.0,2203.4);
\draw[draw=HannoverOber!70!black,line width=.10pt] (1792.1,2205.0) -- (1783.8,2199.4);
\draw[draw=HannoverOber!70!black,line width=.10pt] (1791.9,2201.0) -- (1783.6,2195.4);
\draw[draw=HannoverOber!70!black,line width=.10pt] (1791.7,2197.0) -- (1783.4,2191.4);
\draw[draw=HannoverOber!70!black,line width=.10pt] (1791.5,2193.0) -- (1783.2,2187.4);
\draw[draw=HannoverOber!70!black,line width=.10pt] (1791.3,2189.0) -- (1783.0,2183.5);
\draw[draw=HannoverOber!70!black,line width=.10pt] (1791.1,2185.0) -- (1782.8,2179.5);
\draw[draw=HannoverOber!70!black,line width=.10pt] (1790.9,2181.1) -- (1782.6,2175.5);
\draw[draw=HannoverOber!70!black,line width=.10pt] (1790.7,2177.1) -- (1782.4,2171.5);
\draw[draw=HannoverOber!70!black,line width=.10pt] (1790.5,2173.1) -- (1782.2,2167.5);
\draw[draw=HannoverOber!70!black,line width=.10pt] (1790.3,2169.1) -- (1782.0,2163.5);
\draw[draw=HannoverOber!70!black,line width=.10pt] (1790.1,2165.1) -- (1781.8,2159.5);
\draw[draw=HannoverOber!70!black,line width=.10pt] (1789.9,2161.1) -- (1781.6,2155.5);
\draw[draw=HannoverOber!70!black,line width=.10pt] (1789.7,2157.1) -- (1781.4,2151.5);
\draw[draw=HannoverOber!70!black,line width=.10pt] (1789.5,2153.1) -- (1781.2,2147.5);
\draw[draw=HannoverOber!70!black,line width=.10pt] (1789.3,2149.1) -- (1781.0,2143.5);
\draw[draw=HannoverOber!70!black,line width=.10pt] (1789.1,2145.1) -- (1780.8,2139.5);
\draw[draw=HannoverOber!70!black,line width=.10pt] (1788.9,2141.1) -- (1780.6,2135.5);
\draw[draw=HannoverOber!70!black,line width=.10pt] (1788.7,2137.1) -- (1780.4,2131.5);
\draw[draw=HannoverOber!70!black,line width=.10pt] (1788.5,2133.1) -- (1780.2,2127.5);
\draw[draw=HannoverOber!70!black,line width=.10pt] (1788.3,2129.1) -- (1780.0,2123.5);
\draw[draw=HannoverOber!70!black,line width=.10pt] (1788.1,2125.1) -- (1779.8,2119.5);
\draw[draw=HannoverOber!70!black,line width=.10pt] (1787.9,2121.1) -- (1779.6,2115.5);
\draw[draw=HannoverOber!70!black,line width=.10pt] (1787.7,2117.1) -- (1779.4,2111.5);
\draw[draw=HannoverOber!70!black,line width=.10pt] (1787.5,2113.1) -- (1779.2,2107.5);
\draw[mapBand,draw=HannoverOber,line width=.61mm] (hohehagen) -- (weper) -- (meridianzeichen);
\draw[draw=HannoverOber!70!black,line width=.10pt] (1839.1,2292.9) -- (1832.7,2285.3);
\draw[draw=HannoverOber!70!black,line width=.10pt] (1840.0,2289.0) -- (1833.6,2281.4);
\draw[draw=HannoverOber!70!black,line width=.10pt] (1840.9,2285.1) -- (1834.5,2277.5);
\draw[draw=HannoverOber!70!black,line width=.10pt] (1841.8,2281.2) -- (1835.4,2273.6);
\draw[draw=HannoverOber!70!black,line width=.10pt] (1842.7,2277.3) -- (1836.3,2269.7);
\draw[draw=HannoverOber!70!black,line width=.10pt] (1843.6,2273.4) -- (1837.2,2265.8);
\draw[draw=HannoverOber!70!black,line width=.10pt] (1844.5,2269.5) -- (1838.1,2261.9);
\draw[draw=HannoverOber!70!black,line width=.10pt] (1845.4,2265.6) -- (1839.0,2258.0);
\draw[draw=HannoverOber!70!black,line width=.10pt] (1846.3,2261.7) -- (1839.9,2254.1);
\draw[draw=HannoverOber!70!black,line width=.10pt] (1847.2,2257.8) -- (1840.8,2250.2);
\draw[draw=HannoverOber!70!black,line width=.10pt] (1848.1,2253.9) -- (1841.7,2246.3);
\draw[draw=HannoverOber!70!black,line width=.10pt] (1849.0,2250.0) -- (1842.5,2242.4);
\draw[draw=HannoverOber!70!black,line width=.10pt] (1849.9,2246.1) -- (1843.4,2238.5);
\draw[draw=HannoverOber!70!black,line width=.10pt] (1850.8,2242.2) -- (1844.3,2234.6);
\draw[draw=HannoverOber!70!black,line width=.10pt] (1851.7,2238.4) -- (1845.2,2230.7);
\draw[draw=HannoverOber!70!black,line width=.10pt] (1852.6,2234.5) -- (1846.1,2226.8);
\draw[draw=HannoverOber!70!black,line width=.10pt] (1853.5,2230.6) -- (1847.0,2222.9);
\draw[draw=HannoverOber!70!black,line width=.10pt] (1854.4,2226.7) -- (1847.9,2219.0);
\draw[draw=HannoverOber!70!black,line width=.10pt] (1855.3,2222.8) -- (1848.8,2215.1);
\draw[draw=HannoverOber!70!black,line width=.10pt] (1856.2,2218.9) -- (1849.7,2211.2);
\draw[draw=HannoverOber!70!black,line width=.10pt] (1857.1,2215.0) -- (1850.6,2207.3);
\draw[draw=HannoverOber!70!black,line width=.10pt] (1858.0,2211.1) -- (1851.5,2203.4);
\draw[draw=HannoverOber!70!black,line width=.10pt] (1858.9,2207.2) -- (1852.4,2199.5);
\draw[draw=HannoverOber!70!black,line width=.10pt] (1859.8,2203.3) -- (1853.3,2195.6);
\draw[draw=HannoverOber!70!black,line width=.10pt] (1860.7,2199.4) -- (1854.2,2191.7);
\draw[draw=HannoverOber!70!black,line width=.10pt] (1861.6,2195.5) -- (1855.1,2187.8);
\draw[draw=HannoverOber!70!black,line width=.10pt] (1862.5,2191.6) -- (1856.0,2183.9);
\draw[draw=HannoverOber!70!black,line width=.10pt] (1863.4,2187.7) -- (1856.9,2180.0);
\draw[draw=HannoverOber!70!black,line width=.10pt] (1864.3,2183.8) -- (1857.8,2176.1);
\draw[draw=HannoverOber!70!black,line width=.10pt] (1865.2,2179.9) -- (1858.7,2172.2);
\draw[draw=HannoverOber!70!black,line width=.10pt] (1866.1,2176.0) -- (1859.6,2168.3);
\draw[draw=HannoverOber!70!black,line width=.10pt] (1867.0,2172.1) -- (1860.5,2164.4);
\draw[draw=HannoverOber!70!black,line width=.10pt] (1867.8,2168.2) -- (1861.4,2160.5);
\draw[draw=HannoverOber!70!black,line width=.10pt] (1868.7,2164.3) -- (1862.3,2156.6);
\draw[draw=HannoverOber!70!black,line width=.10pt] (1869.6,2160.4) -- (1863.2,2152.8);
\draw[draw=HannoverOber!70!black,line width=.10pt] (1870.5,2156.5) -- (1864.1,2148.9);
\draw[draw=HannoverOber!70!black,line width=.10pt] (1871.4,2152.6) -- (1865.0,2145.0);
\draw[draw=HannoverOber!70!black,line width=.10pt] (1872.3,2148.7) -- (1865.9,2141.1);
\draw[draw=HannoverOber!70!black,line width=.10pt] (1873.2,2144.8) -- (1866.8,2137.2);
\draw[draw=HannoverOber!70!black,line width=.10pt] (1874.1,2140.9) -- (1867.7,2133.3);
\draw[draw=HannoverOber!70!black,line width=.10pt] (1875.0,2137.0) -- (1868.6,2129.4);
\draw[draw=HannoverOber!70!black,line width=.10pt] (1870.4,2131.2) -- (1880.4,2131.3);
\draw[draw=HannoverOber!70!black,line width=.10pt] (1872.7,2134.4) -- (1882.7,2134.5);
\draw[draw=HannoverOber!70!black,line width=.10pt] (1875.1,2137.6) -- (1885.1,2137.7);
\draw[draw=HannoverOber!70!black,line width=.10pt] (1877.5,2140.8) -- (1887.5,2140.9);
\draw[draw=HannoverOber!70!black,line width=.10pt] (1879.8,2144.1) -- (1889.8,2144.2);
\draw[draw=HannoverOber!70!black,line width=.10pt] (1882.2,2147.3) -- (1892.2,2147.4);
\draw[draw=HannoverOber!70!black,line width=.10pt] (1884.6,2150.5) -- (1894.6,2150.6);
\draw[draw=HannoverOber!70!black,line width=.10pt] (1886.9,2153.7) -- (1896.9,2153.8);
\draw[draw=HannoverOber!70!black,line width=.10pt] (1889.3,2157.0) -- (1899.3,2157.1);
\draw[draw=HannoverOber!70!black,line width=.10pt] (1891.7,2160.2) -- (1901.7,2160.3);
\draw[draw=HannoverOber!70!black,line width=.10pt] (1894.0,2163.4) -- (1904.0,2163.5);
\draw[draw=HannoverOber!70!black,line width=.10pt] (1896.4,2166.6) -- (1906.4,2166.7);
\draw[draw=HannoverOber!70!black,line width=.10pt] (1898.8,2169.9) -- (1908.8,2170.0);
\draw[draw=HannoverOber!70!black,line width=.10pt] (1901.1,2173.1) -- (1911.1,2173.2);
\draw[draw=HannoverOber!70!black,line width=.10pt] (1903.5,2176.3) -- (1913.5,2176.4);
\draw[draw=HannoverOber!70!black,line width=.10pt] (1905.9,2179.5) -- (1915.9,2179.6);
\draw[draw=HannoverOber!70!black,line width=.10pt] (1908.2,2182.8) -- (1918.2,2182.9);
\draw[draw=HannoverOber!70!black,line width=.10pt] (1910.6,2186.0) -- (1920.6,2186.1);
\draw[draw=HannoverOber!70!black,line width=.10pt] (1913.0,2189.2) -- (1923.0,2189.3);
\draw[draw=HannoverOber!70!black,line width=.10pt] (1915.3,2192.4) -- (1925.3,2192.5);
\draw[draw=HannoverOber!70!black,line width=.10pt] (1917.7,2195.7) -- (1927.7,2195.8);
\draw[draw=HannoverOber!70!black,line width=.10pt] (1920.1,2198.9) -- (1930.1,2199.0);
\draw[draw=HannoverOber!70!black,line width=.10pt] (1922.4,2202.1) -- (1932.4,2202.2);
\draw[mapBand,draw=HannoverAller,line width=.61mm] (eckberg) -- (brelingerberg) -- (falkenberg) -- (celle) -- (osterberg-e) -- (hannover) -- (eckberg);
\draw[draw=HannoverAller!70!black,line width=.10pt] (1641.2,1348.9) -- (1646.8,1340.7);
\draw[draw=HannoverAller!70!black,line width=.10pt] (1645.2,1348.7) -- (1650.8,1340.5);
\draw[draw=HannoverAller!70!black,line width=.10pt] (1649.2,1348.6) -- (1654.8,1340.3);
\draw[draw=HannoverAller!70!black,line width=.10pt] (1653.2,1348.4) -- (1658.8,1340.1);
\draw[draw=HannoverAller!70!black,line width=.10pt] (1657.2,1348.2) -- (1662.8,1339.9);
\draw[draw=HannoverAller!70!black,line width=.10pt] (1661.2,1348.0) -- (1666.8,1339.7);
\draw[draw=HannoverAller!70!black,line width=.10pt] (1665.2,1347.8) -- (1670.8,1339.5);
\draw[draw=HannoverAller!70!black,line width=.10pt] (1669.2,1347.6) -- (1674.8,1339.3);
\draw[draw=HannoverAller!70!black,line width=.10pt] (1673.2,1347.4) -- (1678.8,1339.1);
\draw[draw=HannoverAller!70!black,line width=.10pt] (1677.2,1347.2) -- (1682.8,1338.9);
\draw[draw=HannoverAller!70!black,line width=.10pt] (1681.1,1347.0) -- (1686.7,1338.7);
\draw[draw=HannoverAller!70!black,line width=.10pt] (1685.1,1346.8) -- (1690.7,1338.5);
\draw[draw=HannoverAller!70!black,line width=.10pt] (1689.1,1346.6) -- (1694.7,1338.3);
\draw[draw=HannoverAller!70!black,line width=.10pt] (1693.1,1346.4) -- (1698.7,1338.1);
\draw[draw=HannoverAller!70!black,line width=.10pt] (1697.1,1346.2) -- (1702.7,1337.9);
\draw[draw=HannoverAller!70!black,line width=.10pt] (1701.1,1346.0) -- (1706.7,1337.7);
\draw[draw=HannoverAller!70!black,line width=.10pt] (1705.1,1345.8) -- (1710.7,1337.5);
\draw[draw=HannoverAller!70!black,line width=.10pt] (1709.1,1345.6) -- (1714.7,1337.3);
\draw[draw=HannoverAller!70!black,line width=.10pt] (1713.1,1345.4) -- (1718.7,1337.1);
\draw[draw=HannoverAller!70!black,line width=.10pt] (1717.1,1345.2) -- (1722.7,1336.9);
\draw[draw=HannoverAller!70!black,line width=.10pt] (1721.1,1345.0) -- (1726.7,1336.7);
\draw[draw=HannoverAller!70!black,line width=.10pt] (1725.1,1344.8) -- (1730.7,1336.5);
\draw[draw=HannoverAller!70!black,line width=.10pt] (1729.1,1344.6) -- (1734.7,1336.3);
\draw[draw=HannoverAller!70!black,line width=.10pt] (1733.1,1344.4) -- (1738.7,1336.1);
\draw[draw=HannoverAller!70!black,line width=.10pt] (1737.1,1344.2) -- (1742.7,1335.9);
\draw[draw=HannoverAller!70!black,line width=.10pt] (1741.1,1344.0) -- (1746.7,1335.7);
\draw[draw=HannoverAller!70!black,line width=.10pt] (1745.1,1343.8) -- (1750.7,1335.5);
\draw[draw=HannoverAller!70!black,line width=.10pt] (1749.1,1343.6) -- (1754.7,1335.3);
\draw[draw=HannoverAller!70!black,line width=.10pt] (1753.1,1343.4) -- (1758.7,1335.1);
\draw[draw=HannoverAller!70!black,line width=.10pt] (1757.1,1343.2) -- (1762.7,1334.9);
\draw[draw=HannoverAller!70!black,line width=.10pt] (1761.1,1343.0) -- (1766.6,1334.7);
\draw[draw=HannoverAller!70!black,line width=.10pt] (1765.0,1342.8) -- (1770.6,1334.5);
\draw[draw=HannoverAller!70!black,line width=.10pt] (1769.0,1342.6) -- (1774.6,1334.3);
\draw[draw=HannoverAller!70!black,line width=.10pt] (1773.0,1342.4) -- (1778.6,1334.1);
\draw[draw=HannoverAller!70!black,line width=.10pt] (1777.0,1342.2) -- (1782.6,1334.0);
\draw[draw=HannoverAller!70!black,line width=.10pt] (1781.0,1342.0) -- (1786.6,1333.8);
\draw[draw=HannoverAller!70!black,line width=.10pt] (1785.0,1341.8) -- (1790.6,1333.6);
\draw[draw=HannoverAller!70!black,line width=.10pt] (1789.0,1341.6) -- (1794.6,1333.4);
\draw[draw=HannoverAller!70!black,line width=.10pt] (1793.0,1341.4) -- (1798.6,1333.2);
\draw[draw=HannoverAller!70!black,line width=.10pt] (1806.1,1337.6) -- (1801.0,1329.0);
\draw[draw=HannoverAller!70!black,line width=.10pt] (1807.6,1333.9) -- (1802.5,1325.3);
\draw[draw=HannoverAller!70!black,line width=.10pt] (1809.1,1330.2) -- (1804.0,1321.6);
\draw[draw=HannoverAller!70!black,line width=.10pt] (1810.7,1326.5) -- (1805.6,1317.9);
\draw[draw=HannoverAller!70!black,line width=.10pt] (1812.2,1322.8) -- (1807.1,1314.2);
\draw[draw=HannoverAller!70!black,line width=.10pt] (1813.7,1319.1) -- (1808.6,1310.5);
\draw[draw=HannoverAller!70!black,line width=.10pt] (1815.2,1315.4) -- (1810.1,1306.8);
\draw[draw=HannoverAller!70!black,line width=.10pt] (1816.8,1311.7) -- (1811.7,1303.1);
\draw[draw=HannoverAller!70!black,line width=.10pt] (1818.3,1308.0) -- (1813.2,1299.4);
\draw[draw=HannoverAller!70!black,line width=.10pt] (1819.8,1304.3) -- (1814.7,1295.7);
\draw[draw=HannoverAller!70!black,line width=.10pt] (1821.4,1300.6) -- (1816.3,1292.0);
\draw[draw=HannoverAller!70!black,line width=.10pt] (1822.9,1296.9) -- (1817.8,1288.3);
\draw[draw=HannoverAller!70!black,line width=.10pt] (1824.4,1293.2) -- (1819.3,1284.6);
\draw[draw=HannoverAller!70!black,line width=.10pt] (1825.9,1289.6) -- (1820.8,1280.9);
\draw[draw=HannoverAller!70!black,line width=.10pt] (1827.5,1285.9) -- (1822.4,1277.3);
\draw[draw=HannoverAller!70!black,line width=.10pt] (1829.0,1282.2) -- (1823.9,1273.6);
\draw[draw=HannoverAller!70!black,line width=.10pt] (1830.5,1278.5) -- (1825.4,1269.9);
\draw[draw=HannoverAller!70!black,line width=.10pt] (1832.1,1274.8) -- (1827.0,1266.2);
\draw[draw=HannoverAller!70!black,line width=.10pt] (1833.6,1271.1) -- (1828.5,1262.5);
\draw[draw=HannoverAller!70!black,line width=.10pt] (1835.1,1267.4) -- (1830.0,1258.8);
\draw[draw=HannoverAller!70!black,line width=.10pt] (1836.6,1263.7) -- (1831.5,1255.1);
\draw[draw=HannoverAller!70!black,line width=.10pt] (1838.2,1260.0) -- (1833.1,1251.4);
\draw[draw=HannoverAller!70!black,line width=.10pt] (1839.7,1256.3) -- (1834.6,1247.7);
\draw[draw=HannoverAller!70!black,line width=.10pt] (1841.2,1252.6) -- (1836.1,1244.0);
\draw[draw=HannoverAller!70!black,line width=.10pt] (1842.8,1248.9) -- (1837.7,1240.3);
\draw[draw=HannoverAller!70!black,line width=.10pt] (1844.3,1245.2) -- (1839.2,1236.6);
\draw[draw=HannoverAller!70!black,line width=.10pt] (1845.8,1241.5) -- (1840.7,1232.9);
\draw[draw=HannoverAller!70!black,line width=.10pt] (1847.3,1237.8) -- (1842.2,1229.2);
\draw[draw=HannoverAller!70!black,line width=.10pt] (1848.9,1234.1) -- (1843.8,1225.5);
\draw[draw=HannoverAller!70!black,line width=.10pt] (1850.4,1230.4) -- (1845.3,1221.8);
\draw[draw=HannoverAller!70!black,line width=.10pt] (1851.9,1226.7) -- (1846.8,1218.1);
\draw[draw=HannoverAller!70!black,line width=.10pt] (1853.5,1223.0) -- (1848.4,1214.4);
\draw[draw=HannoverAller!70!black,line width=.10pt] (1855.0,1219.3) -- (1849.9,1210.7);
\draw[draw=HannoverAller!70!black,line width=.10pt] (1856.5,1215.6) -- (1851.4,1207.0);
\draw[draw=HannoverAller!70!black,line width=.10pt] (1858.0,1211.9) -- (1852.9,1203.3);
\draw[draw=HannoverAller!70!black,line width=.10pt] (1859.6,1208.2) -- (1854.5,1199.6);
\draw[draw=HannoverAller!70!black,line width=.10pt] (1861.1,1204.5) -- (1856.0,1195.9);
\draw[draw=HannoverAller!70!black,line width=.10pt] (1862.6,1200.8) -- (1857.5,1192.2);
\draw[draw=HannoverAller!70!black,line width=.10pt] (1864.2,1197.1) -- (1859.1,1188.5);
\draw[draw=HannoverAller!70!black,line width=.10pt] (1865.7,1193.4) -- (1860.6,1184.8);
\draw[draw=HannoverAller!70!black,line width=.10pt] (1867.2,1189.7) -- (1862.1,1181.1);
\draw[draw=HannoverAller!70!black,line width=.10pt] (1868.7,1186.1) -- (1863.6,1177.4);
\draw[draw=HannoverAller!70!black,line width=.10pt] (1870.3,1182.4) -- (1865.2,1173.8);
\draw[draw=HannoverAller!70!black,line width=.10pt] (1871.8,1178.7) -- (1866.7,1170.1);
\draw[draw=HannoverAller!70!black,line width=.10pt] (1873.3,1175.0) -- (1868.2,1166.4);
\draw[draw=HannoverAller!70!black,line width=.10pt] (1874.9,1171.3) -- (1869.8,1162.7);
\draw[draw=HannoverAller!70!black,line width=.10pt] (1876.4,1167.6) -- (1871.3,1159.0);
\draw[draw=HannoverAller!70!black,line width=.10pt] (1877.9,1163.9) -- (1872.8,1155.3);
\draw[draw=HannoverAller!70!black,line width=.10pt] (1879.4,1160.2) -- (1874.3,1151.6);
\draw[draw=HannoverAller!70!black,line width=.10pt] (1881.0,1156.5) -- (1875.9,1147.9);
\draw[draw=HannoverAller!70!black,line width=.10pt] (1882.5,1152.8) -- (1877.4,1144.2);
\draw[draw=HannoverAller!70!black,line width=.10pt] (1884.0,1149.1) -- (1878.9,1140.5);
\draw[draw=HannoverAller!70!black,line width=.10pt] (1885.6,1145.4) -- (1880.5,1136.8);
\draw[draw=HannoverAller!70!black,line width=.10pt] (1887.1,1141.7) -- (1882.0,1133.1);
\draw[draw=HannoverAller!70!black,line width=.10pt] (1888.6,1138.0) -- (1883.5,1129.4);
\draw[draw=HannoverAller!70!black,line width=.10pt] (1890.1,1134.3) -- (1885.0,1125.7);
\draw[draw=HannoverAller!70!black,line width=.10pt] (1891.7,1130.6) -- (1886.6,1122.0);
\draw[draw=HannoverAller!70!black,line width=.10pt] (1893.2,1126.9) -- (1888.1,1118.3);
\draw[draw=HannoverAller!70!black,line width=.10pt] (1894.7,1123.2) -- (1889.6,1114.6);
\draw[draw=HannoverAller!70!black,line width=.10pt] (1896.3,1119.5) -- (1891.2,1110.9);
\draw[draw=HannoverAller!70!black,line width=.10pt] (1897.8,1115.8) -- (1892.7,1107.2);
\draw[draw=HannoverAller!70!black,line width=.10pt] (1899.3,1112.1) -- (1894.2,1103.5);
\draw[draw=HannoverAller!70!black,line width=.10pt] (1900.8,1108.4) -- (1895.7,1099.8);
\draw[draw=HannoverAller!70!black,line width=.10pt] (1897.0,1102.9) -- (1907.0,1104.0);
\draw[draw=HannoverAller!70!black,line width=.10pt] (1899.1,1106.3) -- (1909.0,1107.5);
\draw[draw=HannoverAller!70!black,line width=.10pt] (1901.1,1109.8) -- (1911.0,1111.0);
\draw[draw=HannoverAller!70!black,line width=.10pt] (1903.1,1113.3) -- (1913.0,1114.4);
\draw[draw=HannoverAller!70!black,line width=.10pt] (1905.1,1116.7) -- (1915.0,1117.9);
\draw[draw=HannoverAller!70!black,line width=.10pt] (1907.1,1120.2) -- (1917.0,1121.3);
\draw[draw=HannoverAller!70!black,line width=.10pt] (1909.1,1123.6) -- (1919.0,1124.8);
\draw[draw=HannoverAller!70!black,line width=.10pt] (1911.1,1127.1) -- (1921.0,1128.3);
\draw[draw=HannoverAller!70!black,line width=.10pt] (1913.1,1130.5) -- (1923.0,1131.7);
\draw[draw=HannoverAller!70!black,line width=.10pt] (1915.1,1134.0) -- (1925.1,1135.2);
\draw[draw=HannoverAller!70!black,line width=.10pt] (1917.1,1137.5) -- (1927.1,1138.6);
\draw[draw=HannoverAller!70!black,line width=.10pt] (1919.1,1140.9) -- (1929.1,1142.1);
\draw[draw=HannoverAller!70!black,line width=.10pt] (1921.2,1144.4) -- (1931.1,1145.6);
\draw[draw=HannoverAller!70!black,line width=.10pt] (1923.2,1147.8) -- (1933.1,1149.0);
\draw[draw=HannoverAller!70!black,line width=.10pt] (1925.2,1151.3) -- (1935.1,1152.5);
\draw[draw=HannoverAller!70!black,line width=.10pt] (1927.2,1154.8) -- (1937.1,1155.9);
\draw[draw=HannoverAller!70!black,line width=.10pt] (1929.2,1158.2) -- (1939.1,1159.4);
\draw[draw=HannoverAller!70!black,line width=.10pt] (1931.2,1161.7) -- (1941.1,1162.8);
\draw[draw=HannoverAller!70!black,line width=.10pt] (1933.2,1165.1) -- (1943.1,1166.3);
\draw[draw=HannoverAller!70!black,line width=.10pt] (1935.2,1168.6) -- (1945.1,1169.8);
\draw[draw=HannoverAller!70!black,line width=.10pt] (1937.2,1172.1) -- (1947.2,1173.2);
\draw[draw=HannoverAller!70!black,line width=.10pt] (1939.2,1175.5) -- (1949.2,1176.7);
\draw[draw=HannoverAller!70!black,line width=.10pt] (1941.2,1179.0) -- (1951.2,1180.1);
\draw[draw=HannoverAller!70!black,line width=.10pt] (1943.2,1182.4) -- (1953.2,1183.6);
\draw[draw=HannoverAller!70!black,line width=.10pt] (1945.3,1185.9) -- (1955.2,1187.1);
\draw[draw=HannoverAller!70!black,line width=.10pt] (1947.3,1189.3) -- (1957.2,1190.5);
\draw[draw=HannoverAller!70!black,line width=.10pt] (1949.3,1192.8) -- (1959.2,1194.0);
\draw[draw=HannoverAller!70!black,line width=.10pt] (1951.3,1196.3) -- (1961.2,1197.4);
\draw[draw=HannoverAller!70!black,line width=.10pt] (1953.3,1199.7) -- (1963.2,1200.9);
\draw[draw=HannoverAller!70!black,line width=.10pt] (1955.3,1203.2) -- (1965.2,1204.4);
\draw[draw=HannoverAller!70!black,line width=.10pt] (1957.3,1206.6) -- (1967.2,1207.8);
\draw[draw=HannoverAller!70!black,line width=.10pt] (1959.3,1210.1) -- (1969.3,1211.3);
\draw[draw=HannoverAller!70!black,line width=.10pt] (1961.3,1213.6) -- (1971.3,1214.7);
\draw[draw=HannoverAller!70!black,line width=.10pt] (1963.3,1217.0) -- (1973.3,1218.2);
\draw[draw=HannoverAller!70!black,line width=.10pt] (1965.3,1220.5) -- (1975.3,1221.6);
\draw[draw=HannoverAller!70!black,line width=.10pt] (1967.4,1223.9) -- (1977.3,1225.1);
\draw[draw=HannoverAller!70!black,line width=.10pt] (1969.4,1227.4) -- (1979.3,1228.6);
\draw[draw=HannoverAller!70!black,line width=.10pt] (1971.4,1230.9) -- (1981.3,1232.0);
\draw[draw=HannoverAller!70!black,line width=.10pt] (1973.4,1234.3) -- (1983.3,1235.5);
\draw[draw=HannoverAller!70!black,line width=.10pt] (1975.4,1237.8) -- (1985.3,1238.9);
\draw[draw=HannoverAller!70!black,line width=.10pt] (1977.4,1241.2) -- (1987.3,1242.4);
\draw[draw=HannoverAller!70!black,line width=.10pt] (1979.4,1244.7) -- (1989.3,1245.9);
\draw[draw=HannoverAller!70!black,line width=.10pt] (1981.4,1248.1) -- (1991.4,1249.3);
\draw[draw=HannoverAller!70!black,line width=.10pt] (1983.4,1251.6) -- (1993.4,1252.8);
\draw[draw=HannoverAller!70!black,line width=.10pt] (1985.4,1255.1) -- (1995.4,1256.2);
\draw[draw=HannoverAller!70!black,line width=.10pt] (1987.4,1258.5) -- (1997.4,1259.7);
\draw[draw=HannoverAller!70!black,line width=.10pt] (1989.5,1262.0) -- (1999.4,1263.2);
\draw[draw=HannoverAller!70!black,line width=.10pt] (1991.5,1265.4) -- (2001.4,1266.6);
\draw[draw=HannoverAller!70!black,line width=.10pt] (1993.5,1268.9) -- (2003.4,1270.1);
\draw[draw=HannoverAller!70!black,line width=.10pt] (1995.5,1272.4) -- (2005.4,1273.5);
\draw[draw=HannoverAller!70!black,line width=.10pt] (1997.5,1275.8) -- (2007.4,1277.0);
\draw[draw=HannoverAller!70!black,line width=.10pt] (1999.5,1279.3) -- (2009.4,1280.4);
\draw[draw=HannoverAller!70!black,line width=.10pt] (2001.5,1282.7) -- (2011.4,1283.9);
\draw[draw=HannoverAller!70!black,line width=.10pt] (2003.5,1286.2) -- (2013.4,1287.4);
\draw[draw=HannoverAller!70!black,line width=.10pt] (2005.5,1289.7) -- (2015.5,1290.8);
\draw[draw=HannoverAller!70!black,line width=.10pt] (2007.5,1293.1) -- (2017.5,1294.3);
\draw[draw=HannoverAller!70!black,line width=.10pt] (2011.3,1296.2) -- (2013.6,1305.9);
\draw[draw=HannoverAller!70!black,line width=.10pt] (2008.7,1299.3) -- (2011.0,1309.0);
\draw[draw=HannoverAller!70!black,line width=.10pt] (2006.2,1302.3) -- (2008.5,1312.1);
\draw[draw=HannoverAller!70!black,line width=.10pt] (2003.6,1305.4) -- (2005.9,1315.1);
\draw[draw=HannoverAller!70!black,line width=.10pt] (2001.0,1308.5) -- (2003.3,1318.2);
\draw[draw=HannoverAller!70!black,line width=.10pt] (1998.5,1311.6) -- (2000.8,1321.3);
\draw[draw=HannoverAller!70!black,line width=.10pt] (1995.9,1314.6) -- (1998.2,1324.4);
\draw[draw=HannoverAller!70!black,line width=.10pt] (1993.3,1317.7) -- (1995.6,1327.4);
\draw[draw=HannoverAller!70!black,line width=.10pt] (1990.8,1320.8) -- (1993.1,1330.5);
\draw[draw=HannoverAller!70!black,line width=.10pt] (1988.2,1323.8) -- (1990.5,1333.6);
\draw[draw=HannoverAller!70!black,line width=.10pt] (1985.7,1326.9) -- (1987.9,1336.6);
\draw[draw=HannoverAller!70!black,line width=.10pt] (1983.1,1330.0) -- (1985.4,1339.7);
\draw[draw=HannoverAller!70!black,line width=.10pt] (1980.5,1333.0) -- (1982.8,1342.8);
\draw[draw=HannoverAller!70!black,line width=.10pt] (1978.0,1336.1) -- (1980.3,1345.9);
\draw[draw=HannoverAller!70!black,line width=.10pt] (1975.4,1339.2) -- (1977.7,1348.9);
\draw[draw=HannoverAller!70!black,line width=.10pt] (1972.8,1342.3) -- (1975.1,1352.0);
\draw[draw=HannoverAller!70!black,line width=.10pt] (1970.3,1345.3) -- (1972.6,1355.1);
\draw[draw=HannoverAller!70!black,line width=.10pt] (1967.7,1348.4) -- (1970.0,1358.1);
\draw[draw=HannoverAller!70!black,line width=.10pt] (1965.1,1351.5) -- (1967.4,1361.2);
\draw[draw=HannoverAller!70!black,line width=.10pt] (1962.6,1354.5) -- (1964.9,1364.3);
\draw[draw=HannoverAller!70!black,line width=.10pt] (1960.0,1357.6) -- (1962.3,1367.3);
\draw[draw=HannoverAller!70!black,line width=.10pt] (1957.4,1360.7) -- (1959.7,1370.4);
\draw[draw=HannoverAller!70!black,line width=.10pt] (1954.9,1363.8) -- (1957.2,1373.5);
\draw[draw=HannoverAller!70!black,line width=.10pt] (1952.3,1366.8) -- (1954.6,1376.6);
\draw[draw=HannoverAller!70!black,line width=.10pt] (1949.8,1369.9) -- (1952.1,1379.6);
\draw[draw=HannoverAller!70!black,line width=.10pt] (1947.2,1373.0) -- (1949.5,1382.7);
\draw[draw=HannoverAller!70!black,line width=.10pt] (1944.6,1376.0) -- (1946.9,1385.8);
\draw[draw=HannoverAller!70!black,line width=.10pt] (1942.1,1379.1) -- (1944.4,1388.8);
\draw[draw=HannoverAller!70!black,line width=.10pt] (1939.5,1382.2) -- (1941.8,1391.9);
\draw[draw=HannoverAller!70!black,line width=.10pt] (1936.9,1385.2) -- (1939.2,1395.0);
\draw[draw=HannoverAller!70!black,line width=.10pt] (1930.4,1393.0) -- (1932.2,1402.9);
\draw[draw=HannoverAller!70!black,line width=.10pt] (1927.7,1396.0) -- (1929.5,1405.8);
\draw[draw=HannoverAller!70!black,line width=.10pt] (1925.0,1398.9) -- (1926.8,1408.8);
\draw[draw=HannoverAller!70!black,line width=.10pt] (1922.3,1401.9) -- (1924.1,1411.7);
\draw[draw=HannoverAller!70!black,line width=.10pt] (1919.6,1404.9) -- (1921.4,1414.7);
\draw[draw=HannoverAller!70!black,line width=.10pt] (1916.9,1407.8) -- (1918.7,1417.6);
\draw[draw=HannoverAller!70!black,line width=.10pt] (1914.2,1410.8) -- (1916.0,1420.6);
\draw[draw=HannoverAller!70!black,line width=.10pt] (1911.5,1413.7) -- (1913.3,1423.5);
\draw[draw=HannoverAller!70!black,line width=.10pt] (1908.8,1416.7) -- (1910.6,1426.5);
\draw[draw=HannoverAller!70!black,line width=.10pt] (1906.1,1419.6) -- (1907.9,1429.4);
\draw[draw=HannoverAller!70!black,line width=.10pt] (1903.4,1422.6) -- (1905.2,1432.4);
\draw[draw=HannoverAller!70!black,line width=.10pt] (1900.7,1425.5) -- (1902.5,1435.3);
\draw[draw=HannoverAller!70!black,line width=.10pt] (1898.0,1428.5) -- (1899.9,1438.3);
\draw[draw=HannoverAller!70!black,line width=.10pt] (1895.3,1431.4) -- (1897.2,1441.3);
\draw[draw=HannoverAller!70!black,line width=.10pt] (1892.6,1434.4) -- (1894.5,1444.2);
\draw[draw=HannoverAller!70!black,line width=.10pt] (1889.9,1437.3) -- (1891.8,1447.2);
\draw[draw=HannoverAller!70!black,line width=.10pt] (1887.2,1440.3) -- (1889.1,1450.1);
\draw[draw=HannoverAller!70!black,line width=.10pt] (1884.5,1443.2) -- (1886.4,1453.1);
\draw[draw=HannoverAller!70!black,line width=.10pt] (1881.8,1446.2) -- (1883.7,1456.0);
\draw[draw=HannoverAller!70!black,line width=.10pt] (1879.1,1449.1) -- (1881.0,1459.0);
\draw[draw=HannoverAller!70!black,line width=.10pt] (1876.4,1452.1) -- (1878.3,1461.9);
\draw[draw=HannoverAller!70!black,line width=.10pt] (1873.7,1455.0) -- (1875.6,1464.9);
\draw[draw=HannoverAller!70!black,line width=.10pt] (1871.0,1458.0) -- (1872.9,1467.8);
\draw[draw=HannoverAller!70!black,line width=.10pt] (1868.3,1461.0) -- (1870.2,1470.8);
\draw[draw=HannoverAller!70!black,line width=.10pt] (1865.6,1463.9) -- (1867.5,1473.7);
\draw[draw=HannoverAller!70!black,line width=.10pt] (1862.9,1466.9) -- (1864.8,1476.7);
\draw[draw=HannoverAller!70!black,line width=.10pt] (1860.2,1469.8) -- (1862.1,1479.6);
\draw[draw=HannoverAller!70!black,line width=.10pt] (1857.5,1472.8) -- (1859.4,1482.6);
\draw[draw=HannoverAller!70!black,line width=.10pt] (1854.8,1475.7) -- (1856.7,1485.5);
\draw[draw=HannoverAller!70!black,line width=.10pt] (1852.1,1478.7) -- (1854.0,1488.5);
\draw[draw=HannoverAller!70!black,line width=.10pt] (1849.4,1481.6) -- (1851.3,1491.5);
\draw[draw=HannoverAller!70!black,line width=.10pt] (1846.7,1484.6) -- (1848.6,1494.4);
\draw[draw=HannoverAller!70!black,line width=.10pt] (1844.0,1487.5) -- (1845.9,1497.4);
\draw[draw=HannoverAller!70!black,line width=.10pt] (1841.3,1490.5) -- (1843.2,1500.3);
\draw[draw=HannoverAller!70!black,line width=.10pt] (1838.6,1493.4) -- (1840.5,1503.3);
\draw[draw=HannoverAller!70!black,line width=.10pt] (1835.9,1496.4) -- (1837.8,1506.2);
\draw[draw=HannoverAller!70!black,line width=.10pt] (1833.2,1499.3) -- (1835.1,1509.2);
\draw[draw=HannoverAller!70!black,line width=.10pt] (1830.5,1502.3) -- (1832.4,1512.1);
\draw[draw=HannoverAller!70!black,line width=.10pt] (1829.9,1507.3) -- (1820.1,1509.4);
\draw[draw=HannoverAller!70!black,line width=.10pt] (1826.9,1504.7) -- (1817.1,1506.7);
\draw[draw=HannoverAller!70!black,line width=.10pt] (1823.9,1502.0) -- (1814.1,1504.1);
\draw[draw=HannoverAller!70!black,line width=.10pt] (1820.9,1499.4) -- (1811.1,1501.4);
\draw[draw=HannoverAller!70!black,line width=.10pt] (1817.9,1496.7) -- (1808.1,1498.8);
\draw[draw=HannoverAller!70!black,line width=.10pt] (1814.9,1494.1) -- (1805.1,1496.1);
\draw[draw=HannoverAller!70!black,line width=.10pt] (1811.9,1491.5) -- (1802.1,1493.5);
\draw[draw=HannoverAller!70!black,line width=.10pt] (1808.9,1488.8) -- (1799.1,1490.8);
\draw[draw=HannoverAller!70!black,line width=.10pt] (1805.9,1486.2) -- (1796.1,1488.2);
\draw[draw=HannoverAller!70!black,line width=.10pt] (1802.9,1483.5) -- (1793.1,1485.5);
\draw[draw=HannoverAller!70!black,line width=.10pt] (1799.9,1480.9) -- (1790.1,1482.9);
\draw[draw=HannoverAller!70!black,line width=.10pt] (1796.9,1478.2) -- (1787.1,1480.2);
\draw[draw=HannoverAller!70!black,line width=.10pt] (1793.9,1475.6) -- (1784.1,1477.6);
\draw[draw=HannoverAller!70!black,line width=.10pt] (1790.9,1472.9) -- (1781.1,1474.9);
\draw[draw=HannoverAller!70!black,line width=.10pt] (1787.9,1470.3) -- (1778.1,1472.3);
\draw[draw=HannoverAller!70!black,line width=.10pt] (1784.9,1467.6) -- (1775.1,1469.7);
\draw[draw=HannoverAller!70!black,line width=.10pt] (1781.9,1465.0) -- (1772.1,1467.0);
\draw[draw=HannoverAller!70!black,line width=.10pt] (1778.9,1462.3) -- (1769.1,1464.4);
\draw[draw=HannoverAller!70!black,line width=.10pt] (1775.9,1459.7) -- (1766.1,1461.7);
\draw[draw=HannoverAller!70!black,line width=.10pt] (1772.9,1457.0) -- (1763.1,1459.1);
\draw[draw=HannoverAller!70!black,line width=.10pt] (1769.9,1454.4) -- (1760.1,1456.4);
\draw[draw=HannoverAller!70!black,line width=.10pt] (1766.9,1451.7) -- (1757.1,1453.8);
\draw[draw=HannoverAller!70!black,line width=.10pt] (1763.9,1449.1) -- (1754.1,1451.1);
\draw[draw=HannoverAller!70!black,line width=.10pt] (1760.9,1446.4) -- (1751.1,1448.5);
\draw[draw=HannoverAller!70!black,line width=.10pt] (1757.9,1443.8) -- (1748.1,1445.8);
\draw[draw=HannoverAller!70!black,line width=.10pt] (1754.9,1441.2) -- (1745.1,1443.2);
\draw[draw=HannoverAller!70!black,line width=.10pt] (1751.9,1438.5) -- (1742.1,1440.5);
\draw[draw=HannoverAller!70!black,line width=.10pt] (1748.9,1435.9) -- (1739.1,1437.9);
\draw[draw=HannoverAller!70!black,line width=.10pt] (1745.9,1433.2) -- (1736.1,1435.2);
\draw[draw=HannoverAller!70!black,line width=.10pt] (1742.9,1430.6) -- (1733.2,1432.6);
\draw[draw=HannoverAller!70!black,line width=.10pt] (1739.9,1427.9) -- (1730.2,1429.9);
\draw[draw=HannoverAller!70!black,line width=.10pt] (1736.9,1425.3) -- (1727.2,1427.3);
\draw[draw=HannoverAller!70!black,line width=.10pt] (1733.9,1422.6) -- (1724.2,1424.6);
\draw[draw=HannoverAller!70!black,line width=.10pt] (1731.0,1420.0) -- (1721.2,1422.0);
\draw[draw=HannoverAller!70!black,line width=.10pt] (1728.0,1417.3) -- (1718.2,1419.3);
\draw[draw=HannoverAller!70!black,line width=.10pt] (1725.0,1414.7) -- (1715.2,1416.7);
\draw[draw=HannoverAller!70!black,line width=.10pt] (1722.0,1412.0) -- (1712.2,1414.1);
\draw[draw=HannoverAller!70!black,line width=.10pt] (1719.0,1409.4) -- (1709.2,1411.4);
\draw[draw=HannoverAller!70!black,line width=.10pt] (1716.0,1406.7) -- (1706.2,1408.8);
\draw[draw=HannoverAller!70!black,line width=.10pt] (1713.0,1404.1) -- (1703.2,1406.1);
\draw[draw=HannoverAller!70!black,line width=.10pt] (1710.0,1401.4) -- (1700.2,1403.5);
\draw[draw=HannoverAller!70!black,line width=.10pt] (1707.0,1398.8) -- (1697.2,1400.8);
\draw[draw=HannoverAller!70!black,line width=.10pt] (1704.0,1396.1) -- (1694.2,1398.2);
\draw[draw=HannoverAller!70!black,line width=.10pt] (1701.0,1393.5) -- (1691.2,1395.5);
\draw[draw=HannoverAller!70!black,line width=.10pt] (1698.0,1390.8) -- (1688.2,1392.9);
\draw[draw=HannoverAller!70!black,line width=.10pt] (1695.0,1388.2) -- (1685.2,1390.2);
\draw[draw=HannoverAller!70!black,line width=.10pt] (1692.0,1385.6) -- (1682.2,1387.6);
\draw[draw=HannoverAller!70!black,line width=.10pt] (1689.0,1382.9) -- (1679.2,1384.9);
\draw[draw=HannoverAller!70!black,line width=.10pt] (1686.0,1380.3) -- (1676.2,1382.3);
\draw[draw=HannoverAller!70!black,line width=.10pt] (1683.0,1377.6) -- (1673.2,1379.6);
\draw[draw=HannoverAller!70!black,line width=.10pt] (1680.0,1375.0) -- (1670.2,1377.0);
\draw[draw=HannoverAller!70!black,line width=.10pt] (1677.0,1372.3) -- (1667.2,1374.3);
\draw[draw=HannoverAller!70!black,line width=.10pt] (1674.0,1369.7) -- (1664.2,1371.7);
\draw[draw=HannoverAller!70!black,line width=.10pt] (1671.0,1367.0) -- (1661.2,1369.0);
\draw[draw=HannoverAller!70!black,line width=.10pt] (1668.0,1364.4) -- (1658.2,1366.4);
\draw[draw=HannoverAller!70!black,line width=.10pt] (1665.0,1361.7) -- (1655.2,1363.8);
\draw[draw=HannoverAller!70!black,line width=.10pt] (1662.0,1359.1) -- (1652.2,1361.1);
\draw[draw=HannoverAller!70!black,line width=.10pt] (1659.0,1356.4) -- (1649.2,1358.5);
\draw[draw=HannoverAller!70!black,line width=.10pt] (1656.0,1353.8) -- (1646.2,1355.8);
\draw[draw=HannoverAller!70!black,line width=.10pt] (1653.0,1351.1) -- (1643.2,1353.2);
\draw[draw=HannoverAller!70!black,line width=.10pt] (1650.0,1348.5) -- (1640.2,1350.5);
\draw[mapBand,draw=HannoverHaupt,line width=0.58mm] (903.8,446.1) -- (1021.8,601.1);
\draw[mapBand,draw=HannoverWest,line width=0.56mm] (896.5,451.7) -- (1014.5,606.7);
\draw[mapBand,draw=HannoverHaupt,line width=0.58mm] (1350.4,886.4) -- (1599.4,964.4);
\draw[mapBand,draw=HannoverWest,line width=0.56mm] (1347.5,895.9) -- (1596.5,973.9);
\draw[mapBand,draw=HannoverHaupt,line width=0.58mm] (2112.0,1737.5) -- (2291.0,2015.5);
\draw[mapBand,draw=HannoverLueneburg,line width=0.56mm] (2119.8,1732.6) -- (2298.8,2010.6);
\draw[mapBand,draw=HannoverHaupt,line width=0.58mm] (2264.6,640.3) -- (2180.6,747.3);
\draw[mapBand,draw=HannoverLueneburg,line width=0.56mm] (2271.2,645.5) -- (2187.2,752.5);
\draw[mapBand,draw=HannoverHaupt,line width=0.58mm] (1969.7,494.6) -- (1779.7,635.6);
\draw[mapBand,draw=HannoverBremen,line width=0.56mm] (1964.5,487.6) -- (1774.5,628.6);
\draw[mapBand,draw=HannoverWest,line width=0.58mm] (1370.8,1323.6) -- (1162.8,1592.6);
\draw[mapBand,draw=HannoverMittel,line width=0.56mm] (1377.3,1328.6) -- (1169.3,1597.6);
\draw[mapBand,draw=HannoverOst,line width=0.58mm] (657.2,367.2) -- (901.2,445.2);
\draw[mapBand,draw=HannoverWest,line width=0.56mm] (654.8,374.8) -- (898.8,452.8);
\draw[mapBand,draw=HannoverWest,dash pattern=on 1.9pt off 1.0pt] (kirchhesepe) -- (quekenberg) -- (mordkuhlenberg);
\draw[mapBand,draw=HannoverBremen,dash pattern=on 1.7pt off 1.0pt] (silberberg) -- (basdahl);
\draw[mapBand,draw=HannoverHaupt] (1980,2498) -- (2014,2552);
\draw[mapBand,draw=HannoverHaupt] (2231,2498) -- (2228,2554);
\draw[mapMain,->] (1977,2496) -- (2014,2555);
\draw[mapMain,->] (2227,2496) -- (2226,2561);

% Triangle sides, with the three stroke classes from the signs key.
\draw[mapDash] (ahaus) -- (bocholt);
\draw[mapDash] (basis-s) -- (hohenhorn);
\draw[mapDash] (bazenberg) -- (meisner);
\draw[mapDash] (beilen) -- (gieten);
\draw[mapDash] (beilen) -- (koeverden);
\draw[mapDash] (beilen) -- (sleen);
\draw[mapDash] (bentheim) -- (ahaus);
\draw[mapDash] (borkum) -- (uithuiser);
\draw[mapDash] (bremen) -- (oldenburg);
\draw[mapDash] (bremen) -- (wildeshausen);
\draw[mapDash] (deisselberg) -- (staufenberg);
\draw[mapDash] (desenberg) -- (deisselberg);
\draw[mapDash] (desenberg) -- (dommel);
\draw[mapDash] (desenberg) -- (meisner);
\draw[mapDash] (desenberg) -- (staufenberg);
\draw[mapDash] (doerenberg) -- (huenenburg);
\draw[mapDash] (emden) -- (holwierda);
\draw[mapDash] (emden) -- (midwolde);
\draw[mapDash] (festberg) -- (huyseburg);
\draw[mapDash] (gieten) -- (sleen);
\draw[mapDash] (gieten) -- (ter-apel);
\draw[mapDash] (goettingen) -- (meisner);
\draw[mapDash] (groll) -- (ahaus);
\draw[mapDash] (groll) -- (bocholt);
\draw[mapDash] (groningen) -- (midwolde);
\draw[mapDash] (groningen) -- (osterwolde);
\draw[mapDash] (groningen) -- (rolde);
\draw[mapDash] (hage) -- (558,505);
\draw[mapMain] (558,505) -- (emden);
\draw[mapDash] (hamburg) -- (basis-n);
\draw[mapDash] (hamburg) -- (basis-s);
\draw[mapDash] (hamburg) -- (segeberg);
\draw[mapDash] (hamburg) -- (syk);
\draw[mapDash] (hannover) -- (1784,1588);
\draw[mapFine] (1784,1588) -- (deister-i);
\draw[mapDash] (hannover) -- (1851,1719);
\draw[mapFine] (1851,1719) -- (hils);
\draw[mapDash] (hannover) -- (1977,1627);
\draw[mapFine] (1977,1627) -- (lichtenberg);
\draw[mapDash] (harikerberg) -- (ahaus);
\draw[mapDash] (harikerberg) -- (groll);
\draw[mapDash] (hausheide) -- (desenberg);
\draw[mapDash] (hausheide) -- (dommel);
\draw[mapDash] (hausheide) -- (spitzwarte);
\draw[mapDash] (hohehagen) -- (meisner);
\draw[mapDash] (hohehagen) -- (staufenberg);
\draw[mapDash] (hohenhorn) -- (2141,712);
\draw[mapMain] (2141,712) -- (nindorf);
\draw[mapDash] (hohenhorn) -- (2042,706);
\draw[mapMain] (2042,706) -- (wilsede);
\draw[mapDash] (holwierda) -- (groningen);
\draw[mapDash] (holwierda) -- (midwolde);
\draw[mapDash] (hornhuizen) -- (groningen);
\draw[mapDash] (hornhuizen) -- (uithuiser);
\draw[mapDash] (huenenburg) -- (hausheide);
\draw[mapDash] (huenenburg) -- (koeterberg);
\draw[mapDash] (huenenburg) -- (spitzwarte);
\draw[mapDash] (huyseburg) -- (brocken);
\draw[mapDash] (kirchhesepe) -- (koeverden);
\draw[mapDash] (kirchhesepe) -- (uelsen);
\draw[mapDash] (koeterberg) -- (desenberg);
\draw[mapDash] (koeterberg) -- (hausheide);
\draw[mapDash] (koeterberg) -- (hohehagen);
\draw[mapDash] (koeverden) -- (lemelerberg);
\draw[mapDash] (koeverden) -- (uelsen);
\draw[mapDash] (leer) -- (midwolde);
\draw[mapDash] (lemelerberg) -- (harikerberg);
\draw[mapDash] (lemelerberg) -- (uelsen);
\draw[mapDash] (luebeck) -- (syk);
\draw[mapDash] (midwolde) -- (onstwedde);
\draw[mapDash] (midwolde) -- (rolde);
\draw[mapDash] (neuwerk) -- (1232,309);
\draw[mapFine] (1232,309) -- (bremerlehe);
\draw[mapFine] (neuwerk) -- (1050,209);
\draw[mapDash] (1050,209) -- (wangeroog);
\draw[mapDash] (nonnenstein) -- (huenenburg);
\draw[mapDash] (oldenburg) -- (krapendorf);
\draw[mapDash] (oldenzaal) -- (ahaus);
\draw[mapDash] (oldenzaal) -- (bentheim);
\draw[mapDash] (oldenzaal) -- (harikerberg);
\draw[mapDash] (onstwedde) -- (gieten);
\draw[mapDash] (onstwedde) -- (ter-apel);
\draw[mapDash] (osterwolde) -- (beilen);
\draw[mapDash] (osterwolde) -- (rolde);
\draw[mapDash] (pilsum) -- (498,560);
\draw[mapMain] (498,560) -- (emden);
\draw[mapDash] (pilsum) -- (holwierda);
\draw[mapDash] (pilsum) -- (uithuiser);
\draw[mapDash] (rolde) -- (beilen);
\draw[mapDash] (rolde) -- (gieten);
\draw[mapDash] (rolde) -- (sleen);
\draw[mapDash] (segeberg) -- (luebeck);
\draw[mapDash] (segeberg) -- (syk);
\draw[mapDash] (sleen) -- (koeverden);
\draw[mapDash] (spitzwarte) -- (dommel);
\draw[mapDash] (steinberg-s) -- (meisner);
\draw[mapDash] (steinberg-s) -- (staufenberg);
\draw[mapDash] (syk) -- (basis-n);
\draw[mapDash] (syk) -- (basis-s);
\draw[mapDash] (twistringen) -- (wildeshausen);
\draw[mapDash] (uelsen) -- (bentheim);
\draw[mapDash] (uelsen) -- (harikerberg);
\draw[mapDash] (uelsen) -- (oldenzaal);
\draw[mapDash] (uithuiser) -- (groningen);
\draw[mapDash] (uithuiser) -- (holwierda);
\draw[mapDash] (wangeroog) -- (1078,367);
\draw[mapFine] (1078,367) -- (bremerlehe);
\draw[mapDash] (wangeroog) -- (891,350);
\draw[mapFine] (891,350) -- (jever);
\draw[mapDash] (wangeroog) -- (1000,346);
\draw[mapFine] (1000,346) -- (langwarden);
\draw[mapDash] (wangeroog) -- (843,264);
\draw[mapFine] (843,264) -- (spiekeroog);
\draw[mapDash] (westerstede) -- (oldenburg);
\draw[mapDash] (wildeshausen) -- (krapendorf);
\draw[mapDash] (wildeshausen) -- (oldenburg);
\draw[mapDash] (wittekindstein) -- (huenenburg);
\draw[mapFine] (altenwalde) -- (bederkesa);
\draw[mapFine] (altenwalde) -- (krempel);
\draw[mapFine] (bazenberg) -- (sonnenstein);
\draw[mapFine] (bederkesa) -- (basdahl);
\draw[mapFine] (bederkesa) -- (dolosenberg);
\draw[mapFine] (bederkesa) -- (silberberg);
\draw[mapFine] (bederkesa) -- (wuestenwohlde);
\draw[mapFine] (bederkesa-i) -- (haye);
\draw[mapFine] (bederkesa-i) -- (silberberg);
\draw[mapFine] (bocksberg) -- (fuhrenberg);
\draw[mapFine] (bocksberg) -- (haspelkopf);
\draw[mapFine] (brelingerberg) -- (deister-i);
\draw[mapFine] (brelingerberg) -- (hannover);
\draw[mapFine] (brelingerberg) -- (osterberg-e);
\draw[mapFine] (bremerlehe) -- (holssel);
\draw[mapFine] (bremerlehe) -- (wuestenwohlde);
\draw[mapFine] (brocken) -- (bazenberg);
\draw[mapFine] (brocken) -- (bocksberg);
\draw[mapFine] (brocken) -- (eversberg);
\draw[mapFine] (brocken) -- (lange-ecke);
\draw[mapFine] (brocken) -- (schalliethe);
\draw[mapFine] (brocken) -- (sonnenstein);
\draw[mapFine] (celle) -- (brelingerberg);
\draw[mapFine] (celle) -- (osterberg-e);
\draw[mapFine] (dolosenberg) -- (basdahl);
\draw[mapFine] (dolosenberg) -- (silberberg);
\draw[mapFine] (dolosenberg) -- (stade);
\draw[mapFine] (eckberg) -- (brelingerberg);
\draw[mapFine] (eckberg) -- (hannover);
\draw[mapFine] (eversberg) -- (lange-ecke);
\draw[mapFine] (eversberg) -- (schalliethe);
\draw[mapFine] (eversberg) -- (sonnenstein);
\draw[mapFine] (fahlenstollen) -- (hohehagen);
\draw[mapFine] (fahlenstollen) -- (kleper);
\draw[mapFine] (fahlenstollen) -- (steinberg-s);
\draw[mapFine] (falkenberg) -- (brelingerberg);
\draw[mapFine] (falkenberg) -- (osterberg-e);
\draw[mapFine] (fuhrenberg) -- (haspelkopf);
\draw[mapFine] (gieseberg) -- (steinberg-s);
\draw[mapFine] (goettingen) -- (gieseberg);
\draw[mapDash] (hage) -- (632,469);
\draw[mapMain] (632,469) -- (aurich);
\draw[mapFine] (haspelkopf) -- (eversberg);
\draw[mapFine] (haspelkopf) -- (sonnenstein);
\draw[mapFine] (hauseberg) -- (bazenberg);
\draw[mapFine] (hauseberg) -- (hellberg);
\draw[mapFine] (hauseberg) -- (sonnenstein);
\draw[mapFine] (haye) -- (holssel);
\draw[mapFine] (haye) -- (krempel);
\draw[mapFine] (hellberg) -- (bazenberg);
\draw[mapFine] (hellberg) -- (sonnenstein);
\draw[mapFine] (hils) -- (tockenburg);
\draw[mapFine] (hohehagen) -- (gieseberg);
\draw[mapFine] (hohehagen) -- (steinberg-s);
\draw[mapFine] (holssel) -- (bederkesa-i);
\draw[mapFine] (holssel) -- (krempel);
\draw[mapFine] (kleper) -- (gieseberg);
\draw[mapFine] (kleper) -- (hohehagen);
\draw[mapFine] (kleper) -- (steinberg-s);
\draw[mapFine] (krempel) -- (bederkesa-i);
\draw[mapFine] (krempel) -- (silberberg);
\draw[mapFine] (lange-ecke) -- (schalliethe);
\draw[mapFine] (lange-ecke) -- (sonnenstein);
\draw[mapFine] (meridianzeichen) -- (gieseberg);
\draw[mapFine] (meridianzeichen) -- (kleper);
\draw[mapFine] (osterberg-e) -- (hannover);
\draw[mapDash] (pilsum) -- (577,519);
\draw[mapMain] (577,519) -- (aurich);
\draw[mapFine] (schalliethe) -- (sonnenstein);
\draw[mapFine] (sonnenstein) -- (hohehagen);
\draw[mapFine] (tockenburg) -- (bazenberg);
\draw[mapFine] (tockenburg) -- (hauseberg);
\draw[mapFine] (tockenburg) -- (wulften);
\draw[mapFine] (weper) -- (hohehagen);
\draw[mapFine] (weper) -- (meridianzeichen);
\draw[mapFine] (wuestenwohlde) -- (basdahl);
\draw[mapFine] (wuestenwohlde) -- (dolosenberg);
\draw[mapFine] (wuestenwohlde) -- (silberberg);
\draw[mapFine] (wulften) -- (bazenberg);
\draw[mapFine] (wulften) -- (hauseberg);
\draw[mapFine] (wulften) -- (hellberg);
\draw[mapFine] (wulften) -- (sonnenstein);
\draw[mapMain] (altenwalde) -- (haye);
\draw[mapMain] (altenwalde) -- (silberberg);
\draw[mapMain] (asendorf) -- (knickberg);
\draw[mapMain] (asendorf) -- (twistringen);
\draw[mapMain] (aurich) -- (emden);
\draw[mapMain] (aurich) -- (leer);
\draw[mapMain] (baltrum) -- (dornum);
\draw[mapMain] (baltrum) -- (esens);
\draw[mapMain] (baltrum) -- (hage);
\draw[mapMain] (baltrum) -- (langeoog);
\draw[mapMain] (basdahl) -- (stade);
\draw[mapMain] (basis-n) -- (basis-s);
\draw[mapMain] (bentheim) -- (doerenberg);
\draw[mapMain] (borkum) -- (juist);
\draw[mapMain] (borkum) -- (pilsum);
\draw[mapMain] (bottel) -- (bullerberg);
\draw[mapMain] (bottel) -- (steinberg-n);
\draw[mapMain] (bottel) -- (wilsede);
\draw[mapMain] (breithorn) -- (scharnhorst);
\draw[mapMain] (bremen) -- (asendorf);
\draw[mapMain] (bremen) -- (bottel);
\draw[mapMain] (bremen) -- (bruettendorf);
\draw[mapMain] (bremen) -- (steinberg-n);
\draw[mapMain] (bremen) -- (twistringen);
\draw[mapMain] (bremen) -- (zeven);
\draw[mapMain] (bremerlehe) -- (basdahl);
\draw[mapMain] (bremerlehe) -- (brillit);
\draw[mapMain] (bremerlehe) -- (garlste);
\draw[mapMain] (bremerlehe) -- (varel);
\draw[mapMain] (bretze) -- (glienitz);
\draw[mapMain] (bretze) -- (taetendorf);
\draw[mapMain] (brillit) -- (bremen);
\draw[mapMain] (brillit) -- (garlste);
\draw[mapMain] (brillit) -- (zeven);
\draw[mapMain] (brocken) -- (hohehagen);
\draw[mapMain] (bruettendorf) -- (bottel);
\draw[mapMain] (bruettendorf) -- (bullerberg);
\draw[mapMain] (bruettendorf) -- (litberg);
\draw[mapMain] (bruettendorf) -- (steinberg-n);
\draw[mapMain] (bruettendorf) -- (wilsede);
\draw[mapMain] (crempe) -- (hamburg);
\draw[mapMain] (crempe) -- (stade);
\draw[mapMain] (deister-i) -- (garssen);
\draw[mapMain] (deister-i) -- (hils);
\draw[mapMain] (deister-i) -- (lichtenberg);
\draw[mapMain] (deister-ii) -- (osterwald);
\draw[mapMain] (doerenberg) -- (nonnenstein);
\draw[mapMain] (dornum) -- (aurich);
\draw[mapMain] (dornum) -- (esens);
\draw[mapMain] (dornum) -- (hage);
\draw[mapMain] (dornum) -- (jever);
\draw[mapMain] (elmhorst) -- (bullerberg);
\draw[mapMain] (elmhorst) -- (falkenberg);
\draw[mapMain] (emden) -- (leer);
\draw[mapMain] (emden) -- (onstwedde);
\draw[mapMain] (esens) -- (aurich);
\draw[mapMain] (esens) -- (hage);
\draw[mapMain] (esens) -- (jever);
\draw[mapMain] (falkenberg) -- (breithorn);
\draw[mapMain] (falkenberg) -- (celle);
\draw[mapMain] (falkenberg) -- (deister-i);
\draw[mapMain] (falkenberg) -- (garssen);
\draw[mapMain] (falkenberg) -- (hauselberg);
\draw[mapMain] (falkenberg) -- (lichtenberg);
\draw[mapMain] (falkenberg) -- (scharnhorst);
\draw[mapMain] (festberg) -- (brocken);
\draw[mapMain] (garlste) -- (bremen);
\draw[mapMain] (garssen) -- (celle);
\draw[mapMain] (garssen) -- (lichtenberg);
\draw[mapMain] (glienitz) -- (hoe-beck);
\draw[mapMain] (glienitz) -- (hohenmechtin);
\draw[mapMain] (glienitz) -- (taetendorf);
\draw[mapMain] (goettingen) -- (hohehagen);
\draw[mapMain] (hage) -- (pilsum);
\draw[mapMain] (hamburg) -- (hohenhorn);
\draw[mapMain] (hamburg) -- (litberg);
\draw[mapMain] (hamburg) -- (lueneburg);
\draw[mapMain] (hamburg) -- (nindorf);
\draw[mapMain] (hamburg) -- (timpenberg);
\draw[mapMain] (hamburg) -- (wilsede);
\draw[mapMain] (hankensbuettel) -- (wohlenberg);
\draw[mapMain] (hauselberg) -- (breithorn);
\draw[mapMain] (haye) -- (basdahl);
\draw[mapMain] (haye) -- (bremerlehe);
\draw[mapMain] (haye) -- (silberberg);
\draw[mapMain] (hils) -- (brocken);
\draw[mapMain] (hils) -- (hohehagen);
\draw[mapMain] (hils) -- (lichtenberg);
\draw[mapMain] (hils) -- (meridianzeichen);
\draw[mapMain] (hoe-beck) -- (hohenmechtin);
\draw[mapMain] (hoe-beck) -- (pugelatz);
\draw[mapMain] (hohenhorn) -- (lauenburg);
\draw[mapMain] (hohenhorn) -- (lueneburg);
\draw[mapMain] (hohenmechtin) -- (pugelatz);
\draw[mapMain] (hohenmechtin) -- (wieren);
\draw[mapMain] (holxerberg) -- (wieren);
\draw[mapMain] (jever) -- (aurich);
\draw[mapMain] (jever) -- (langwarden);
\draw[mapMain] (jever) -- (varel);
\draw[mapMain] (jever) -- (westerstede);
\draw[mapMain] (juist) -- (hage);
\draw[mapMain] (juist) -- (norderney);
\draw[mapMain] (juist) -- (pilsum);
\draw[mapMain] (kirchhesepe) -- (bentheim);
\draw[mapMain] (kirchhesepe) -- (quekenberg);
\draw[mapMain] (knickberg) -- (nonnenstein);
\draw[mapMain] (knickberg) -- (osterberg-w);
\draw[mapMain] (knickberg) -- (wittekindstein);
\draw[mapMain] (koeterberg) -- (deister-ii);
\draw[mapMain] (koeterberg) -- (hils);
\draw[mapMain] (koeterberg) -- (osterwald);
\draw[mapMain] (krapendorf) -- (mordkuhlenberg);
\draw[mapMain] (krapendorf) -- (quekenberg);
\draw[mapMain] (krapendorf) -- (windberg);
\draw[mapMain] (langeoog) -- (dornum);
\draw[mapMain] (langeoog) -- (esens);
\draw[mapMain] (langeoog) -- (spiekeroog);
\draw[mapMain] (langwarden) -- (bremerlehe);
\draw[mapMain] (langwarden) -- (varel);
\draw[mapMain] (lauenburg) -- (bretze);
\draw[mapMain] (lauenburg) -- (lueneburg);
\draw[mapMain] (leer) -- (onstwedde);
\draw[mapMain] (leer) -- (windberg);
\draw[mapMain] (lichtenberg) -- (brocken);
\draw[mapMain] (lichtenberg) -- (festberg);
\draw[mapMain] (litberg) -- (elmhorst);
\draw[mapMain] (litberg) -- (wilsede);
\draw[mapMain] (lueneburg) -- (bretze);
\draw[mapMain] (lueneburg) -- (nindorf);
\draw[mapMain] (lueneburg) -- (taetendorf);
\draw[mapMain] (meridianzeichen) -- (goettingen);
\draw[mapMain] (meridianzeichen) -- (hohehagen);
\draw[mapMain] (mordkuhlenberg) -- (doerenberg);
\draw[mapMain] (mordkuhlenberg) -- (knickberg);
\draw[mapMain] (mordkuhlenberg) -- (nonnenstein);
\draw[mapMain] (neuwerk) -- (langwarden);
\draw[mapMain] (nindorf) -- (timpenberg);
\draw[mapMain] (nonnenstein) -- (wittekindstein);
\draw[mapMain] (norderney) -- (baltrum);
\draw[mapMain] (norderney) -- (dornum);
\draw[mapMain] (norderney) -- (hage);
\draw[mapMain] (osterberg-w) -- (deister-ii);
\draw[mapMain] (osterberg-w) -- (wittekindstein);
\draw[mapMain] (osterwald) -- (hils);
\draw[mapMain] (pugelatz) -- (hankensbuettel);
\draw[mapMain] (quekenberg) -- (bentheim);
\draw[mapMain] (quekenberg) -- (doerenberg);
\draw[mapMain] (quekenberg) -- (mordkuhlenberg);
\draw[mapMain] (scharnhorst) -- (garssen);
\draw[mapMain] (silberberg) -- (basdahl);
\draw[mapMain] (silberberg) -- (crempe);
\draw[mapMain] (silberberg) -- (stade);
\draw[mapMain] (spiekeroog) -- (dornum);
\draw[mapMain] (spiekeroog) -- (esens);
\draw[mapMain] (spiekeroog) -- (jever);
\draw[mapMain] (stade) -- (hamburg);
\draw[mapMain] (stade) -- (litberg);
\draw[mapMain] (steinberg-n) -- (asendorf);
\draw[mapMain] (steinberg-n) -- (wilsede);
\draw[mapMain] (taetendorf) -- (hohenmechtin);
\draw[mapMain] (taetendorf) -- (holxerberg);
\draw[mapMain] (taetendorf) -- (pugelatz);
\draw[mapMain] (taetendorf) -- (wieren);
\draw[mapMain] (timpenberg) -- (wulfsode);
\draw[mapMain] (twistringen) -- (knickberg);
\draw[mapMain] (twistringen) -- (mordkuhlenberg);
\draw[mapMain] (varel) -- (garlste);
\draw[mapMain] (varel) -- (westerstede);
\draw[mapMain] (westerstede) -- (aurich);
\draw[mapMain] (westerstede) -- (krapendorf);
\draw[mapMain] (westerstede) -- (leer);
\draw[mapMain] (westerstede) -- (windberg);
\draw[mapMain] (wieren) -- (hankensbuettel);
\draw[mapMain] (wieren) -- (pugelatz);
\draw[mapMain] (wilsede) -- (breithorn);
\draw[mapMain] (wilsede) -- (bullerberg);
\draw[mapMain] (wilsede) -- (elmhorst);
\draw[mapMain] (wilsede) -- (falkenberg);
\draw[mapMain] (wilsede) -- (hauselberg);
\draw[mapMain] (wilsede) -- (lueneburg);
\draw[mapMain] (wilsede) -- (nindorf);
\draw[mapMain] (wilsede) -- (timpenberg);
\draw[mapMain] (wilsede) -- (wulfsode);
\draw[mapMain] (windberg) -- (kirchhesepe);
\draw[mapMain] (windberg) -- (onstwedde);
\draw[mapMain] (windberg) -- (quekenberg);
\draw[mapMain] (wittekindstein) -- (deister-ii);
\draw[mapMain] (wittekindstein) -- (koeterberg);
\draw[mapMain] (wohlenberg) -- (festberg);
\draw[mapMain] (wohlenberg) -- (lichtenberg);
\draw[mapMain] (wulfsode) -- (falkenberg);
\draw[mapMain] (wulfsode) -- (hauselberg);
\draw[mapMain] (zeven) -- (bottel);
\draw[mapMain] (zeven) -- (bruettendorf);
\draw[mapMain] (zeven) -- (litberg);
\draw[mapMain] (zeven) -- (wilsede);
\draw[mapFine] ([xshift=-.14mm]basis-n) -- ([xshift=-.14mm]basis-s);
\draw[mapFine] ([xshift=.14mm]basis-n) -- ([xshift=.14mm]basis-s);
\draw[mapMain] (hohehagen) -- (1935,2440);
\draw[mapMain] (brocken) -- (2236,2440);

% Survey marks and their witnessed labels.
\HannoverCircle{(neuwerk)}{6.1}
\HannoverLight{(neuwerk)}
\node[mapMinor,anchor=west] at (1192,148) {Neuwerk};
\HannoverCircle{(wangeroog)}{6.1}
\HannoverLight{(wangeroog)}
\node[mapMinor,anchor=west] at (901,231) {Wangeroog};
\HannoverCircle{(altenwalde)}{6.1}
\node[mapTown,anchor=east] at (1318,242) {Altenwalde};
\HannoverCircle{(borkum)}{6.1}
\HannoverLight{(borkum)}
\node[mapTown,anchor=east] at (248,411) {Borkum};
\HannoverCircle{(juist)}{6.1}
\node[mapTown,anchor=east] at (428,320) {Juist};
\HannoverCircle{(norderney)}{6.1}
\node[mapTown,anchor=east] at (541,289) {Norderney};
\HannoverCircle{(baltrum)}{6.1}
\node[mapTown,anchor=center] at (612,281) {Baltrum};
\HannoverCircle{(langeoog)}{6.1}
\node[mapTown,anchor=center] at (698,265) {Langeoog};
\HannoverCircle{(spiekeroog)}{6.1}
\node[mapTown,anchor=center] at (784,245) {Spiekeroog};
\HannoverCircle{(dornum)}{6.1}
\HannoverCross{(dornum)}
\node[mapTown,anchor=east] at (644,367) {Dornum};
\HannoverCircle{(esens)}{6.1}
\HannoverCross{(esens)}
\node[mapTown,anchor=west] at (776,372) {Esens};
\HannoverCircle{(hage)}{6.1}
\HannoverCross{(hage)}
\node[mapTown,anchor=east] at (571,410) {Hage};
\HannoverCircle{(pilsum)}{6.1}
\node[mapTown,anchor=west] at (470,528) {Pilsum};
\HannoverCircle{(jever)}{6.1}
\HannoverCastle{(jever)}
\node[mapTown,anchor=west] at (918,454) {Jever};
\HannoverCircle{(langwarden)}{6.1}
\HannoverCross{(langwarden)}
\node[mapTown,anchor=east] at (1096,409) {Langwarden};
\HannoverCircle{(haye)}{6.1}
\node[mapTown,anchor=east] at (1281,424) {Haye};
\HannoverCircle{(bremerlehe)}{6.1}
\HannoverCross{(bremerlehe)}
\node[mapTown,anchor=west] at (1267,461) {Bremerlehe};
\HannoverCircle{(holssel)}{4}
\node[mapMinor,anchor=east] at (1258,355) {Holssel};
\HannoverCircle{(krempel)}{4}
\node[mapMinor,anchor=west] at (1324,321) {Krempel};
\HannoverCircle{(bederkesa-i)}{4}
\node[mapMinor,anchor=center] at (1374,401) {I};
\HannoverCircle{(bederkesa)}{4}
\node[mapMinor,anchor=west] at (1400,411) {Bederkesa};
\HannoverCircle{(wuestenwohlde)}{4}
\node[mapMinor,anchor=center] at (1448,477) {Wüstenwohlde};
\HannoverCircle{(dolosenberg)}{4}
\node[mapMinor,anchor=west] at (1523,427) {Dolosenberg};
\HannoverCircle{(silberberg)}{6.1}
\node[mapTown,anchor=west] at (1511,330) {Silberberg};
\HannoverCircle{(basdahl)}{6.1}
\node[mapTown,anchor=west] at (1483,576) {Basdahl};
\HannoverCircle{(brillit)}{6.1}
\node[mapTown,anchor=east] at (1447,599) {Brillit};
\HannoverCircle{(varel)}{6.1}
\HannoverCross{(varel)}
\node[mapTown,anchor=east] at (993,601) {Varel};
\HannoverCircle{(garlste)}{6.1}
\node[mapTown,anchor=east] at (1292,729) {Garlste};
\HannoverCircle{(westerstede)}{6.1}
\HannoverCross{(westerstede)}
\node[mapTown,anchor=west] at (918,716) {Westerstede};
\HannoverCircle{(aurich)}{6.1}
\HannoverCross{(aurich)}
\node[mapTown,anchor=west] at (723,533) {Aurich};
\HannoverCircle{(emden)}{6.1}
\node[mapTown,anchor=west] at (561,610) {Emden};
\HannoverCircle{(leer)}{6.1}
\HannoverCross{(leer)}
\node[mapTown,anchor=west] at (678,749) {Leer};
\HannoverCircle{(bremen)}{6.1}
\HannoverCross{(bremen)}
\node[mapTown,anchor=west] at (1377,880) {Bremen};
\HannoverCircle{(zeven)}{6.1}
\HannoverCross{(zeven)}
\node[mapTown,anchor=center] at (1623,681) {Zeven};
\HannoverCircle{(bruettendorf)}{6.1}
\node[mapTown,anchor=east] at (1585,733) {Brüttendorf};
\HannoverCircle{(bottel)}{6.1}
\node[mapTown,anchor=east] at (1597,903) {Bottel};
\HannoverCircle{(steinberg-n)}{6.1}
\node[mapTown,anchor=west] at (1612,985) {Steinberg};
\HannoverCircle{(asendorf)}{6.1}
\HannoverCross{(asendorf)}
\node[mapTown,anchor=west] at (1483,1162) {Asendorf};
\HannoverCircle{(twistringen)}{6.1}
\HannoverCross{(twistringen)}
\node[mapTown,anchor=west] at (1277,1117) {Twistringen};
\HannoverCircle{(wildeshausen)}{6.1}
\HannoverCross{(wildeshausen)}
\node[mapTown,anchor=east] at (1148,1028) {Wildeshausen};
\HannoverCircle{(oldenburg)}{6.1}
\HannoverCastle{(oldenburg)}
\node[mapTown,anchor=west] at (1069,817) {Oldenburg};
\HannoverCircle{(krapendorf)}{6.1}
\HannoverCross{(krapendorf)}
\node[mapTown,anchor=west] at (962,1093) {Krapendorf};
\HannoverCircle{(windberg)}{6.1}
\node[mapTown,anchor=west] at (713,1031) {Windberg};
\HannoverCircle{(kirchhesepe)}{6.1}
\HannoverCross{(kirchhesepe)}
\node[mapTown,anchor=west] at (534,1250) {Kirchhesepe};
\HannoverCircle{(quekenberg)}{6.1}
\node[mapTown,anchor=east] at (764,1343) {Quekenberg};
\HannoverCircle{(mordkuhlenberg)}{6.1}
\node[mapTown,anchor=east] at (1020,1351) {Mordkuhlenberg};
\HannoverCircle{(knickberg)}{6.1}
\node[mapTown,anchor=east] at (1358,1312) {Knickberg};
\HannoverCircle{(osterberg-w)}{6.1}
\node[mapTown,anchor=center] at (1582,1271) {Osterberg};
\HannoverCircle{(eckberg)}{4}
\node[mapMinor,anchor=east] at (1638,1360) {Eckberg};
\HannoverCircle{(hornhuizen)}{6.1}
\HannoverCross{(hornhuizen)}
\node[mapTown,anchor=west] at (68,584) {Hornhuizen};
\HannoverCircle{(uithuiser)}{6.1}
\HannoverCross{(uithuiser)}
\node[mapTown,anchor=east] at (266,549) {Uithuiser\\Medem};
\HannoverCircle{(holwierda)}{6.1}
\HannoverCross{(holwierda)}
\node[mapTown,anchor=east] at (349,625) {Holwierda};
\HannoverCircle{(groningen)}{6.1}
\HannoverCross{(groningen)}
\node[mapTown,anchor=west] at (118,739) {Groningen};
\HannoverCircle{(midwolde)}{6.1}
\HannoverCross{(midwolde)}
\node[mapTown,anchor=east] at (413,761) {Midwolde};
\HannoverCircle{(onstwedde)}{6.1}
\HannoverCross{(onstwedde)}
\node[mapTown,anchor=east] at (449,916) {Onstwedde};
\HannoverCircle{(osterwolde)}{6.1}
\HannoverCross{(osterwolde)}
\node[mapTown,anchor=west] at (56,892) {Osterwolde};
\HannoverCircle{(rolde)}{6.1}
\HannoverCross{(rolde)}
\node[mapTown,anchor=east] at (228,935) {Rolde};
\HannoverCircle{(gieten)}{6.1}
\HannoverCross{(gieten)}
\node[mapTown,anchor=west] at (293,882) {Gieten};
\HannoverCircle{(beilen)}{6.1}
\HannoverCross{(beilen)}
\node[mapTown,anchor=west] at (173,1031) {Beilen};
\HannoverCircle{(ter-apel)}{6.1}
\HannoverCross{(ter-apel)}
\node[mapTown,anchor=east] at (431,1052) {Ter Apel};
\HannoverCircle{(sleen)}{6.1}
\HannoverCross{(sleen)}
\node[mapTown,anchor=west] at (310,1118) {Sleen};
\HannoverCircle{(koeverden)}{6.1}
\HannoverCross{(koeverden)}
\node[mapTown,anchor=east] at (242,1208) {Koeverden};
\HannoverCircle{(lemelerberg)}{6.1}
\node[mapTown,anchor=west] at (86,1381) {Lemelerberg};
\HannoverCircle{(uelsen)}{6.1}
\HannoverCross{(uelsen)}
\node[mapTown,anchor=west] at (362,1347) {Uelsen};
\HannoverCircle{(crempe)}{6.1}
\HannoverCross{(crempe)}
\node[mapTown,anchor=west] at (1743,234) {Crempe};
\HannoverCircle{(stade)}{6.1}
\HannoverCross{(stade)}
\node[mapTown,anchor=west] at (1706,466) {Stade};
\HannoverCircle{(hamburg)}{6.1}
\HannoverCross{(hamburg)}
\node[mapTown,anchor=east] at (1961,458) {Hamburg};
\HannoverCircle{(segeberg)}{6.1}
\node[mapTown,anchor=east] at (2127,156) {Segeberg};
\HannoverCircle{(luebeck)}{6.1}
\HannoverCross{(luebeck)}
\node[mapTown,anchor=west] at (2353,227) {Lübeck};
\HannoverCircle{(syk)}{6.1}
\HannoverCross{(syk)}
\node[mapTown,anchor=west] at (2143,428) {Syk};
\HannoverCircle{(basis-n)}{6.1}
\HannoverCircle{(basis-s)}{6.1}
\HannoverCircle{(hohenhorn)}{6.1}
\HannoverCross{(hohenhorn)}
\node[mapTown,anchor=west] at (2177,552) {Hohenhorn};
\HannoverCircle{(lauenburg)}{6.1}
\node[mapTown,anchor=west] at (2255,623) {Lauenburg};
\HannoverCircle{(litberg)}{6.1}
\node[mapTown,anchor=west] at (1808,630) {Litberg};
\HannoverCircle{(wilsede)}{6.1}
\node[mapTown,anchor=center] at (1945,799) {Wilsede};
\HannoverCircle{(lueneburg)}{6.1}
\HannoverCross{(lueneburg)}
\node[mapTown,anchor=west] at (2203,763) {Lüneburg};
\HannoverCircle{(bretze)}{6.1}
\node[mapTown,anchor=west] at (2323,719) {Bretze};
\HannoverCircle{(glienitz)}{6.1}
\node[mapTown,anchor=west] at (2443,764) {Glienitz};
\HannoverCircle{(hoe-beck)}{6.1}
\node[mapTown,anchor=west] at (2703,881) {Höbeck};
\HannoverCircle{(taetendorf)}{6.1}
\node[mapTown,anchor=west] at (2275,940) {Taetendorf};
\HannoverCircle{(hohenmechtin)}{6.1}
\node[mapTown,anchor=west] at (2459,946) {Hohenmechtin};
\HannoverCircle{(nindorf)}{6.1}
\node[mapTown,anchor=west] at (2142,836) {Nindorf};
\HannoverCircle{(timpenberg)}{6.1}
\node[mapTown,anchor=west] at (2128,867) {Timpenberg};
\HannoverCircle{(wulfsode)}{6.1}
\node[mapTown,anchor=east] at (2088,913) {Wulfsode};
\HannoverCircle{(elmhorst)}{6.1}
\node[mapTown,anchor=west] at (1800,972) {Elmhorst};
\HannoverCircle{(falkenberg)}{6.1}
\node[mapTown,anchor=east] at (1888,1101) {Falkenberg};
\HannoverCircle{(hauselberg)}{6.1}
\node[mapTown,anchor=east] at (2060,1072) {Hauselberg};
\HannoverCircle{(breithorn)}{6.1}
\node[mapTown,anchor=west] at (2129,1130) {Breithorn};
\HannoverCircle{(scharnhorst)}{6.1}
\node[mapTown,anchor=west] at (2125,1211) {Scharnhorst};
\HannoverCircle{(garssen)}{6.1}
\node[mapTown,anchor=west] at (2063,1272) {Garssen};
\HannoverCircle{(celle)}{6.1}
\HannoverCross{(celle)}
\node[mapTown,anchor=east] at (2006,1291) {Celle};
\HannoverCircle{(holxerberg)}{6.1}
\node[mapTown,anchor=east] at (2216,1050) {Holxerberg};
\HannoverCircle{(wieren)}{6.1}
\node[mapTown,anchor=east] at (2294,1100) {Wieren};
\HannoverCircle{(pugelatz)}{6.1}
\node[mapTown,anchor=west] at (2393,1118) {Pugelatz};
\HannoverCircle{(hankensbuettel)}{6.1}
\node[mapTown,anchor=west] at (2272,1227) {Hankensbüttel};
\HannoverCircle{(bullerberg)}{6.1}
\node[mapTown,anchor=center] at (1694,797) {Bullerberg};
\HannoverCircle{(brelingerberg)}{4}
\node[mapMinor,anchor=east] at (1812,1324) {Brelingerberg};
\HannoverCircle{(osterberg-e)}{4}
\node[mapMinor,anchor=east] at (1915,1398) {Osterberg};
\HannoverCircle{(oldenzaal)}{6.1}
\HannoverCross{(oldenzaal)}
\node[mapTown,anchor=east] at (341,1501) {Oldenzaal};
\HannoverCircle{(bentheim)}{6.1}
\HannoverCastle{(bentheim)}
\node[mapTown,anchor=west] at (489,1531) {Bentheim};
\HannoverCircle{(harikerberg)}{6.1}
\node[mapTown,anchor=west] at (157,1578) {Harikerberg};
\HannoverCircle{(groll)}{6.1}
\HannoverCross{(groll)}
\node[mapTown,anchor=west] at (173,1751) {Groll};
\HannoverCircle{(ahaus)}{6.1}
\HannoverCross{(ahaus)}
\node[mapTown,anchor=west] at (390,1712) {Ahaus};
\HannoverCircle{(bocholt)}{6.1}
\HannoverCross{(bocholt)}
\node[mapTown,anchor=west] at (188,1908) {Bocholt};
\HannoverCircle{(doerenberg)}{6.1}
\node[mapTown,anchor=west] at (941,1680) {Dörenberg};
\HannoverCircle{(nonnenstein)}{6.1}
\node[mapTown,anchor=west] at (1171,1580) {Nonnenstein};
\HannoverCircle{(wittekindstein)}{6.1}
\HannoverTower{(wittekindstein)}
\node[mapTown,anchor=east] at (1489,1634) {Wittekindstein};
\HannoverCircle{(huenenburg)}{6.1}
\node[mapTown,anchor=east] at (1143,1815) {Hünenburg};
\HannoverCircle{(koeterberg)}{6.1}
\node[mapTown,anchor=east] at (1585,1972) {Köterberg};
\HannoverCircle{(hausheide)}{6.1}
\node[mapTown,anchor=east] at (1408,2065) {Hausheide};
\HannoverCircle{(spitzwarte)}{6.1}
\node[mapTown,anchor=west] at (1164,2239) {Spitzwarte};
\HannoverCircle{(desenberg)}{6.1}
\node[mapTown,anchor=east] at (1538,2241) {Desenberg};
\HannoverCircle{(deisselberg)}{6.1}
\node[mapTown,anchor=west] at (1620,2177) {Deisselberg};
\HannoverCircle{(dommel)}{6.1}
\node[mapTown,anchor=east] at (1221,2403) {Dommel};
\HannoverCircle{(hannover)}{4}
\HannoverCross{(hannover)}
\node[mapMinor,anchor=west] at (1841,1513) {Hannover};
\HannoverCircle{(deister-ii)}{6.1}
\node[mapTown,anchor=east] at (1722,1604) {Deister \(^{\mathrm{II}}\)};
\HannoverCircle{(deister-i)}{6.1}
\node[mapTown,anchor=west] at (1768,1628) {\(^{\mathrm I}\) Deister};
\HannoverCircle{(osterwald)}{6.1}
\node[mapTown,anchor=west] at (1774,1730) {Osterwald};
\HannoverCircle{(hils)}{6.1}
\node[mapTown,anchor=east] at (1854,1921) {Hils};
\HannoverCircle{(wohlenberg)}{6.1}
\node[mapTown,anchor=west] at (2194,1449) {Wohlenberg};
\HannoverCircle{(lichtenberg)}{6.1}
\node[mapTown,anchor=west] at (2131,1750) {Lichtenberg};
\HannoverCircle{(festberg)}{6.1}
\node[mapTown,anchor=west] at (2321,1690) {Festberg};
\HannoverCircle{(huyseburg)}{6.1}
\HannoverCross{(huyseburg)}
\node[mapTown,anchor=west] at (2529,1865) {Huyseburg};
\HannoverCircle{(brocken)}{6.1}
\HannoverTower{(brocken)}
\node[mapTown,anchor=west] at (2305,1990) {Brocken};
\HannoverCircle{(bocksberg)}{4}
\node[mapMinor,anchor=center] at (2131,1947) {Bocksberg};
\HannoverCircle{(fuhrenberg)}{4}
\node[mapMinor,anchor=east] at (2033,1996) {Fuhrenberg};
\HannoverCircle{(haspelkopf)}{4}
\node[mapMinor,anchor=east] at (2143,2074) {Haspelkopf};
\HannoverCircle{(eversberg)}{4}
\node[mapMinor,anchor=west] at (2305,2101) {Eversberg};
\HannoverCircle{(lange-ecke)}{4}
\node[mapMinor,anchor=east] at (2255,2140) {Lange Ecke};
\HannoverCircle{(schalliethe)}{4}
\node[mapMinor,anchor=west] at (2473,2144) {Schalliethe};
\HannoverCircle{(tockenburg)}{4}
\node[mapMinor,anchor=center] at (1965,2148) {Tockenberg};
\HannoverCircle{(wulften)}{4}
\node[mapMinor,anchor=east] at (2036,2104) {Wulften};
\HannoverCircle{(hauseberg)}{4}
\node[mapMinor,anchor=west] at (2012,2191) {Hauseberg};
\HannoverCircle{(hellberg)}{4}
\node[mapMinor,anchor=west] at (2104,2213) {Hellberg};
\HannoverCircle{(bazenberg)}{4}
\node[mapMinor,anchor=center] at (2097,2257) {Batzenberg};
\HannoverCircle{(sonnenstein)}{4}
\node[mapMinor,anchor=west] at (2176,2275) {Sonnenstein};
\HannoverCircle{(fahlenstollen)}{4}
\node[mapMinor,anchor=center] at (1767,2090) {Fahlenstollen};
\HannoverCircle{(weper)}{4}
\node[mapMinor,anchor=west] at (1857,2111) {Weper};
\HannoverCircle{(meridianzeichen)}{6.1}
\node[mapTown,anchor=east] at (1997,2188) {Meridianzeichen};
\HannoverCircle{(kleper)}{4}
\node[mapMinor,anchor=west] at (1957,2237) {Kleper};
\HannoverCircle{(goettingen)}{6.1}
\node[mapTown,font=\fontsize{7.3}{8}\selectfont\bfseries,anchor=west] at (1890,2265) {GÖTTINGEN\\{\fontsize{4.6}{5.1}\selectfont STERNWARTE}};
\HannoverCircle{(hohehagen)}{6.1}
\node[mapTown,anchor=west] at (1762,2315) {Hohehagen};
\HannoverCircle{(gieseberg)}{4}
\node[mapMinor,anchor=west] at (1900,2333) {Gieseberg};
\HannoverCircle{(steinberg-s)}{4}
\node[mapMinor,anchor=center] at (1762,2395) {Steinberg};
\HannoverCircle{(meisner)}{6.1}
\node[mapTown,anchor=center] at (1857,2483) {Meisner};
\HannoverCircle{(staufenberg)}{6.1}
\node[mapTown,anchor=east] at (1758,2281) {Staufenberg};
\node[mapTick] at (1388,434) {II};
\node[mapTick,rotate=-51] at (2096,444) {Basis};
\node[mapMinor,rotate=78] at (2250,2360) {Inselsberg};
\draw[mapFine] (1951,130) -- (1951,420);
\node[mapTick,rotate=90] at (1943,294) {Meridian der Göttinger Sternwarte};

% Graduated frame. Openings accommodate the marginal survey marks.
\draw[line width=.75pt] (45,868) -- (45,62) -- (3050,68) -- (3044,2500)
  -- (1910,2498);
\draw[line width=.75pt] (1847,2498) -- (42,2492) -- (45,941);
\draw[mapFine] (55,852) -- (55,72) -- (3040,78) -- (3034,2489) -- (1907,2488);
\draw[mapFine] (1850,2488) -- (53,2482) -- (55,951);
\draw[mapFine] (90,106) -- (850,106) -- (1450,108) -- (1576,109) -- (1880,113) -- (1987,113) -- (2150,110) -- (2411,114) -- (2663,111) -- (2920,112) -- (3007,113);
\draw[mapFine] (99,115) -- (850,115) -- (1450,117) -- (1576,118) -- (1880,122) -- (1987,122) -- (2150,119) -- (2411,123) -- (2663,120) -- (2920,121) -- (2998,121.9);
\draw[mapFine] (109,125) -- (850,125) -- (1450,127) -- (1576,128) -- (1880,132) -- (1987,132) -- (2150,129) -- (2411,133) -- (2663,130) -- (2920,131) -- (2988,131.8);
\draw[mapFine] (84,2454) -- (600,2455) -- (1010,2457) -- (1417,2459) -- (1507,2458) -- (1575,2456) -- (1687,2459) -- (1777,2459);
\draw[mapFine] (1940,2457) -- (2180,2458) -- (2410,2459) -- (2900,2459) -- (2999,2459);
\draw[mapFine] (93,2445) -- (600,2446) -- (1010,2448) -- (1417,2450) -- (1507,2449) -- (1575,2447) -- (1687,2450) -- (1777,2450) -- (1778,2450);
\draw[mapFine] (1939,2448) -- (2180,2449) -- (2410,2450) -- (2900,2450) -- (2990,2450);
\draw[mapFine] (103,2435) -- (600,2436) -- (1010,2438) -- (1417,2440) -- (1507,2439) -- (1575,2437) -- (1687,2440) -- (1757,2440);
\draw[mapFine] (1939,2438) -- (2180,2439) -- (2410,2440) -- (2900,2440) -- (2980,2440);
\draw[mapFine] (89.5,106) -- (89,545);
\draw[mapFine] (88.9,594) -- (88.5,826);
\draw[mapFine] (88.1,974) -- (87,1328);
\draw[mapFine] (86.7,1443) -- (84.3,2454);
\draw[mapFine] (3007,113) -- (3006,610) -- (3004.6,960) -- (3003.4,1390) -- (3001.9,1800) -- (2999.6,2459);
\draw[mapFine] (98.5,115) -- (98,545);
\draw[mapFine] (97.9,595) -- (97.5,817);
\draw[mapFine] (97,985) -- (96,1325);
\draw[mapFine] (95.6,1477) -- (93.3,2445);
\draw[mapFine] (2998,122) -- (2997,610) -- (2995.6,960) -- (2994.4,1390) -- (2992.9,1800) -- (2990.6,2450);
\draw[mapFine] (108.5,125) -- (108,545);
\draw[mapFine] (107.9,595) -- (107.5,805);
\draw[mapFine] (107,993) -- (106,1317);
\draw[mapFine] (105.6,1498) -- (103.3,2435);
\draw[mapFine] (2988,132) -- (2987,610) -- (2985.6,960) -- (2984.4,1390) -- (2982.9,1800) -- (2980.7,2440);
\draw[mapFine] (90,106) -- (109,125);
\draw[mapFine] (3007,113) -- (2988,131.8);
\draw[mapFine] (84,2454) -- (103,2435);
\draw[mapFine] (2999,2459) -- (2980,2440);
\draw[mapFine] (193,110.5) -- (277,110.5);
\draw[mapFine] (363,110.5) -- (447,110.5);
\draw[mapFine] (533,110.5) -- (617,110.5);
\draw[mapFine] (703,110.5) -- (790,110.5);
\draw[mapFine] (875,110.6) -- (960,110.9);
\draw[mapFine] (1045,111.2) -- (1131,111.4);
\draw[mapFine] (1217,111.7) -- (1303,112);
\draw[mapFine] (1390,112.3) -- (1450,112.5) -- (1475,112.7);
\draw[mapFine] (1560,113.4) -- (1576,113.5) -- (1645,114.4);
\draw[mapFine] (1730,115.5) -- (1816,116.7);
\draw[mapFine] (1898,117.5) -- (1987,117.5);
\draw[mapFine] (2073,115.9) -- (2150,114.5);
\draw[mapFine] (2237,115.8) -- (2324,117.2);
\draw[mapFine] (2411,118.5) -- (2496,117.5);
\draw[mapFine] (2581,116.5) -- (2663,115.5);
\draw[mapFine] (2751,115.8) -- (2836,116.2);
\draw[mapFine] (2920,116.5) -- (3003,117.5);
\draw[mapFine] (90,2449.5) -- (157,2449.6);
\draw[mapFine] (246,2449.8) -- (336,2450);
\draw[mapFine] (426,2450.2) -- (515,2450.3);
\draw[mapFine] (604,2450.5) -- (696,2451);
\draw[mapFine] (788,2451.4) -- (877,2451.9);
\draw[mapFine] (967,2452.3) -- (1010,2452.5) -- (1056,2452.7);
\draw[mapFine] (1147,2453.2) -- (1237,2453.6);
\draw[mapFine] (1327,2454.1) -- (1417,2454.5);
\draw[mapFine] (1507,2453.5) -- (1575,2451.5) -- (1598,2452.1);
\draw[mapFine] (1687,2454.5) -- (1777,2454.5);
\draw[mapFine] (1940,2452.5) -- (1957,2452.6);
\draw[mapFine] (2047,2452.9) -- (2136,2453.3);
\draw[mapFine] (2227,2453.7) -- (2317,2454.1);
\draw[mapFine] (2409,2454.5) -- (2410,2454.5) -- (2497,2454.5);
\draw[mapFine] (2587,2454.5) -- (2686,2454.5);
\draw[mapFine] (2776,2454.5) -- (2866,2454.5);
\draw[mapFine] (2960,2454.5) -- (2999,2454.5);
\draw[mapFine] (93.9,184) -- (93.7,332);
\draw[mapFine] (93.6,471) -- (93.5,545);
\draw[mapFine] (93.4,595) -- (93.4,610) -- (93.4,618);
\draw[mapFine] (93.1,760) -- (93,816);
\draw[mapFine] (92.3,1051) -- (91.9,1198);
\draw[mapFine] (90.8,1629) -- (90.4,1777);
\draw[mapFine] (90.1,1921) -- (89.7,2069);
\draw[mapFine] (89.4,2212) -- (89,2360);
\draw[mapFine] (3002.2,238) -- (3002,383);
\draw[mapFine] (3001.7,526) -- (3001.5,610) -- (3001.3,670);
\draw[mapFine] (3000.7,812) -- (3000.1,959);
\draw[mapFine] (2999.7,1101) -- (2999.3,1248);
\draw[mapFine] (2998.9,1391) -- (2998.4,1538);
\draw[mapFine] (2997.8,1679) -- (2997.4,1800) -- (2997.3,1827);
\draw[mapFine] (2996.8,1969) -- (2996.3,2117);
\draw[mapFine] (2995.8,2259) -- (2995.3,2407);
\draw[mapFine] (363,106) -- (363,125);
\node[mapTick] at (363,91) {\(30'\)};
\draw[mapFine] (617,106) -- (617,125);
\node[mapTick] at (617,91) {\(25^0\)};
\draw[mapFine] (875,106.1) -- (875,125.1);
\node[mapTick] at (875,91.0833) {\(30'\)};
\draw[mapFine] (1131,106.9) -- (1131,125.9);
\node[mapTick] at (1131,91.9367) {\(26^0\)};
\draw[mapFine] (1390,107.8) -- (1390,126.8);
\node[mapTick] at (1390,92.8) {\(30'\)};
\draw[mapFine] (1645,109.9) -- (1645,128.9);
\node[mapTick] at (1645,94.9079) {\(27^0\)};
\draw[mapFine] (1898,113) -- (1898,132);
\node[mapTick] at (1898,98) {\(30'\)};
\draw[mapFine] (2150,110) -- (2150,129);
\node[mapTick] at (2150,95) {\(28^0\)};
\draw[mapFine] (2411,114) -- (2411,133);
\node[mapTick] at (2411,99) {\(30'\)};
\draw[mapFine] (2663,111) -- (2663,130);
\node[mapTick] at (2663,96) {\(29^0\)};
\draw[mapFine] (2920,112) -- (2920,131);
\node[mapTick] at (2920,97) {\(30'\)};
\draw[mapFine] (193,106) -- (193,115);
\draw[mapFine] (277,106) -- (277,115);
\draw[mapFine] (447,106) -- (447,115);
\draw[mapFine] (533,106) -- (533,115);
\draw[mapFine] (703,106) -- (703,115);
\draw[mapFine] (790,106) -- (790,115);
\draw[mapFine] (960,106.4) -- (960,115.4);
\draw[mapFine] (1045,106.7) -- (1045,115.7);
\draw[mapFine] (1217,107.2) -- (1217,116.2);
\draw[mapFine] (1303,107.5) -- (1303,116.5);
\draw[mapFine] (1475,108.2) -- (1475,117.2);
\draw[mapFine] (1560,108.9) -- (1560,117.9);
\draw[mapFine] (1730,111) -- (1730,120);
\draw[mapFine] (1816,112.2) -- (1816,121.2);
\draw[mapFine] (1987,113) -- (1987,122);
\draw[mapFine] (2073,111.4) -- (2073,120.4);
\draw[mapFine] (2237,111.3) -- (2237,120.3);
\draw[mapFine] (2324,112.7) -- (2324,121.7);
\draw[mapFine] (2496,113) -- (2496,122);
\draw[mapFine] (2581,112) -- (2581,121);
\draw[mapFine] (2751,111.3) -- (2751,120.3);
\draw[mapFine] (2836,111.7) -- (2836,120.7);
\draw[mapFine] (3003,113) -- (3003,122);
\draw[mapFine] (246,2454.3) -- (246,2435.3);
\node[mapTick] at (246,2474.31) {\(30'\)};
\draw[mapFine] (515,2454.8) -- (515,2435.8);
\node[mapTick] at (515,2474.84) {\(25^0\)};
\draw[mapFine] (788,2455.9) -- (788,2436.9);
\node[mapTick] at (788,2475.92) {\(30'\)};
\draw[mapFine] (1056,2457.2) -- (1056,2438.2);
\node[mapTick] at (1056,2477.23) {\(26^0\)};
\draw[mapFine] (1327,2458.6) -- (1327,2439.6);
\node[mapTick] at (1327,2478.56) {\(30'\)};
\draw[mapFine] (1598,2456.6) -- (1598,2437.6);
\node[mapTick] at (1598,2476.62) {\(27^0\)};
\draw[mapFine] (2136,2457.8) -- (2136,2438.8);
\node[mapTick] at (2136,2477.82) {\(28^0\)};
\draw[mapFine] (2409,2459) -- (2409,2440);
\node[mapTick] at (2409,2479) {\(30'\)};
\draw[mapFine] (2686,2459) -- (2686,2440);
\node[mapTick] at (2686,2479) {\(29^0\)};
\draw[mapFine] (2960,2459) -- (2960,2440);
\node[mapTick] at (2960,2479) {\(30'\)};
\draw[mapFine] (157,2454.1) -- (157,2445.1);
\draw[mapFine] (336,2454.5) -- (336,2445.5);
\draw[mapFine] (426,2454.7) -- (426,2445.7);
\draw[mapFine] (604,2455) -- (604,2446);
\draw[mapFine] (696,2455.5) -- (696,2446.5);
\draw[mapFine] (877,2456.4) -- (877,2447.4);
\draw[mapFine] (967,2456.8) -- (967,2447.8);
\draw[mapFine] (1147,2457.7) -- (1147,2448.7);
\draw[mapFine] (1237,2458.1) -- (1237,2449.1);
\draw[mapFine] (1417,2459) -- (1417,2450);
\draw[mapFine] (1507,2458) -- (1507,2449);
\draw[mapFine] (1687,2459) -- (1687,2450);
\draw[mapFine] (1777,2459) -- (1777,2450);
\draw[mapFine] (1957,2457.1) -- (1957,2448.1);
\draw[mapFine] (2047,2457.4) -- (2047,2448.4);
\draw[mapFine] (2227,2458.2) -- (2227,2449.2);
\draw[mapFine] (2317,2458.6) -- (2317,2449.6);
\draw[mapFine] (2497,2459) -- (2497,2450);
\draw[mapFine] (2587,2459) -- (2587,2450);
\draw[mapFine] (2776,2459) -- (2776,2450);
\draw[mapFine] (2866,2459) -- (2866,2450);
\draw[mapFine] (89.1,473) -- (108.1,473);
\node[mapTick] at (72,473) {\(30'\)};
\node[mapTick] at (72,1336) {\(30'\)};
\draw[mapFine] (85.9,1777) -- (104.9,1777);
\node[mapTick] at (72,1777) {\(52^0\)};
\draw[mapFine] (84.9,2210) -- (103.9,2210);
\node[mapTick] at (72,2210) {\(30'\)};
\draw[mapFine] (89.4,184) -- (98.4,184);
\draw[mapFine] (89.2,332) -- (98.2,332);
\draw[mapFine] (88.9,618) -- (97.9,618);
\draw[mapFine] (88.6,760) -- (97.6,760);
\draw[mapFine] (87.8,1051) -- (96.8,1051);
\draw[mapFine] (87.4,1198) -- (96.4,1198);
\draw[mapFine] (86.3,1629) -- (95.3,1629);
\draw[mapFine] (85.6,1921) -- (94.6,1921);
\draw[mapFine] (85.2,2069) -- (94.2,2069);
\draw[mapFine] (84.5,2360) -- (93.5,2360);
\draw[mapFine] (3006.2,526) -- (2987.2,526);
\node[mapTick] at (3024,526) {\(30'\)};
\draw[mapFine] (3004.6,959) -- (2985.6,959);
\node[mapTick] at (3024,959) {\(53^0\)};
\draw[mapFine] (3003.4,1393) -- (2984.4,1393);
\node[mapTick] at (3024,1393) {\(30'\)};
\draw[mapFine] (3001.8,1826) -- (2982.8,1826);
\node[mapTick] at (3024,1826) {\(52^0\)};
\draw[mapFine] (3000.3,2264) -- (2981.3,2264);
\node[mapTick] at (3024,2264) {\(30'\)};
\draw[mapFine] (3006.7,238) -- (2997.7,238);
\draw[mapFine] (3006.5,383) -- (2997.5,383);
\draw[mapFine] (3005.8,670) -- (2996.8,670);
\draw[mapFine] (3005.2,812) -- (2996.2,812);
\draw[mapFine] (3004.2,1101) -- (2995.2,1101);
\draw[mapFine] (3003.8,1248) -- (2994.8,1248);
\draw[mapFine] (3002.9,1538) -- (2993.9,1538);
\draw[mapFine] (3002.3,1679) -- (2993.3,1679);
\draw[mapFine] (3001.3,1969) -- (2992.3,1969);
\draw[mapFine] (3000.8,2117) -- (2991.8,2117);
\draw[mapFine] (2999.8,2407) -- (2990.8,2407);

% The signs key and its samples.
\node[mapCaption] at (2738,192) {Zeichenerklärung};
\draw[mapMain] (2508,254) -- (2616,254);
\node[mapLegend,anchor=west] at (2624,250) {Seite eines Hauptdreiecks.};
\draw[mapFine] (2508,303) -- (2616,303);
\node[mapLegend,anchor=west] at (2640,296) {„\qquad „\quad eingefügten, übrigens mit};
\node[mapLegend,anchor=west] at (2624,327) {nahezu gleicher Genauigkeit gemessenen};
\node[mapLegend,anchor=west] at (2624,356) {Dreiecks.};
\draw[mapDash] (2508,391) -- (2616,391);
\node[mapLegend,anchor=west] at (2624,385) {Seite eines Dreiecks der Systeme der Nach-};
\node[mapLegend,anchor=west] at (2618,431) {barländer.};
\coordinate (key-church) at (2564,477);
\coordinate (key-light) at (2564,523);
\coordinate (key-castle) at (2564,568);
\coordinate (key-tower) at (2564,614);
\coordinate (key-circle) at (2564,659);
\foreach \n in {key-church,key-light,key-castle,key-tower,key-circle}{\HannoverCircle{(\n)}{6}}
\HannoverCross{(key-church)}
\HannoverLight{(key-light)}
\HannoverCastle{(key-castle)}
\HannoverTower{(key-tower)}
\node[mapLegend,anchor=west] at (2624,477) {Kirchturm.};
\node[mapLegend,anchor=west] at (2624,523) {Leuchtturm.};
\node[mapLegend,anchor=west] at (2624,568) {Schlossturm.};
\node[mapLegend,anchor=west] at (2624,614) {Turm (Brocken, Wittekindstein)};
\node[mapLegend,anchor=west] at (2624,659) {Standpunkte, meistens durch Signale};
\node[mapLegend,anchor=west] at (2624,705) {bezeichnet.};

% Summary and explanatory text, with the original line divisions.
\node[mapCaption,font=\fontsize{9.8}{10.5}\selectfont\bfseries\itshape] at (685,1821)
  {Übersicht der gemessenen Dreieckssysteme.};
\node[mapLegend,anchor=east] at (480,1868) {I. Das Hauptsystem};
\node[mapLegend,anchor=east] at (425,1938) {II. Westphalen};
\node[mapLegend,anchor=east] at (419,1997) {III. Ostfriesland};
\node[mapLegend,anchor=east] at (371,2023) {IV. Bremen};
\node[mapLegend,anchor=east] at (433,2070) {V. Lüneburg etc.};
\node[mapLegend,anchor=east] at (424,2094) {VI. Mittelweser};
\node[mapLegend,anchor=east] at (399,2122) {VII. Eichsfeld};
\node[mapLegend,anchor=east] at (360,2149) {VIII. Harz};
\node[mapLegend,anchor=east] at (413,2175) {IX. Oberweser};
\node[mapLegend,anchor=east] at (356,2201) {X. Aller};
\foreach \x/\X/\y/\c in {
350/486/1885/HannoverHaupt,305/426/1955/HannoverWest,305/409/2010/HannoverOst,
305/396/2038/HannoverBremen,305/424/2083/HannoverLueneburg,305/426/2110/HannoverMittel,
305/399/2136/HannoverEichs,305/362/2163/HannoverHarz,305/414/2190/HannoverOber,305/360/2217/HannoverAller}{
  \draw[draw=\c,line width=.57mm] (\x,\y) -- (\X,\y);
}
\draw[draw=HannoverHaupt,line width=.13mm] (395,2137) -- (399,2137);
\draw[draw=HannoverOber!70!black,line width=.10pt] (306,2193) -- (310,2187);
\draw[draw=HannoverOber!70!black,line width=.10pt] (310,2193) -- (314,2187);
\draw[draw=HannoverOber!70!black,line width=.10pt] (314,2193) -- (318,2187);
\draw[draw=HannoverOber!70!black,line width=.10pt] (318,2193) -- (322,2187);
\draw[draw=HannoverOber!70!black,line width=.10pt] (322,2193) -- (326,2187);
\draw[draw=HannoverOber!70!black,line width=.10pt] (326,2193) -- (330,2187);
\draw[draw=HannoverOber!70!black,line width=.10pt] (330,2193) -- (334,2187);
\draw[draw=HannoverOber!70!black,line width=.10pt] (334,2193) -- (338,2187);
\draw[draw=HannoverOber!70!black,line width=.10pt] (338,2193) -- (342,2187);
\draw[draw=HannoverOber!70!black,line width=.10pt] (342,2193) -- (346,2187);
\draw[draw=HannoverOber!70!black,line width=.10pt] (346,2193) -- (350,2187);
\draw[draw=HannoverOber!70!black,line width=.10pt] (350,2193) -- (354,2187);
\draw[draw=HannoverOber!70!black,line width=.10pt] (354,2193) -- (358,2187);
\draw[draw=HannoverOber!70!black,line width=.10pt] (358,2193) -- (362,2187);
\draw[draw=HannoverOber!70!black,line width=.10pt] (362,2193) -- (366,2187);
\draw[draw=HannoverOber!70!black,line width=.10pt] (366,2193) -- (370,2187);
\draw[draw=HannoverOber!70!black,line width=.10pt] (370,2193) -- (374,2187);
\draw[draw=HannoverOber!70!black,line width=.10pt] (374,2193) -- (378,2187);
\draw[draw=HannoverOber!70!black,line width=.10pt] (378,2193) -- (382,2187);
\draw[draw=HannoverOber!70!black,line width=.10pt] (382,2193) -- (386,2187);
\draw[draw=HannoverOber!70!black,line width=.10pt] (386,2193) -- (390,2187);
\draw[draw=HannoverOber!70!black,line width=.10pt] (390,2193) -- (394,2187);
\draw[draw=HannoverOber!70!black,line width=.10pt] (394,2193) -- (398,2187);
\draw[draw=HannoverOber!70!black,line width=.10pt] (398,2193) -- (402,2187);
\draw[draw=HannoverOber!70!black,line width=.10pt] (402,2193) -- (406,2187);
\draw[draw=HannoverOber!70!black,line width=.10pt] (406,2193) -- (410,2187);
\draw[draw=HannoverOber!70!black,line width=.10pt] (410,2193) -- (414,2187);
\draw[draw=HannoverAller!70!black,line width=.10pt] (306,2220) -- (310,2214);
\draw[draw=HannoverAller!70!black,line width=.10pt] (310,2220) -- (314,2214);
\draw[draw=HannoverAller!70!black,line width=.10pt] (314,2220) -- (318,2214);
\draw[draw=HannoverAller!70!black,line width=.10pt] (318,2220) -- (322,2214);
\draw[draw=HannoverAller!70!black,line width=.10pt] (322,2220) -- (326,2214);
\draw[draw=HannoverAller!70!black,line width=.10pt] (326,2220) -- (330,2214);
\draw[draw=HannoverAller!70!black,line width=.10pt] (330,2220) -- (334,2214);
\draw[draw=HannoverAller!70!black,line width=.10pt] (334,2220) -- (338,2214);
\draw[draw=HannoverAller!70!black,line width=.10pt] (338,2220) -- (342,2214);
\draw[draw=HannoverAller!70!black,line width=.10pt] (342,2220) -- (346,2214);
\draw[draw=HannoverAller!70!black,line width=.10pt] (346,2220) -- (350,2214);
\draw[draw=HannoverAller!70!black,line width=.10pt] (350,2220) -- (354,2214);
\draw[draw=HannoverAller!70!black,line width=.10pt] (354,2220) -- (358,2214);
\node[mapLegend,anchor=west] at (491,1868)
  {enth. 81 Dreiecksseiten, gemessen 1821-25 von Hofrat Gauss. Die};
\node[mapLegend,anchor=west] at (491,1888)
  {Dreiecksreihe bis zu den Seiten Falkenberg-Wilsede-Hamburg ist};
\node[mapLegend,anchor=west] at (491,1909) {die Gradmessung.};
\node[mapLegend,anchor=west] at (489,1936)
  {enth. 41 Dreiecksseiten, gem. 1829 von Leutnant Gauss bis zu};
\node[mapLegend,anchor=west] at (489,1957)
  {den Seiten Mordkuhlenberg-Quekenberg-Kirchhesepe, vollend. nord-};
\node[mapLegend,anchor=west] at (489,1977) {wärts 1830 und 31 von Leutnant Hartmann.};
\node[mapLegend,anchor=west] at (488,1998)
  {enth. 21 Dreiecksseiten, gemessen 1841 von Hauptmann Müller.};
\node[mapLegend,anchor=west] at (488,2024)
  {enth. 36 Dreiecksseiten, gem. 1839 die westl. Hälfte von Hauptmann};
\node[mapLegend,anchor=west] at (488,2045)
  {Müller, 1843 u. 44 bis Basdahl u. Silberbg. von Leutnant Gauss.};
\node[mapLegend,anchor=west] at (488,2070)
  {enth. 25 Dreiecksseiten, gem. 1830 von Hptm. Müller u. Lt. Gauss.};
\node[mapLegend,anchor=west] at (487,2096)
  {enth. 12 Dreiecksseiten, gem. 1833 von\quad „\quad Müller „\quad „\quad Gauss.};
\node[mapLegend,anchor=west] at (488,2122)
  {enth. 14 Dreiecksseiten, gemessen 1828 von Hauptmann Müller.};
\node[mapLegend,anchor=west] at (487,2148)
  {enth. 16 Dreiecksseiten, gemessen 1833 von Leutnant Hartmann.};
\node[mapLegend,anchor=west] at (487,2174)
  {enth. 13 Dreiecksseiten, gemessen 1836 von Hauptmann Müller.};
\node[mapLegend,anchor=west] at (486,2201)
  {enth. 10 Dreiecksseiten, gemessen 1838 von Hauptmann Müller.};
\node[mapLegend,anchor=west] at (213,2233)
  {Das absolute Längenmaß für die Berechnung aller dieser Dreiecke ist der für die};
\node[mapLegend,anchor=west] at (176,2260)
  {dänische Triangulierung nahe bei Hamburg gemessenen Basis entnommen. Die obigen Dreiecke};
\node[mapLegend,anchor=west] at (176,2287)
  {enthalten 89 Standpunkte. Aus den auf diesen und etwa 100 weiteren Standpunkten 2. Klasse,};
\node[mapLegend,anchor=west] at (176,2313)
  {die keine Dreiecke bilden, gemachten Winkelmessungen sind die rechtwinkligen Koordinaten};
\node[mapLegend,anchor=west] at (176,2340)
  {von rund 3000 Kirchtürmen, Windmühlen und sonstigen Objekten berechnet und in der großen};
\node[mapLegend,anchor=west] at (176,2366)
  {Karte [des Papenschen Atlas] aufgetragen. Die Hofrat Gauss eigene Berechnungs-};
\node[mapLegend,anchor=west] at (176,2393)
  {art dieser Koordinaten, deren Achsen der Göttinger Meridian und eine Senkrechte darauf sind, ent-};
\node[mapLegend,anchor=west] at (176,2420)
  {hält zugleich die Projektion der Karte.};

% Scale with geographic miles and the witnessed fractional notation.
\node[mapTown,anchor=east] at (1243,2539) {Maßstab von};
\draw[mapFine] (1250,2538) -- (1425,2538) -- (1515,2537) -- (1575,2534.5) -- (1657,2538) -- (1832,2537.5);
\draw[mapFine] (1250,2543) -- (1425,2543) -- (1515,2542) -- (1575,2539.5) -- (1657,2543) -- (1832,2542.5);
\draw[mapFine] (1250,2538) -- (1250,2543);
\draw[mapFine] (1250,2540.5) -- (1308.2,2540.5);
\draw[mapFine] (1366.4,2540.5) -- (1424.6,2540.5);
\draw[mapFine] (1482.8,2539.9) -- (1515,2539.5) -- (1541,2538.4);
\draw[mapFine] (1599.2,2538) -- (1657,2540.5) -- (1657.4,2540.5);
\draw[mapFine] (1715.6,2540.3) -- (1773.8,2540.2);
\draw[mapFine] (1308.2,2538) -- (1308.2,2543);
\node[mapTick] at (1308.2,2531) {1};
\draw[mapFine] (1366.4,2538) -- (1366.4,2543);
\node[mapTick] at (1366.4,2531) {2};
\draw[mapFine] (1424.6,2538) -- (1424.6,2543);
\node[mapTick] at (1424.6,2531) {3};
\draw[mapFine] (1482.8,2537.4) -- (1482.8,2542.4);
\node[mapTick] at (1482.8,2530.36) {4};
\draw[mapFine] (1541,2535.9) -- (1541,2540.9);
\node[mapTick] at (1541,2528.92) {5};
\draw[mapFine] (1599.2,2535.5) -- (1599.2,2540.5);
\node[mapTick] at (1599.2,2528.53) {6};
\draw[mapFine] (1657.4,2538) -- (1657.4,2543);
\node[mapTick] at (1657.4,2531) {7};
\draw[mapFine] (1715.6,2537.8) -- (1715.6,2542.8);
\node[mapTick] at (1715.6,2530.83) {8};
\draw[mapFine] (1773.8,2537.7) -- (1773.8,2542.7);
\node[mapTick] at (1773.8,2530.67) {9};
\draw[mapFine] (1832,2537.5) -- (1832,2542.5);
\node[mapTick] at (1832,2530.5) {10};
\node[mapTown,anchor=west] at (1837,2540) {geogr. Meilen.};
\node[mapTown] at (1536,2564) {\(\dfrac{1}{1000000}\) der wahren Länge.};
\end{tikzpicture}
\end{center}
\endgroup
\end{landscape}
\restoregeometry
\pdfpagewidth=\HannoverPageWidth\relax
\pdfpageheight=\HannoverPageHeight\relax
\endgroup


% Unnumbered title leaf [435]; principal PDF 442; Google PDF 446.
\clearpage\null\vspace*{5.7cm}
\begin{center}{\fontsize{24}{30}\selectfont HÖHENMESSUNGEN.}\end{center}
\par\vspace*{11.2cm}\noindent\makebox[.86\linewidth][r]{55\textsuperscript{*}}\par
% Blank reverse [436]; principal PDF 443; Google PDF 447.
\clearpage\null
% Unnumbered opening [437]; principal PDF 444; Google PDF 448.
\clearpage
\vspace*{1.5cm}
\begin{center}
{\large\textbf{[Der Refraktionskoeffizient aus den Höhenmessungen bei der hannoverschen Gradmessung.]}}\par\bigskip
\begin{tabular}{c}\hline
\textbf{Astronomisches Jahrbuch für das Jahr 1826.\quad Berlin 1823.\quad S. 89--92.}\\\hline
\end{tabular}\par\bigskip
{\large [1.]}\par\medskip
{\large Beobachtete und berechnete Triangulation im Hannoverschen, Braunschweigschen und Lüneburgschen, von Herrn Hofrat und Ritter \textsc{Gauss} in Göttingen unter dem Datum des 22. Januar 1823 eingesandt.}
\end{center}
\par\noindent\hspace*{\parindent}Obwohl meine Triangulation noch nicht abgeschlossen ist, kann ich sie dank des Anschlusses an die von \textsc{Zach}sche Basis bereits vorläufig berechnen, und die folgenden geografischen Bestimmungen werden schon alle Genauigkeit besitzen, die man sich wünschen kann.
\begin{center}
\begingroup\setlength{\tabcolsep}{4pt}\renewcommand{\arraystretch}{1.3}
\begin{tabular}{l|rrr|rrrr}
&\multicolumn{3}{c|}{Breite}&\multicolumn{4}{c}{\shortstack{[Östl.] Länge\\von Göttingen}}\\
Sternwarte Seeberg&$50^\circ$&$56'$&$6''{,}7$&$+$&$0^\circ$&$47'$&$19''{,}2$\\
Göttingen, neue Sternwarte&51&31&48,7&\multicolumn{4}{c}{0}\\
\quad---\quad Platz der alten Sternwarte&51&31&55,0&$-$&0&0&28,1\\
Brockenhausturm&51&48&2,7&$+$&0&40&22,9\\
Hildesheim, Turm 1&52&9&11,9&$+$&0&0&27,1\\
\quad---\qquad\quad---\quad 2&52&9&16,7&$+$&0&0&3,5\\
\quad---\qquad\quad---\quad 3&52&9&19,4&$+$&0&0&29,4\\
Braunschweig, Martinsturm&52&15&51,5&$+$&0&34&24,6\\
\qquad---\qquad Petrusturm&52&16&4,4&$+$&0&34&22,6\\
\qquad---\qquad Catharinenturm&52&16&9,3&$+$&0&34&57,9\\
\qquad---\qquad Andreasturm&52&16&10,8&$+$&0&34&37,8\\
\end{tabular}\endgroup
\end{center}
% The table continues on printed 438.
